% !TEX program = xelatex
% =====================================================================
%  Introduction to Salam
%  The Official Salam Programming Language Tutorial Book (English)
% ---------------------------------------------------------------------
%  Build:
%     xelatex intro-to-salam.tex     % first pass
%     xelatex intro-to-salam.tex     % second pass (ToC + references)
% =====================================================================
\documentclass[11pt,a4paper,openright]{book}

\usepackage{fontspec}
\usepackage{polyglossia}
\setdefaultlanguage{english}
\setotherlanguage{farsi}
\usepackage{xcolor}
\usepackage{graphicx}
\usepackage{listings}
\usepackage{tcolorbox}
\usepackage{amsmath}
\usepackage{amssymb}
\usepackage{longtable}
\usepackage{booktabs}
\usepackage{array}
\usepackage{enumitem}
\usepackage{microtype}
\usepackage{titlesec}
\usepackage{fancyhdr}
\usepackage[hidelinks]{hyperref}

\tcbuselibrary{skins,breakable}

\setmainfont{Latin Modern Roman}
\setsansfont{Latin Modern Sans}

\newfontfamily\salamcodefont[Path=./../../../../fonts/, Extension=.ttf,
  UprightFont=DejaVuSansMono, BoldFont=DejaVuSansMono-Bold,
  ItalicFont=DejaVuSansMono-Oblique, BoldItalicFont=DejaVuSansMono-BoldOblique,
  Scale=0.86]{DejaVuSansMono}

\newfontfamily\farsifont[Path=./../../fonts/, Extension=.ttf, Script=Arabic,
  Scale=1.0]{XB Roya}
\newfontfamily\arabicfont[Path=./../../fonts/, Extension=.ttf, Script=Arabic,
  Scale=1.0]{XB Roya}
\newfontfamily\salamcodefontfa[Path=./../../fonts/, Extension=.ttf, Script=Arabic,
  Scale=1.0]{XB Roya}

\definecolor{salamKeyword}{HTML}{0000FF}
\definecolor{salamType}{HTML}{267F99}
\definecolor{salamString}{HTML}{A31515}
\definecolor{salamComment}{HTML}{008000}
\definecolor{salamNumber}{HTML}{098658}
\definecolor{salamOp}{HTML}{AF00DB}
\definecolor{salamBg}{HTML}{F9F9F7}
\definecolor{salamFrame}{HTML}{DCDCD7}
\definecolor{noteBg}{RGB}{235,244,252}
\definecolor{noteFrame}{RGB}{90,150,210}
\definecolor{tipBg}{RGB}{235,250,238}
\definecolor{tipFrame}{RGB}{80,170,110}
\definecolor{warnBg}{RGB}{253,243,232}
\definecolor{warnFrame}{RGB}{220,150,60}
\definecolor{deepBg}{RGB}{244,240,250}
\definecolor{deepFrame}{RGB}{150,110,190}
\definecolor{outBg}{RGB}{244,244,244}
\definecolor{outFrame}{RGB}{200,200,200}

\newcommand{\salamcodeENfont}{\salamcodefont}%   English code: DejaVu Mono
\newcommand{\salamcodeFAfont}{\salamcodefontfa}% Persian code: XB Roya
\usepackage{listings}
\usepackage{fancyvrb}
\usepackage{xstring}

\providecommand{\salamcodeENfont}{\ttfamily}
\providecommand{\salamcodeFAfont}{\normalfont}

\definecolor{bg}{HTML}{F4F4F2}
\definecolor{kw}{HTML}{1E40FF}   % keywords
\definecolor{ty}{HTML}{267F99}   % types
\definecolor{cm}{HTML}{008000}   % comments
\definecolor{st}{HTML}{A31515}   % strings
\definecolor{nm}{HTML}{098658}   % numbers

\newcommand{\kw}[1]{\textcolor{kw}{\textbf{#1}}}
\newcommand{\tyword}[1]{\textcolor{ty}{#1}}
\newcommand{\cmtword}[1]{\textcolor{cm}{\textit{#1}}}
\newcommand{\stword}[1]{\textcolor{st}{#1}}
\newcommand{\nmword}[1]{\textcolor{nm}{#1}}

\lstdefinelanguage{Salam}{
    morekeywords={func,ret,if,else,while,for,mut,const,type,struct,enum,end,%
        import,as,true,false,null,this,break,continue,layout,package,print,%
        println,printerr,printerrln,input,defer,operator,extern,interface,%
        pub,component,repeat,impl,to,every,main,sizeof,spawn,join,link},
    morekeywords=[2]{void,bool,char,str,int,uint,float,auto,byte,%
        HashMap,Vector,MapIter,Option,Stack,Queue,Thread,Mutex,WaitGroup},
    sensitive=true,
    morecomment=[l]{//},
    morecomment=[n]{/*}{*/},
    morestring=[b]",
    morestring=[b]',
    morestring=[b]`,
    literate=%
        {i8}{{\tyword{i8}}}2   {i16}{{\tyword{i16}}}3 {i32}{{\tyword{i32}}}3%
        {i64}{{\tyword{i64}}}3 {u8}{{\tyword{u8}}}2   {u16}{{\tyword{u16}}}3%
        {u32}{{\tyword{u32}}}3 {u64}{{\tyword{u64}}}3 {f32}{{\tyword{f32}}}3%
        {f64}{{\tyword{f64}}}3%
        {0}{{\nmword{0}}}1 {1}{{\nmword{1}}}1 {2}{{\nmword{2}}}1%
        {3}{{\nmword{3}}}1 {4}{{\nmword{4}}}1 {5}{{\nmword{5}}}1%
        {6}{{\nmword{6}}}1 {7}{{\nmword{7}}}1 {8}{{\nmword{8}}}1%
        {9}{{\nmword{9}}}1,
}

\lstdefinestyle{salamcardstyle}{
    language=Salam,
    basicstyle=\salamcodeENfont\small,
    backgroundcolor=\color{bg},
    frame=single,
    numbers=left,
    keywordstyle=\color{kw}\bfseries,
    keywordstyle=[2]\color{ty},
    commentstyle=\color{cm}\itshape,
    stringstyle=\color{st},
    showstringspaces=false,
    breaklines=true,
    breakatwhitespace=false,
    tabsize=4,
    columns=fullflexible,
    keepspaces=true,
    upquote=true,
    extendedchars=true,
}

\makeatletter

\@for\salamfa@x:=%
  func,ret,if,else,while,for,mut,const,type,struct,enum,end,import,as,true,%
  false,null,this,break,continue,layout,package,print,println,printerr,%
  printerrln,input,defer,operator,extern,interface,pub,component,repeat,impl,%
  to,every,main,sizeof,spawn,join,link,%
  تابع,بازگشت,اگر,وگرنه,تاوقتی,برای,متغیر,ثابت,نوع,ساختار,شمارش,پایان,%
  واردکردن,بعنوان,درست,نادرست,پوچ,این,بشکن,ادامه,صفحه,بسته,بنویس,چاپ,%
  خطابنویس,خطاچاپ,بخوان,تعلیق,عملگر,خارجی,واسط,عمومی,مولفه,تکرار,پیاده‌سازی,%
  اصلی,تا,هر\do{%
    \expandafter\let\csname salamkw@\salamfa@x\endcsname\@empty}

\@for\salamfa@x:=%
  void,bool,char,str,i8,i16,i32,int,i64,u8,u16,u32,uint,u64,f32,float,f64,%
  auto,byte,HashMap,Vector,MapIter,Option,Stack,Queue,Thread,Mutex,WaitGroup,%
  تهی,منطقی,نویسه,رشته,صحیح۸,صحیح۱۶,صحیح۳۲,صحیح,صحیح۶۴,صحیح8,صحیح16,صحیح32,%
  صحیح64,طبیعی۸,طبیعی۱۶,طبیعی۳۲,طبیعی,طبیعی۶۴,طبیعی8,طبیعی16,طبیعی32,طبیعی64,%
  اعشاری۳۲,اعشاری,اعشاری۶۴,اعشاری32,اعشاری64,خودکار,وکتور,نگاشت\do{%
    \expandafter\let\csname salamty@\salamfa@x\endcsname\@empty}

\@for\salamfa@x:=0,1,2,3,4,5,6,7,8,9,%
  ۰,۱,۲,۳,۴,۵,۶,۷,۸,۹,٠,١,٢,٣,٤,٥,٦,٧,٨,٩\do{%
    \expandafter\let\csname salamfa@dig@\salamfa@x\endcsname\@empty}

\newdimen\salamfa@unit \salamfa@unit=0.5em % indent width per leading space
\newcount\salamfa@lead
\newcount\salamfa@nlines
\newcount\salamfa@cmtdepth                 % /* */ nesting depth (global)
\newif\ifsalamfa@incmt                     % inside a block comment (across lines)
\newif\ifsalamfa@linecmt                   % inside a // comment (this line only)
\newif\ifsalamfa@instr                     % inside a "..." string span
\newcommand{\salamfa@code}{}
\newcommand{\salamfa@nums}{}
\long\def\salamfa@gadd#1{\g@addto@macro\salamfa@code{#1}}
\protected\def\salamfa@end{}% non-expandable end-of-line sentinel
\def\salamfa@endmac{\salamfa@end}
\newif\ifsalamfa@inlead

\newcommand{\salamfa@procline}[1]{%
  \global\advance\salamfa@nlines\@ne
  \salamfa@lead=0
  \salamfa@linecmtfalse \salamfa@instrfalse  % these never span lines
  \ifnum\salamfa@nlines>\@ne \salamfa@gadd{\\}\fi
  \salamfa@gadd{\strut}% keep blank lines non-empty (avoids consecutive \\)
  \edef\salamfa@line{\detokenize{#1}\space\salamfa@end\space}%
  \salamfa@inleadtrue
  \expandafter\salamfa@loop\salamfa@line
}
\def\salamfa@loop#1 {%
  \def\salamfa@tmp{#1}%
  \ifx\salamfa@tmp\salamfa@endmac
  \else
    \ifx\salamfa@tmp\@empty
      \ifsalamfa@inlead \advance\salamfa@lead\@ne \else \salamfa@gadd{\space}\fi
    \else
      \ifsalamfa@inlead
        \salamfa@inleadfalse
        \ifnum\salamfa@lead>\z@
          \edef\salamfa@tmp{\noexpand\hspace*{\the\dimexpr\salamfa@lead\salamfa@unit\relax}}%
          \expandafter\salamfa@gadd\expandafter{\salamfa@tmp}%
        \fi
      \else
        \salamfa@gadd{\space}%
      \fi
      \salamfa@classify{#1}%
    \fi
    \expandafter\salamfa@loop
  \fi}
\newcommand{\salamfa@classify}[1]{%
  \ifsalamfa@linecmt
    \salamfa@gadd{\cmtword{#1}}%
  \else\ifsalamfa@incmt
    \salamfa@gadd{\cmtword{#1}}%
    \IfBeginWith{#1}{/*}{\global\advance\salamfa@cmtdepth\@ne}{}%
    \IfEndWith{#1}{*/}{\salamfa@closecmt}{}%
  \else\ifsalamfa@instr
    \salamfa@gadd{\stword{#1}}%
    \IfEndWith{#1}{"}{\salamfa@instrfalse}{}%
  \else
    \IfBeginWith{#1}{//}%
      {\salamfa@linecmttrue \salamfa@gadd{\cmtword{#1}}}%
      {\IfBeginWith{#1}{/*}{\salamfa@opencmt{#1}}{\salamfa@word{#1}}}%
  \fi\fi\fi}
\def\salamfa@closecmt{%
  \global\advance\salamfa@cmtdepth\m@ne
  \ifnum\salamfa@cmtdepth<\@ne \global\salamfa@incmtfalse\fi}
\newcommand{\salamfa@opencmt}[1]{%
  \salamfa@gadd{\cmtword{#1}}%
  \global\salamfa@cmtdepth\@ne \global\salamfa@incmttrue
  \IfEndWith{#1}{*/}{\salamfa@closecmt}{}}
\newcommand{\salamfa@word}[1]{%
  \IfBeginWith{#1}{"}%
    {\IfEndWith{#1}{"}%
       {\salamfa@gadd{\stword{#1}}}%            complete "..."
       {\salamfa@gadd{\stword{#1}}\salamfa@instrtrue}}% opens a span
    {\ifcsname salamkw@#1\endcsname
       \salamfa@gadd{\kw{#1}}%
     \else\ifcsname salamty@#1\endcsname
       \salamfa@gadd{\tyword{#1}}%
     \else
       \StrLeft{#1}{1}[\salamfa@fc]%
       \ifcsname salamfa@dig@\salamfa@fc\endcsname
         \salamfa@gadd{\nmword{#1}}%
       \else
         \salamfa@gadd{#1}%
       \fi
     \fi\fi}}
\renewcommand{\FancyVerbFormatLine}[1]{\salamfa@procline{#1}}
\def\salamfa@buildnums{%
  \gdef\salamfa@nums{}\@tempcnta=0
  \@whilenum\@tempcnta<\salamfa@nlines\do{%
    \advance\@tempcnta\@ne
    \ifnum\@tempcnta=\@ne
      \xdef\salamfa@nums{\the\@tempcnta}%
    \else
      \xdef\salamfa@nums{\salamfa@nums\noexpand\\\the\@tempcnta}%
    \fi}}
\newbox\salamfa@scratch

\newcommand{\salamfacardbegin}{%
  \gdef\salamfa@code{}\salamfa@nlines=0
  \global\salamfa@cmtdepth=0 \global\salamfa@incmtfalse
  \VerbatimEnvironment
  \setbox\salamfa@scratch=\vbox\bgroup\begin{Verbatim}}
\newcommand{\salamfacardend}{%
  \end{Verbatim}\egroup
  \salamfa@buildnums
  \par\noindent
  \hfuzz=30pt% numbers protrude into the margin by design; don't warn
  \begin{LTR}%
    \fcolorbox{black}{bg}{%
      \begin{minipage}[t]{\dimexpr\linewidth-2\fboxsep-2\fboxrule\relax}%
        \begin{RTL}\raggedleft\salamcodeFAfont\salamfa@code\end{RTL}%
      \end{minipage}%
    }%
    \makebox[0pt][l]{%
      \hspace{\fboxsep}%
      \begin{minipage}[t]{3em}%
        \begin{RTL}\raggedright\salamfa@nums\end{RTL}%
      \end{minipage}%
    }%
  \end{LTR}%
  \par}
\makeatother


\lstnewenvironment{salam}[1][]{\setLTR\lstset{style=salamcardstyle,#1}}{}
\newenvironment{salamfa}{\salamfacardbegin}{\salamfacardend}

\lstdefinestyle{shellstyle}{
  basicstyle=\salamcodefont\small,
  backgroundcolor=\color{outBg},
  frame=single,
  rulecolor=\color{outFrame},
  framesep=6pt,
  numbers=none,
  showstringspaces=false,
  breaklines=true,
  columns=fullflexible,
  keepspaces=true,
  extendedchars=true,
}
\lstnewenvironment{shell}[1][]{%
  \lstset{style=shellstyle,#1}%
}{}
\lstnewenvironment{progout}[1][]{%
  \lstset{style=shellstyle,#1}%
}{}

\newcommand{\code}[1]{{\salamcodefont\small #1}}
\newcommand{\fa}[1]{\textfarsi{#1}}

\newtcolorbox{note}[1][Note]{%
  enhanced,breakable,colback=noteBg,colframe=noteFrame,
  fonttitle=\bfseries,title={#1},left=6pt,right=6pt,top=4pt,bottom=4pt,
  boxrule=0.6pt,arc=2pt}

\newtcolorbox{tip}[1][Tip]{%
  enhanced,breakable,colback=tipBg,colframe=tipFrame,
  fonttitle=\bfseries,title={#1},left=6pt,right=6pt,top=4pt,bottom=4pt,
  boxrule=0.6pt,arc=2pt}

\newtcolorbox{warn}[1][Watch out]{%
  enhanced,breakable,colback=warnBg,colframe=warnFrame,
  fonttitle=\bfseries,title={#1},left=6pt,right=6pt,top=4pt,bottom=4pt,
  boxrule=0.6pt,arc=2pt}

\newtcolorbox{deep}[1][In Depth]{%
  enhanced,breakable,colback=deepBg,colframe=deepFrame,
  fonttitle=\bfseries,title={#1 \textnormal{\small(optional, for advanced readers)}},
  left=6pt,right=6pt,top=4pt,bottom=4pt,boxrule=0.6pt,arc=2pt}

\newcounter{exercise}[chapter]
\renewcommand{\theexercise}{\thechapter.\arabic{exercise}}
\newcommand{\exercise}[1]{\refstepcounter{exercise}\par\medskip%
  \noindent\textbf{Exercise \theexercise.}\ #1\par\medskip}

\pagestyle{fancy}
\fancyhf{}
\fancyhead[LO]{\small\rightmark}
\fancyhead[RE]{\small\leftmark}
\fancyhead[RO,LE]{\small\thepage}
\renewcommand{\headrulewidth}{0.4pt}
\fancypagestyle{plain}{\fancyhf{}\fancyfoot[C]{\thepage}\renewcommand{\headrulewidth}{0pt}}

\hypersetup{
  colorlinks=true,
  linkcolor=salamKeyword,
  urlcolor=salamString,
  citecolor=salamType,
  pdftitle={Introduction to Salam},
  pdfauthor={The Salam Language Team},
  bookmarksnumbered=true,
}

\linespread{1.08}
\setlength{\parskip}{0.35em}
\setlength{\parindent}{0pt}

\begin{document}

\begin{titlepage}
\centering
\vspace*{2.2cm}
{\Huge\bfseries Introduction to Salam\par}
\vspace{0.7cm}
{\LARGE The Official Tutorial\par}
\vspace{1.4cm}
{\large Programming in Your Own Language ---\\[4pt]
from your first program to professional development\par}
\vspace{2.6cm}
\rule{0.65\textwidth}{0.6pt}\par
\vspace{0.8cm}
{\large A bilingual language for everyone:\\[4pt]
students, hobbyists, and professional developers\par}
\vspace{0.8cm}
\rule{0.65\textwidth}{0.6pt}\par
\vfill
{\large The Salam Language Team\par}
\vspace{0.35cm}
{\normalsize First Edition\par}
\end{titlepage}

\thispagestyle{empty}
\noindent\small
This is the official tutorial for the Salam programming language. Every code
listing in this book was tested with the official \code{salam} compiler and is
meant to compile and run as shown. Salam is free software; the language, its
compiler, and its standard library are developed in the open.\par
\vspace{0.3cm}
\noindent The Salam logo and name belong to the Salam Language Team. This book
may be shared and adapted for educational purposes.\par
\vspace{0.3cm}
\noindent First Edition. Typeset with \XeLaTeX{}.\par
\normalsize
\clearpage

\frontmatter
% ===================== Preface =====================
\chapter*{Preface}
\addcontentsline{toc}{chapter}{Preface}
\markboth{Preface}{Preface}

\section*{Why Salam exists}

Almost every programming language in wide use today was designed around
English. Its keywords are English words \code{if}, \code{while},
\code{return} and its tutorials, error messages, and documentation assume
that the reader is comfortable reading English. For hundreds of millions of
people whose first language is Persian or Arabic, that assumption is a wall to
climb before the real learning can even begin.

\textbf{Salam} (from the word \fa{سلام}, meaning \emph{peace}) was created to
lower that wall. It is, at its heart, an ordinary modern programming language:
statically typed, compiled, fast, with structs, generics, interfaces, closures,
a real standard library, and the ability to call C. What makes it unusual is
that it is \emph{bilingual}. The very same language can be written with English
keywords or with Persian keywords, and both produce identical programs. A loop
is spelled \code{while} in one and \fa{تاوقتی} in the other; a function is
\code{func} or \fa{تابع}. The compiler, \code{salam}, understands both.

This means a student in Tehran can write their first program in their own
language, reading their own script, with error messages in their own
words and then, when they are ready, switch the very same concepts into
English to work with the wider software world. The ideas transfer perfectly,
because they are \emph{the same ideas}.

\section*{Who this book is for}

This book assumes \emph{no prior programming experience}. If you have never
written a line of code in your life, start at Chapter~1 and walk forward one
step at a time every term is explained the first time it appears, and every
concept is shown with a small, complete program you can run yourself.

If you are already a programmer coming from C, Python, Go, or Rust, you will
find Salam comfortably familiar. You can move quickly through the early
chapters, treat the \emph{In Depth} boxes as your main course, and use the
appendices as a quick reference. Salam will feel like a language you already
half-know.

The book is therefore written on two levels at once:

\begin{itemize}[leftmargin=1.4em]
  \item The \textbf{main text} is friendly and gradual. It uses plain language,
        small examples, and analogies, and it never assumes you have seen a
        concept before.
  \item The \textbf{In Depth} boxes dig into how things actually work the
        memory model, monomorphization, the C that Salam generates,
        performance trade-offs. They are clearly marked as optional. Skip them
        on a first reading; come back when you are curious.
\end{itemize}

\section*{How to use this book}

Read with your hands, not just your eyes. Programming is a craft, and crafts are
learned by doing. Every chapter contains complete programs type them in, run
them, and then \emph{change} them. Break them on purpose and see what error the
compiler gives. The single most effective way to learn a language is to argue
with its compiler until you both agree.

Each chapter ends with exercises. Attempt them before peeking at
Appendix~D, where solutions to selected exercises are collected. The boxes
sprinkled through the text mean:

\begin{itemize}[leftmargin=1.4em]
  \item \textbf{Note} a clarification or a fact worth remembering.
  \item \textbf{Tip} a practical shortcut or good habit.
  \item \textbf{Watch out} a common mistake; read these twice.
  \item \textbf{In Depth} optional, deeper material for the curious.
\end{itemize}

\section*{A note on scope}

Salam also ships with a built-in \emph{layout} dialect for describing web pages
(it compiles small \code{layout:} blocks straight to HTML and CSS). That side of
the language is a topic of its own and is \emph{not} covered here. This book is
about Salam as a general-purpose programming language: variables, types, control
flow, functions, data structures, the standard library, concurrency, and the
foreign function interface. Everything you learn here is pure Salam
programming.

\begin{note}
All the example programs in this book are drawn from, or directly modelled on,
the real, tested examples that ship with the Salam compiler (under
\code{examples/en/} and \code{tests/en/}). They are not pseudo-code; they
compile and run.
\end{note}

\section*{Conventions}

Salam source code is shown in framed, syntax-highlighted boxes with line
numbers:

\begin{salam}
func main:
    println "Hello, Salam!"
end
\end{salam}

\noindent Terminal commands and program output are shown in plain grey boxes:

\begin{progout}
$ salam exec hello.salam
Hello, Salam!
\end{progout}

\noindent Inline code such as \code{func} or a filename like
\code{hello.salam} appears in a monospaced font. When a Persian keyword is shown
beside its English form, it appears like this: \code{func}~/~\code{\fa{تابع}}.

\bigskip
\noindent Welcome to Salam. Let us begin.

\vspace{1em}
\hfill The Salam Language Team
\clearpage


\tableofcontents
\clearpage

\mainmatter
\chapter{Getting Started}%
\label{ch:getting-started} % chktex 8

In this chapter we take the very first steps: we get the Salam compiler ready,
write our first program, and run it. Do not worry if you have never written a
line of code before this chapter takes you by the hand and walks you forward
one step at a time.

\section{What is a program?}

A computer is a machine that is extraordinarily fast and extraordinarily
literal. It does not know what you want; you must \emph{tell} it, in exact
detail. A list of instructions, written in order so that the computer can carry
out a task, is called a \textbf{program}. The language we write those
instructions in is called a \textbf{programming language}.

Salam is one such language, with one special quality: you can write your
instructions in \emph{Persian} as comfortably as in English. The instruction
``print something on the screen'' is written \code{println "Salam"} in English,
or \code{\fa{چاپ} "\fa{سلام}"} in Persian. Both mean exactly the same thing to
the computer.

\begin{note}
A computer does \emph{precisely} what you tell it no more, no less. So when
a program does not behave the way you expected, the mistake is almost always in
the instructions you gave, not in the machine. This is completely normal; even
professional programmers run into it every single day. Learning to read what the
computer actually did, rather than what you \emph{meant}, is half the craft.
\end{note}

\section{What is a compiler, and why do we need one?}

The program you write is meant to be readable by a \emph{human}. The computer,
however, does not understand it directly deep down it only understands tiny
numeric instructions. We need a program that translates what we wrote into
something the machine can run. That translator is called a \textbf{compiler}.
The compiler for the Salam language is called \code{salam}.

\code{salam} does a great deal of work for us: it reads the text of your
program, checks that it makes sense, reports any mistakes, and finally produces
something you can run.

\begin{deep}[In Depth: the journey of a program]
When \code{salam} processes your program, it passes through several stages:
\emph{lexing} (breaking the text into tokens), \emph{parsing} (building a tree
that represents the structure), \emph{semantic analysis} (checking types and
names), and finally \emph{code generation}. Salam's main backend translates your
program into C source code and then invokes a C compiler (\code{tcc} by
default, but \code{gcc} or \code{clang} work too) to produce a native
executable. There is also a built-in \emph{interpreter} that can run your % chktex 8
program directly without any C toolchain, and an experimental LLVM backend. We
will refer to these stages occasionally, but you never need to know their
details to write programs.
\end{deep}

\section{Getting the compiler}

To work with Salam you need the \code{salam} compiler. It is built from source.
In the root folder of the Salam project, run:

\begin{shell}
sh tools/build-compiler.sh
\end{shell}

When it finishes, an executable called \code{salam} (or \code{salam.exe} on
Windows) appears in the project folder. Alternatively you can build with CMake,
which also builds the runtime library and enables the test suite:

\begin{shell}
cmake -B build && cmake --build build
\end{shell}

\begin{tip}
To confirm the compiler is working, ask it for its version:
\begin{shell}
$ salam version
salam 0.2.6
\end{shell}
\end{tip}

\begin{note}
Throughout this book we assume \code{salam} is reachable from your terminal:
either it sits in your current folder (call it as \code{./salam}) or it is
installed on your system \code{PATH} (call it as \code{salam}). For brevity we
will simply write \code{salam} everywhere.
\end{note}

\section{Your first program: Hello, Salam!}

There is a long-standing tradition among programmers: the very first program you % chktex 8
write in a new language is one that displays the words ``Hello, World!'' Let us
honour that tradition.

Create a file named \code{hello.salam} and type these few lines into it:

\begin{salam}
func main:
    name : str = "Salam"
    println "Hello,", name
    println "2 + 3 =", 2 + 3
end
\end{salam}

\begin{note}
Every Salam program file ends in \code{.salam}. The extension tells the tools
(and you) that the file contains Salam source code.
\end{note}

\subsection{Running the program}

The simplest way to see the result is to use Salam's built-in interpreter, which % chktex 8
runs your program directly no C compiler required. In your terminal, type:

\begin{shell}
salam exec hello.salam
\end{shell}

and the output will be:

\begin{progout}
Hello, Salam
2 + 3 = 5
\end{progout}

If you would rather build a standalone executable (which runs faster and is what
you ship to other people), use the \code{build} command:

\begin{shell}
$ salam build hello.salam --output=hello.exe
$ ./hello.exe
Hello, Salam
2 + 3 = 5
\end{shell}

There is also a one-shot \code{run} command that builds and immediately runs % chktex 8
without leaving any files behind:

\begin{shell}
salam run hello.salam
\end{shell}

\begin{tip}
What is the difference between \code{exec}, \code{run}, and \code{build}?
\code{exec} runs the program \emph{immediately} with the interpreter and needs
no C toolchain it is perfect for learning and quick experiments.
\code{build} produces a native executable you can run again and again, and is
what you use to deliver a finished program. \code{run} is the convenient middle % chktex 12
ground: compile and run once, keep nothing. Throughout most of this book we will
use \code{exec} for simplicity.
\end{tip}

\section{Anatomy of the first program}

Let us now go through the program line by line. Do not skim any part of this ---
these few lines are the doorway to every concept that follows.

\paragraph{Line 1 the start of the program.}
\begin{salam}
func main:
\end{salam}
This line says: ``define a \textbf{function} named \textbf{main}.'' A function is
a named bundle of instructions. The function named \code{main} is the
\emph{starting point} of every Salam program: when you run the program, the
computer begins here. The colon (\code{:}) at the end of the line says that the
body of the function starts on the next line.

\paragraph{Line 2 a variable.}
\begin{salam}
    name : str = "Salam"
\end{salam}
Here we create a \textbf{variable} named \code{name}. Picture a variable as a
labelled box that holds a value. In this box we have placed the text
\code{"Salam"}. The word \code{str} states that the box holds \emph{text} (more
on types in the next chapter). Any piece of text is written between double
quotes (\code{"\,}). Notice that this line begins with a few spaces of
\textbf{indentation}; that whitespace shows the line lives \emph{inside} the
\code{main} function.

\paragraph{Line 3 printing output.}
\begin{salam}
    println "Hello,", name
\end{salam}
The word \code{println} (``print line'') tells the computer to display something
on the screen and then move to the next line. Here we display two things: the
fixed text \code{"Hello,"} and then the contents of the variable \code{name}
(which is \code{"Salam"}). The two are separated by a comma (\code{,}), and Salam
puts a single space between them. That is why the output reads
\code{Hello, Salam}.

\paragraph{Line 4 a little arithmetic.}
\begin{salam}
    println "2 + 3 =", 2 + 3
\end{salam}
Salam does not only print text; it can calculate too. The expression
\code{2 + 3} is worked out \emph{before} printing, and the result (\code{5}) is
displayed. Notice the difference: \code{"2 + 3 ="} is inside quotes, so it is a
piece of \emph{text} printed exactly as written; but \code{2 + 3} is outside
quotes, so it is a \emph{calculation}.

\begin{warn}
Text in quotes is printed verbatim; code outside quotes is evaluated. Writing
\code{println "2 + 3"} would print the three characters \texttt{2 + 3}, not the
number \texttt{5}. The quotes make all the difference.
\end{warn}

\paragraph{Last line the end of the function.}
\begin{salam}
end
\end{salam}
The word \code{end} tells Salam that the body of \code{main} is finished. Any
block that opens with a colon (\code{:}) must be closed with \code{end}. This is
a rule you will see throughout the language: \emph{blocks open with a colon and
close with} \code{end}.

\section{The same program in Persian}

One of the joys of Salam is that any program can be written in Persian too. This
is the exact Persian equivalent of the program above:

\begin{salamfa}
تابع اصلی:
    نام : رشته = "سلام"
    چاپ "سلام،", نام
    چاپ "2 + 3 =", 2 + 3
تمام
\end{salamfa}

The structure is identical; only the words changed:
\code{func}~$\leftrightarrow$~\code{\fa{تابع}},
\code{main}~$\leftrightarrow$~\code{\fa{اصلی}},
\code{str}~$\leftrightarrow$~\code{\fa{رشته}},
\code{println}~$\leftrightarrow$~\code{\fa{چاپ}}, and
\code{end}~$\leftrightarrow$~\code{\fa{تمام}}. A full table of these
equivalents is in Appendix~A. To compile a Persian file, tell the compiler the
language:

\begin{shell}
salam exec hello_fa.salam --lang=fa
\end{shell}

\begin{note}
The entry function's name is itself part of the language: in English it is
\code{main}, in Persian \code{\fa{اصلی}}. We will return to bilingual
programming properly in Chapter~22. For now, just notice that the \emph{ideas}
are the same in both languages only the spelling of the keywords differs.
\end{note}

\section{Comments: notes to yourself}

A \textbf{comment} is text that the compiler ignores completely. Comments are
written for \emph{humans} to explain what a piece of code does, or to leave a
reminder. A line comment starts with \code{//} and runs to the end of the line.
A block comment is wrapped in \code{/* ... */} and may span several lines.

\begin{salam}
func main:
    // This is a line comment: the compiler skips it.
    println "working"          // comments can also follow code

    /* This is a block comment.
       It can stretch across
       as many lines as you like. */
    println "done"
end
\end{salam}

\begin{progout}
working
done
\end{progout}

\begin{tip}
Good comments explain \emph{why}, not \emph{what}. The code already says what it
does; a comment is most valuable when it records the reasoning that the code
cannot show by itself.
\end{tip}

\section{Printing with and without a new line}

Salam has two main ways to display output:

\begin{itemize}[leftmargin=1.4em]
  \item \code{println} displays its arguments and \emph{then moves to the
        next line}.
  \item \code{print} displays its arguments but \emph{stays on the same
        line}; whatever is printed next continues right after it.
\end{itemize}

Look at this example:

\begin{salam}
func main:
    print "Salam "
    print "world"
    println ""
    println "a fresh line"
end
\end{salam}

\begin{progout}
Salam world
a fresh line
\end{progout}

The two \code{print} calls both wrote on the same line, and then
\code{println ""} (which prints an empty piece of text) moved us down to the
next line. This simple trick writing with \code{print} and ending the line
with \code{println ""} will be used again and again in the chapter on loops
to draw shapes.

\begin{deep}[In Depth: writing the arguments of print]
\code{print} and \code{println} are statements, not ordinary functions, so
their arguments are written \textbf{without} wrapping parentheses and
separated by commas: \code{println "hi"} and \code{println name, age}.
Wrapping the whole argument list in parentheses is an error, and Salam will
refuse to compile it. Parentheses inside an argument are of course still
fine, since there they group an expression rather than wrap the list:
\code{println 1, (3 * 4)} and \code{println (2 + 3) * 4} are both correct.
Blocks, likewise, always open with a colon and close with \code{end}; there % chktex 8
is no curly-brace alternative. Remember the \code{salam format} command,
which reformats your code into the canonical style automatically.
\end{deep}

\section{Starting a new project}

When you begin something larger than a single file, the compiler can scaffold a
project folder for you:

\begin{shell}
salam new myproject --lang=en
\end{shell}

This creates a \code{myproject} directory with a starter program already in
place, so you can build and run it immediately. For the small examples in the
coming chapters, a single \code{.salam} file is all you need.

\section{Chapter summary}

In this chapter we learned that:
\begin{itemize}[leftmargin=1.4em]
  \item A program is a list of instructions, and the compiler \code{salam}
        translates it for the computer.
  \item Every Salam program starts at the \code{main} function
        (\code{\fa{اصلی}} in Persian).
  \item Blocks open with a colon (\code{:}) and close with \code{end}.
  \item \code{println} prints and moves to a new line; \code{print} prints and
        stays on the line.
  \item Text goes in double quotes; code outside quotes is evaluated.
  \item Comments (\code{//} and \code{/* ... */}) are ignored by the compiler.
  \item \code{salam exec} runs a program instantly; \code{salam build} makes a
        standalone executable.
  \item Any program can be written in Persian or English with identical meaning.
\end{itemize}

\section{Exercises}

\exercise Write a program that prints your name and your age on two separate
lines.

\exercise Write a program that prints the result of \code{12 * 8} together with a
short label, for example: \texttt{12 * 8 = 96}.

\exercise Write ``Hello, World!'' once in English and once in Persian, and run
both. Is the output the same?

\exercise Using \code{print} (not \code{println}), display your first and last
name on a \emph{single} line with one space between them, then move to a new
line.

\exercise Add a block comment at the top of one of your programs describing what
it does, then run it and confirm the comment changes nothing.

\chapter{Variables and Data Types}
\label{ch:variables-types}

In Chapter~1 we used a variable without dwelling on it. Now we slow down and look
properly at the two ideas that underpin every program: \emph{variables} (named
places to keep values) and \emph{types} (what kind of value a place may hold).
Get these right and everything else in the language falls into place.

\section{Variables: named boxes}

A \textbf{variable} is a name attached to a value. Picture a labelled box: the
label is the name, and the contents are the value. You create a variable by
writing its name, its type, and (usually) an initial value:

\begin{salam}
func main:
    age : int = 30
    name : str = "Sara"
    pi : f64 = 3.14159
    println name, "is", age, "years old"
end
\end{salam}

\begin{progout}
Sara is 30 years old
\end{progout}

Read each declaration left to right: a \emph{name}, then a \emph{type}, then
\code{=} and an \emph{initial value}. So \code{age int = 30} means ``make a
variable called \code{age}, of type \code{int} (a 32-bit integer), holding
\code{30}.''

\begin{note}
Unlike some languages, Salam does \emph{not} use a keyword such as \code{let} or
\code{var} to introduce a variable. The shape ``\textit{name} \textit{type}
\code{=} \textit{value}'' is enough. This keeps declarations short and lines up
the names neatly down the left edge of your code.
\end{note}

\section{Mutable and immutable variables}

By default, a variable in Salam is \textbf{immutable}: once you give it a value,
that value cannot change. This is a deliberate and very useful safety net it
means a value you thought was fixed really is fixed. Try to reassign an immutable
variable and the compiler stops you:

\begin{salam}
func main:
    score : int = 10
    score = 20        // error: 'score' is immutable
end
\end{salam}

When you genuinely need a value to change over time a counter, a running
total, a loop variable mark it \textbf{mutable} with the \code{mut} keyword:

\begin{salam}
func main:
    mut score : int = 10
    score = 20
    score = score + 5
    println score
end
\end{salam}

\begin{progout}
25
\end{progout}

There is also \code{const}, for a value that is fixed \emph{and} known when the
program is compiled a true constant, such as a maximum size or a
mathematical constant:

\begin{salam}
const MAX_USERS = 100

func main:
    println "capacity:", MAX_USERS
end
\end{salam}

\begin{progout}
capacity: 100
\end{progout}

So there are three flavours of binding, from most to least restrictive:

\begin{itemize}[leftmargin=1.4em]
  \item \code{const} fully immutable, and computed at compile time. Neither the
        binding nor its contents can change, not even a single array element or
        struct field. Use for true constants. Conventionally written in
        \code{UPPER\_CASE}.
  \item plain (\emph{bare}) the \emph{binding} is fixed once set: you cannot
        reassign the whole variable. But if it holds a compound value you may
        still change what is \emph{inside} it: array elements and struct fields
        stay writable. This is the default; prefer it.
  \item \code{mut} may be reassigned outright. Use only when a whole value
        really must be replaced.
\end{itemize}

The distinction between a bare binding and \code{const} is about depth. A bare
binding freezes the name; \code{const} freezes everything reachable through it:

\begin{salam}
func main:
    grid : int[3] = [1, 2, 3]
    grid[0] = 10          // fine: the binding is fixed, its elements are not
    // grid = [4, 5, 6]   // error: cannot reassign the whole binding

    const limits : int[3] = [1, 2, 3]
    // limits[0] = 10     // error: a const value is frozen through and through

    println grid[0]
end
\end{salam}

\begin{progout}
10
\end{progout}

\begin{tip}
Reach for \code{mut} only when you must. ``Immutable by default'' is one of
Salam's quiet strengths: the fewer things that can change, the fewer things can
go wrong, and the easier your program is to reason about. If the compiler
complains that something is immutable, pause and ask whether it truly needs to
change before adding \code{mut}.
\end{tip}

\section{Type inference with \texttt{auto}}

Often the type is obvious from the value, and writing it feels redundant. The
\code{auto} keyword asks the compiler to \emph{infer} the type from the
initializer:

\begin{salam}
func main:
    count : auto = 42          // inferred as int
    ratio : auto = 3.14        // inferred as f64
    label : auto = "items"     // inferred as str
    ready : auto = true        // inferred as bool
    println count, ratio, label, ready
end
\end{salam}

\begin{progout}
42 3.14 items true
\end{progout}

\code{auto} is about convenience, not weakness: the variable still has one fixed
type, decided once, at compile time. \code{count} above is exactly an \code{int}
--- you simply did not have to spell it out. You can combine \code{auto} with
\code{mut}: \code{mut total auto = 0} makes a mutable \code{int}.

\begin{warn}
\code{auto} needs an initializer to infer from. \code{x auto} on its own is
meaningless there is nothing to deduce the type from and the compiler
will reject it. When in doubt, or when you want to document intent, write the
type out explicitly.
\end{warn}

\section{The primitive types}

Salam's built-in scalar types are the atoms from which everything larger is
built. Here are the ones you will use constantly.

\subsection{Integers}

Integers are whole numbers. Salam offers them in several sizes, both
\emph{signed} (able to hold negatives) and \emph{unsigned} (zero and up only).
The number in the name is how many bits the value occupies, which fixes its
range:

\begin{center}
\begin{tabular}{@{}llll@{}}
\toprule
\textbf{Type} & \textbf{Signed?} & \textbf{Bits} & \textbf{Range} \\
\midrule
\code{i8}  & signed   & 8  & $-128$ \dots\ $127$ \\
\code{i16} & signed   & 16 & $-32768$ \dots\ $32767$ \\
\code{i32} & signed   & 32 & about $\pm 2.1$ billion \\
\code{i64} & signed   & 64 & about $\pm 9.2$ quintillion \\
\code{u8}  & unsigned & 8  & $0$ \dots\ $255$ \\
\code{u16} & unsigned & 16 & $0$ \dots\ $65535$ \\
\code{u32} & unsigned & 32 & $0$ \dots\ about $4.3$ billion \\
\code{u64} & unsigned & 64 & $0$ \dots\ about $18$ quintillion \\
\bottomrule
\end{tabular}
\end{center}

You do not need to memorise these. For everyday whole numbers, reach for
\code{int} it is the default and is almost always the right choice. Step up to
\code{i64} when you might exceed two billion (timestamps, file sizes, big
counters), and use the small or unsigned types only when you have a specific
reason, such as raw bytes (\code{u8}) or interfacing with C.

\begin{note}
\code{int} \emph{is} \code{i32}: a 32-bit signed integer. It is simply the
friendlier, everyday spelling of the same type, and the compiler treats the two
names as identical. The other two width-carrying types have short spellings as
well: \code{uint} is \code{u32}, and \code{float} is \code{f32}. This book writes
\code{int} and \code{float} in ordinary code, because the width is rarely the
point; it falls back on \code{i32}, \code{u32}, and \code{f32} in the places
where the number of bits genuinely matters, such as the table above or a C
interface. The compiler's own error messages always print the explicit names, so
do not be surprised to see \code{i32} looking back at you from a diagnostic.
\end{note}

\subsection{Floating-point numbers}

For numbers with a fractional part measurements, averages, scientific
values use a floating-point type:

\begin{itemize}[leftmargin=1.4em]
  \item \code{f64} a 64-bit ``double precision'' number. This is the default
        for decimal literals and the one to use unless you have a reason not to.
  \item \code{float} (the explicit spelling is \code{f32}) a 32-bit number;
        less precise, but smaller. Useful for graphics or large arrays where
        space matters.
\end{itemize}

\begin{salam}
func main:
    temperature : f64 = 36.6
    half : auto = 7.0 / 2.0        // 3.5, a floating-point division
    println temperature, half
end
\end{salam}

\begin{progout}
36.6 3.5
\end{progout}

\subsection{Booleans}

A \code{bool} holds one of just two values: \code{true} or \code{false}. Booleans
are the language of decisions every \code{if} and every \code{until} in the
coming chapters tests a \code{bool}.

\begin{salam}
func main:
    raining : bool = false
    weekend : auto = true
    println "raining?", raining
    println "weekend?", weekend
end
\end{salam}

\begin{progout}
raining? false
weekend? true
\end{progout}

\subsection{Characters}

A \code{char} is a single byte one character in the basic Latin range, or any
value from 0 to 255. Character literals are written between \emph{single} quotes:

\begin{salam}
func main:
    grade : char = 'A'
    println grade
    println grade as int        // the numeric code of 'A'
    println 66 as char          // the character with code 66
end
\end{salam}

\begin{progout}
A
65
B
\end{progout}

A \code{char} is really a small number in disguise: \code{'A'} \emph{is} 65. You
can convert between the two with \code{as}, which is how the example turns
\code{'A'} into \code{65} and \code{66} back into \code{'B'}.

\begin{note}
A single quote makes a \code{char}; a double quote makes a \code{str}. So
\code{'A'} is one character, while \code{"A"} is a string that happens to be one
character long. They are different types. We meet strings properly in
Chapter~8.
\end{note}

\subsection{Strings}

A \code{str} is a piece of text a sequence of characters written between
double quotes. Strings are \emph{immutable}: operations on them produce new
strings rather than changing the original. You can join two strings with
\code{+}:

\begin{salam}
func main:
    first : str = "Salam"
    greeting : auto = first + ", world!"
    println greeting
end
\end{salam}

\begin{progout}
Salam, world!
\end{progout}

\subsection{Void}

Finally, \code{void} is the ``no value'' type. You never store a \code{void};
it appears only as the result type of a function that returns nothing (a
function that \emph{does} something rather than \emph{computes} something). We
return to it in Chapter~6.

\section{Literals: writing values directly}

A \textbf{literal} is a value written straight into your program. Each kind of
literal has a natural type:

\begin{center}
\begin{tabular}{@{}lll@{}}
\toprule
\textbf{Literal} & \textbf{Type} & \textbf{Examples} \\
\midrule
whole number      & \code{int} & \code{0}, \code{42}, \code{-7} \\
decimal number    & \code{f64} & \code{3.14}, \code{2.0}, \code{1.5e3} \\
character          & \code{char} & \code{'a'}, \code{'9'} \\
text               & \code{str} & \code{"hi"}, \code{""} \\
truth value        & \code{bool} & \code{true}, \code{false} \\
\bottomrule
\end{tabular}
\end{center}

This leads to one rule that surprises newcomers and is worth learning now.

\begin{warn}
\textbf{A plain whole-number literal is an \code{int}.} Salam does not silently
shrink it to fit a smaller or unsigned type. So this is an \emph{error}:
\begin{salam}
mut b : u8 = 250           // error: cannot assign 'i32' to 'u8'
\end{salam}
The literal \code{250} is an \code{int}, and Salam will not narrow it for you
without being asked. (The compiler spells the type \code{i32} in its message;
that is the same type under its explicit name.) The fix is an explicit cast with
\code{as}:
\begin{salam}
mut b : u8 = 250 as u8     // fine
\end{salam}
This strictness is a feature: it makes every change of type visible in the
source, so a narrowing can never happen by accident.
\end{warn}

\section{Converting between types with \texttt{as}}

Salam never converts a value from one type to another behind your back. When you
\emph{do} want a conversion, you ask for it explicitly with the \code{as}
operator:

\begin{salam}
func main:
    total : int = 7
    count : int = 2
    average : f64 = (total as f64) / (count as f64)
    println "average =", average
end
\end{salam}

\begin{progout}
average = 3.5
\end{progout}

Why the casts? Because \code{total} and \code{count} are integers,
\code{total / count} would be \emph{integer} division it would compute
\code{3}, throwing away the remainder. By converting each to \code{f64} first, we
ask for real division and get \code{3.5}.

\begin{warn}
Integer division truncates toward zero: \code{7 / 2} is \code{3}, not
\code{3.5}, and \code{7 \% 2} (the remainder) is \code{1}. If you want a
fractional result, at least one operand must be a floating-point value. This is
one of the most common beginner surprises keep it in mind.
\end{warn}

\section{Type aliases}

When a type's meaning matters more than its mechanics, you can give it a clearer
name with \code{type}:

\begin{salam}
type Age = u32

func main:
    a : Age = 30 as u32
    println a
end
\end{salam}

\begin{progout}
30
\end{progout}

An alias is \emph{transparent}: \code{Age} and \code{u32} are completely
interchangeable the alias is just a more readable spelling. Aliases come into
their own with the bigger types we build later; for now, simply know the keyword
exists.

\section{Naming variables}

A good name is documentation that never goes out of date. Salam identifiers may
contain letters, digits, and underscores, and may not start with a digit.
Crucially, they may also contain \emph{Unicode} letters so Persian names are
first-class:

\begin{salam}
func main:
    user_count : int = 5
    maxScore : auto = 100
    println user_count, maxScore
end
\end{salam}

\begin{tip}
Choose names that say what the value \emph{means}, not what type it is.
\code{user\_count} tells the reader far more than \code{n} or \code{intval}. The
extra keystrokes pay for themselves the first time you re-read your own code a
week later.
\end{tip}

\section{Persian type names}

Just as the keywords have Persian forms, so do the types. A Persian program
writes the very same declarations with Persian type names:

\begin{center}
\begin{tabular}{@{}ll@{}}
\toprule
\textbf{English} & \textbf{Persian} \\
\midrule
\code{int} & \code{\fa{صحیح}} \\
\code{i64} & \code{\fa{صحیح۶۴}} \\
\code{f64} & \code{\fa{اعشار۶۴}} \\
\code{bool} & \code{\fa{منطقی}} \\
\code{char} & \code{\fa{کاراکتر}} \\
\code{str} & \code{\fa{رشته}} \\
\code{auto} & \code{\fa{خودکار}} \\
\bottomrule
\end{tabular}
\end{center}

\noindent Here is the chapter's first program written in Persian:

\begin{salamfa}
تابع اصلی:
    سن : صحیح = 30
    نام : رشته = "سارا"
    چاپ نام, "is", سن, "years old"
تمام
\end{salamfa}

Just as \code{int} is the everyday spelling in English, \code{\fa{صحیح}} is the
everyday spelling in Persian, and the two name exactly the same 32-bit type. When
you do want to state a width, the Persian name carries it in Persian digits:
\code{\fa{صحیح۶۴}} is \code{i64}, and the \fa{۶۴} inside it is part of the
keyword's spelling. Values, by contrast, are written with ordinary digits (the
\code{30} above), because arithmetic uses the standard digit characters. We will
untangle this fully in Chapter~22; for now, just admire that the structure is
identical to the English version.

\section{In Depth: how values are stored}

\begin{deep}[In Depth: representation and the generated C]
Each Salam scalar maps directly onto a C type when the program is compiled:
\code{i32} (that is, \code{int}) becomes \code{int32\_t}, \code{f64} becomes
\code{double},
\code{bool} becomes a one-byte value, \code{char} a single byte, and \code{str}
a pointer to immutable, NUL-terminated bytes. This 1-to-1 mapping is why Salam
programs run at C speed and why the foreign function interface (Chapter~19) is
so direct. Because the sizes are fixed and known, the compiler can check at
\emph{compile time} that \code{250} does not fit in an \code{i8} and that a
\code{u8} and an \code{i32} are not accidentally mixed the strictness you met
above is the visible edge of this static type system.
\end{deep}

\section{Chapter summary}

\begin{itemize}[leftmargin=1.4em]
  \item A variable is declared as \textit{name} \textit{type} \code{=}
        \textit{value}; no \code{let}/\code{var} keyword is needed.
  \item Variables are \emph{immutable by default}; add \code{mut} to allow
        reassignment, or \code{const} for a compile-time constant.
  \item \code{auto} infers the type from the initializer.
  \item Integers come in signed (\code{i8}--\code{i64}) and unsigned
        (\code{u8}--\code{u64}) sizes; the 32-bit one is the everyday default.
        \code{int}/\code{uint}/\code{float} are the short spellings of
        \code{i32}/\code{u32}/\code{f32}, and this book writes \code{int} and
        \code{float} in ordinary code.
  \item \code{f64} is the default for decimals; \code{bool} is \code{true} or
        \code{false}; \code{char} is one byte; \code{str} is immutable text.
  \item A whole-number literal is an \code{int}; conversions are never implicit ---
        use \code{as}. Integer division truncates.
  \item Types have Persian names too (\code{\fa{صحیح}}, \code{\fa{رشته}},
        \dots).
\end{itemize}

\section{Exercises}

\exercise Declare three variables describing a book: a title (\code{str}), a page
count (\code{int}), and a price (\code{f64}). Print them in one sentence.

\exercise Start with \code{mut total int = 0}, add \code{17} and then \code{25}
to it in two steps, and print the result after each step.

\exercise Compute the average of \code{10}, \code{15}, and \code{20} as an
\code{f64} and print it. (Remember what plain integer division would do.)

\exercise Store the character \code{'Z'} in a \code{char}, then print both the
character and its numeric code.

\exercise Try to assign \code{300} to a \code{u8} variable and read the compiler's
error. Then fix it with a cast and explain in a comment what the cast does.

\exercise Rewrite the book-description program from the first exercise using
Persian type names.

\chapter{Operators and Expressions}%
\label{ch:operators}

So far our programs have mostly stored and printed values. Now they start to
\emph{compute}. An \textbf{operator} is a symbol that combines values to produce
a new one \code{+} adds, \code{<} compares, \code{\&\&} reasons about truth.
A combination of values and operators that produces a result is an
\textbf{expression}. This chapter is a tour of Salam's operators and the rules
that decide what an expression means.

\section{Arithmetic operators}

The five everyday arithmetic operators work as you would expect:

\begin{center}
\begin{tabular}{@{}cll@{}}
\toprule
\textbf{Operator} & \textbf{Meaning} & \textbf{Example} \\
\midrule
\code{+} & addition       & \code{3 + 4}  $\to$ \code{7} \\
\code{-} & subtraction    & \code{10 - 6} $\to$ \code{4} \\
\code{*} & multiplication & \code{6 * 7}  $\to$ \code{42} \\
\code{/} & division       & \code{17 / 5} $\to$ \code{3} \\
\code{\%} & remainder (modulo) & \code{17 \% 5} $\to$ \code{2} \\
\bottomrule
\end{tabular}
\end{center}

\begin{salam}
func main:
    println 3 + 4 * 2          // 11
    println 17 / 5, 17 % 5     // 3 2
    println 10 - 6, 6 * 7      // 4 42
end
\end{salam}

\begin{progout}
11
3 2
4 42
\end{progout}

\subsection{Integer versus floating-point division}

We met this in Chapter~2, but it is worth repeating because it trips up almost
everyone once. When \emph{both} operands of \code{/} are integers, you get
\emph{integer} division: the result is a whole number, with any fractional part
thrown away.

\begin{salam}
func main:
    println 7 / 2          // 3  (integer division)
    println 7.0 / 2.0      // 3.5 (floating-point division)
    println 7 as f64 / 2   // 3.5 (one float operand is enough)
end
\end{salam}

\begin{progout}
3
3.5
3.5
\end{progout}

The \code{\%} operator gives the \emph{remainder} of an integer division:
\code{7 \% 2} is \code{1}, because 7 is 2 threes with 1 left over. Modulo is
surprisingly useful: \code{n \% 2 == 0} tests whether \code{n} is even, and
\code{n \% 10} extracts the last decimal digit of \code{n}.

\subsection{The power operator}

Salam has a dedicated exponentiation operator, \code{**}. It raises the left
operand to the power of the right and produces a \emph{floating-point} result:

\begin{salam}
func main:
    println 2 ** 10        // 1024
    x : f64 = 9 ** 0.5       // 3 -- a square root, since x^(1/2) = sqrt(x)
    println x
end
\end{salam}

\begin{progout}
1024
3
\end{progout}

\begin{note}
Because \code{**} always yields a floating-point number, \code{2 ** 10} is the
floating value \code{1024.0}; it just prints as \code{1024} because there is no
fractional part to show. If you need an integer power and the values are small,
plain multiplication (or a loop, as in Chapter~5) keeps the result an integer.
\end{note}

\section{Compound assignment}

When you update a variable using its own current value a very common pattern ---
a \textbf{compound assignment} operator says it more briefly. Each one combines an
arithmetic operation with assignment:

\begin{center}
\begin{tabular}{@{}cl@{}}
\toprule
\textbf{Shorthand} & \textbf{Means} \\
\midrule
\code{x += y}  & \code{x = x + y} \\
\code{x -= y}  & \code{x = x - y} \\
\code{x *= y}  & \code{x = x * y} \\
\code{x /= y}  & \code{x = x / y} \\
\code{x \%= y} & \code{x = x \% y} \\
\code{x **= y} & \code{x = x ** y} \\
\bottomrule
\end{tabular}
\end{center}

\begin{salam}
func main:
    mut n : int = 10
    n += 5       // n is now 15
    n -= 2       // 13
    n *= 2       // 26
    println n

    mut p : f64 = 2.0
    p **= 10     // same as p = p ** 10
    println p
end
\end{salam}

\begin{progout}
26
1024
\end{progout}

\begin{warn}
Compound assignment only works on a \code{mut} variable it \emph{changes} the
variable, after all. Used on an immutable binding, it is the same immutability
error you saw in Chapter~2.
\end{warn}

Because \code{**} always produces a floating-point result (see the previous
section), \code{x **= y} only type-checks when \code{x} itself is a \code{f64}
(or a \code{f32}, or a type with an \code{operator **} overload from
Chapter~13) for the same reason \code{x = x ** y} would need \code{x} to accept
a float on the right-hand side.

\begin{note}
Every compound assignment is checked with exactly the same rules as writing the
operation out longhand. \code{x \%= y} still requires both sides to be integers,
and \code{x += y} still refuses to mix a \code{bool} with anything numeric, even
though \code{x} and \code{y} are the same type. In other words, \code{x op= y}
is never more permissive than \code{x = x op y}; it is purely a shorter way to
write it.
\end{note}

\section{Comparison operators}

A comparison asks a yes/no question about two values and answers with a
\code{bool}. There are six:

\begin{center}
\begin{tabular}{@{}cll@{}}
\toprule
\textbf{Operator} & \textbf{Question} & \textbf{Example} \\
\midrule
\code{==} & equal?                 & \code{5 == 5} $\to$ \code{true} \\
\code{!=} & not equal?             & \code{5 != 3} $\to$ \code{true} \\
\code{<}  & less than?             & \code{3 < 5}  $\to$ \code{true} \\
\code{>}  & greater than?          & \code{3 > 5}  $\to$ \code{false} \\
\code{<=} & less than or equal?    & \code{5 <= 5} $\to$ \code{true} \\
\code{>=} & greater than or equal? & \code{4 >= 5} $\to$ \code{false} \\
\bottomrule
\end{tabular}
\end{center}

\begin{salam}
func main:
    println 5 > 3, 5 == 5, 5 != 3, 5 <= 5
end
\end{salam}

\begin{progout}
true true true true
\end{progout}

\begin{warn}
Do not confuse \code{=} with \code{==}. A single \code{=} \emph{assigns} a value
to a variable; a double \code{==} \emph{compares} two values and yields a
\code{bool}. Writing \code{=} where you meant \code{==} is a classic slip ---
fortunately Salam's static checks catch most such mistakes for you.
\end{warn}

Comparisons also work on strings: \code{==} and \code{!=} test whether two
strings have the same contents (not merely whether they are the same object),
and \code{<}, \code{>}, and friends compare them in dictionary order.

\begin{salam}
func main:
    println "abc" == "abc"     // true
    println "abc" == "abd"     // false
end
\end{salam}

\begin{progout}
true
false
\end{progout}

\section{Logical operators}

Logical operators combine \code{bool} values the raw material of decisions.
There are three:

\begin{center}
\begin{tabular}{@{}cll@{}}
\toprule
\textbf{Operator} & \textbf{Meaning} & \textbf{True when\dots} \\
\midrule
\code{\&\&} & and & both sides are true \\
\code{||}   & or  & at least one side is true \\
\code{!}    & not & the operand is false \\
\bottomrule
\end{tabular}
\end{center}

\begin{salam}
func main:
    age : int = 20
    has_ticket : bool = true
    can_enter : auto = age >= 18 && has_ticket
    println can_enter
    println true || false, !true
end
\end{salam}

\begin{progout}
true
true false
\end{progout}

\begin{deep}[In Depth: short-circuit evaluation]
\code{\&\&} and \code{||} are \emph{short-circuiting}: they evaluate the
right-hand side only if they must. In \code{a \&\& b}, if \code{a} is already
\code{false}, the whole expression is \code{false} regardless of \code{b}, so
\code{b} is never evaluated. Likewise \code{a || b} skips \code{b} when \code{a}
is \code{true}. This is not just an optimization it lets you write a guard and
a use of it in one line, e.g. testing that a count is positive \emph{before}
dividing by it, knowing the risky part runs only when it is safe.
\end{deep}

\begin{note}
Salam spells the logical operators with symbols (\code{\&\&}, \code{||},
\code{!}), not words. There is no \code{and}, \code{or}, or \code{not} keyword.
\end{note}

\section{A note on bitwise operators}

Some languages provide \emph{bitwise} operators (\code{\&}, \code{|},
\code{<<}, and so on) that manipulate the individual bits of an integer. The
current version of Salam does \emph{not} include them writing \code{6 | 1},
for instance, is a syntax error (the compiler even reminds you that \code{||} is
the logical ``or''). If you need bit manipulation today, you can fall back to
arithmetic (multiplying or dividing by powers of two shifts bits) or call a C
function through the foreign function interface of Chapter~19.

\section{String concatenation}

You have already seen that \code{+} joins two strings end to end. This is the
same \code{+} symbol as numeric addition, but Salam picks the right meaning from
the types of the operands: two numbers add, two strings concatenate.

\begin{salam}
func main:
    first : str = "Salam"
    full : auto = first + ", " + "world!"
    println full
end
\end{salam}

\begin{progout}
Salam, world!
\end{progout}

We devote all of Chapter~8 to strings; for now, just remember that \code{+}
between two \code{str} values builds a new, larger string.

\section{Precedence and grouping}

When several operators appear in one expression, \textbf{precedence} decides who
acts first. The rules follow ordinary mathematics: \code{**} binds tightest, then
\code{*}, \code{/}, and \code{\%}, then \code{+} and \code{-}, then the
comparisons, and finally the logical operators. So

\begin{salam}
func main:
    println 3 + 4 * 2          // 11, not 14: * happens before +
    println 2 + 3 > 4          // true: 2+3 first, then compare
end
\end{salam}

\begin{progout}
11
true
\end{progout}

When in doubt or simply to make your intent unmistakable add parentheses.
They cost nothing and override any precedence rule:

\begin{salam}
func main:
    println (3 + 4) * 2        // 14: parentheses force + first
end
\end{salam}

\begin{progout}
14
\end{progout}

\begin{tip}
You do not need to memorise the full precedence table. Two habits cover almost
everything: remember that multiplication and division come before addition and
subtraction, and reach for parentheses whenever an expression is even slightly
hard to read at a glance. Clear beats clever.
\end{tip}

\section{Putting it together}

Here is a small program that uses several operator families at once: it decides
whether a year is a leap year, a genuinely useful rule built entirely from
\code{\%}, comparison, and logical operators.

\begin{salam}
func main:
    year : int = 2024
    leap : auto = (year % 4 == 0 && year % 100 != 0) || (year % 400 == 0)
    println year, "leap?", leap
end
\end{salam}

\begin{progout}
2024 leap? true
\end{progout}

Read the condition in English: a year is a leap year if it is divisible by 4 but
\emph{not} by 100 unless it is also divisible by 400. The operators let us say
that precisely. In Chapter~4 we will use such conditions to make the program
actually \emph{do} different things.

\section{Chapter summary}

\begin{itemize}[leftmargin=1.4em]
  \item Arithmetic: \code{+ - * / \%}, plus \code{**} for powers (which yields a
        floating-point result).
  \item Integer \code{/} truncates; use a float operand for fractional results.
        \code{\%} gives the remainder.
  \item Compound assignment (\code{+=}, \code{-=}, \code{*=}, \code{/=},
        \code{\%=}, \code{**=}) updates a \code{mut} variable in place, and is
        checked with exactly the same type rules as its longhand form.
  \item Comparisons (\code{== != < > <= >=}) yield a \code{bool} and also work on
        strings.
  \item Logical operators are \code{\&\&}, \code{||}, \code{!}, and they
        short-circuit. There are no word forms and no bitwise operators.
  \item \code{+} concatenates strings.
  \item Precedence follows mathematics; parentheses override it and aid clarity.
\end{itemize}

\section{Exercises}

\exercise Print the results of \code{20 / 3}, \code{20 \% 3}, and
\code{20.0 / 3.0}. Explain in a comment why the first two differ from the third.

\exercise Given \code{mut x int = 100}, use compound assignment to halve it, then
add 7, then print it.

\exercise Write a boolean expression that is \code{true} exactly when an integer
\code{n} is between 1 and 12 inclusive, and test it with a couple of values.

\exercise Without running it first, predict the value of
\code{2 + 3 * 4 > 10 \&\& 1 < 2}. Then run it and check.

\exercise Build a person's full name from \code{first} and \code{last} string
variables with a space between them, using \code{+}, and print it.

\exercise Given \code{mut side : f64 = 3.0}, use \code{**=} to compute the cube
of \code{side} in place, then print it. What error do you get if you try the
same thing starting from \code{mut side : int = 3} instead? Why?

\chapter{Control Flow}
\label{ch:control-flow}

Until now our programs have run straight through, top to bottom, doing the same
thing every time. Real programs make \emph{decisions}: do this when the user is
old enough, do that otherwise. The tool for choosing between paths is the
\code{if} statement, and it turns the boolean expressions of Chapter~3 into
action.

\section{The \texttt{if} statement}

An \code{if} runs a block of code only when a condition is true. The condition is
any expression of type \code{bool}; the block opens with a colon and closes with
\code{end}:

\begin{salam}
func main:
    age : int = 20
    if age >= 18:
        println "You may enter."
    end
    println "Done."
end
\end{salam}

\begin{progout}
You may enter.
Done.
\end{progout}

If \code{age} were \code{15}, the condition \code{age >= 18} would be
\code{false}, the indented block would be skipped entirely, and only
\code{Done.} would print. The structure is always the same: \code{if}, a
condition, a colon, the body, and \code{end}.

\section{\texttt{else}: the other path}

Often you want one thing to happen when the condition holds and a \emph{different}
thing otherwise. That is what \code{else} is for:

\begin{salam}
func main:
    age : int = 15
    if age >= 18:
        println "You may enter."
    else:
        println "Sorry, too young."
    end
end
\end{salam}

\begin{progout}
Sorry, too young.
\end{progout}

Exactly one of the two blocks runs never both, never neither. Think of it as a
fork in the road: the condition decides which way to go.

\section{\texttt{else if}: choosing among many}

When there are more than two possibilities, chain them with \code{else if}. Salam
checks each condition in order and runs the block belonging to the \emph{first}
one that is true; if none match, the final \code{else} (if present) runs.

\begin{tip}
Notice that Salam does not repeat the word \code{if} after \code{else}. Put a
condition directly after \code{else} and it behaves as \code{else if}; leave
\code{else} with no condition and it is the final, catch-all branch.
\end{tip}

\begin{salam}
func main:
    score : int = 83
    if score >= 90:
        println "Grade: A"
    else score >= 80:
        println "Grade: B"
    else score >= 70:
        println "Grade: C"
    else:
        println "Grade: F"
    end
end
\end{salam}

\begin{progout}
Grade: B
\end{progout}

The order matters. \code{83} satisfies \code{>= 80}, so we print \code{B} and
stop the remaining branches are not even tested. Because the checks run top
to bottom, arrange them from the most specific or highest threshold downward, as
here.

\section{The inline form}

When a branch's body is a single short statement, you may write it on the same
line as the colon. This keeps tidy little decisions compact:

\begin{salam}
func main:
    n : int = 7
    if n % 2 == 0: println "even"
    else:          println "odd"
    end
end
\end{salam}

\begin{progout}
odd
\end{progout}

The \code{end} is still required; only the line breaks change. Many of the
example programs that ship with Salam use this form to line up short branches
neatly, like a small table. Use it when it improves readability, and the
multi-line form when a body grows beyond a single statement.

\section{Conditions are just boolean expressions}

There is nothing special about the condition in an \code{if} it is any
expression that produces a \code{bool}, exactly the kind you built in Chapter~3.
That means you can combine tests with \code{\&\&}, \code{||}, and \code{!}:

\begin{salam}
func main:
    age : int = 20
    has_ticket : bool = true
    if age >= 18 && has_ticket:
        println "Welcome to the show."
    else:
        println "Entry denied."
    end
end
\end{salam}

\begin{progout}
Welcome to the show.
\end{progout}

You can also store a condition in a \code{bool} variable first and test that ---
sometimes a well-named boolean makes the intent far clearer than a long condition
crammed into the \code{if}:

\begin{salam}
func main:
    age : int = 20
    has_ticket : bool = true
    let_in : auto = age >= 18 && has_ticket
    if let_in:
        println "Welcome to the show."
    end
end
\end{salam}

\begin{warn}
Salam's conditions must be genuine \code{bool} values. Unlike some languages, you
cannot use a number as a shortcut for ``truth'' there is no rule that
\code{0} means false and everything else means true. Write the comparison out:
\code{count != 0} rather than just \code{count}. This keeps conditions explicit
and removes a whole category of subtle bugs.
\end{warn}

\section{Nesting}

The body of an \code{if} is an ordinary block, so it may contain another
\code{if}. This is called \emph{nesting}, and it lets you ask a follow-up
question only after a first one has been answered:

\begin{salam}
func main:
    logged_in : bool = true
    is_admin : bool = false
    if logged_in:
        if is_admin:
            println "Welcome, administrator."
        else:
            println "Welcome, user."
        end
    else:
        println "Please log in."
    end
end
\end{salam}

\begin{progout}
Welcome, user.
\end{progout}

\begin{tip}
Deep nesting quickly becomes hard to follow. Two habits keep it shallow: combine
related tests with \code{\&\&} instead of nesting two \code{if}s, and once you
reach Chapter~6 use an early \code{ret} to handle a special case and leave the
main path un-indented. If you find yourself four levels deep, that is usually a
sign to restructure.
\end{tip}

\section{Decisions inside functions}

Conditions become especially powerful inside functions (the subject of
Chapter~6), where a branch can \emph{return} a result. Here is a function that
classifies a score, using \code{else if} with the inline \code{ret} form:

\begin{salam}
func grade(score : int) : str:
    if score >= 90: ret "A"
    else score >= 80: ret "B"
    else score >= 70: ret "C"
    else: ret "F"
    end
end

func main:
    println grade(95), grade(83), grade(71), grade(50)
end
\end{salam}

\begin{progout}
A B C F
\end{progout}

Each branch hands back a different string, and the first matching one wins. We
will return to functions properly soon; for now, notice how naturally
\code{if}/\code{else if} expresses ``pick the right answer.''

\section{In Persian}

Conditionals read just as naturally in Persian. The keywords are \code{\fa{اگر}}
for \code{if}, \code{\fa{وگرنه}} for \code{else}, and \code{\fa{تمام}} for
\code{end}:

\begin{salamfa}
تابع اصلی:
    سن : صحیح = 15
    اگر سن >= 18:
        چاپ "بزرگسال"
    وگرنه:
        چاپ "کم‌سن"
    تمام
تمام
\end{salamfa}

An \code{else if} chain is written the same way as in English: put a condition
right after \code{\fa{وگرنه}}, with no separate word for \code{if}.

\section{Chapter summary}

\begin{itemize}[leftmargin=1.4em]
  \item \code{if cond:} runs a block only when \code{cond} (a \code{bool}) is
        true; close it with \code{end}.
  \item \code{else} provides the alternative path; exactly one branch runs.
  \item \code{else if} chains several conditions; the first true one wins.
  \item A single-statement branch may be written inline after the colon.
  \item Conditions are ordinary boolean expressions and must be genuine
        \code{bool}s numbers are not truthy.
  \item \code{if}s may be nested, but prefer \code{\&\&} and early returns to keep
        nesting shallow.
\end{itemize}

\section{Exercises}

\exercise Write a program that reads a fixed integer variable \code{temp} and
prints \code{"freezing"} if it is below 0, \code{"cold"} if below 15,
\code{"warm"} if below 30, and \code{"hot"} otherwise.

\exercise Using a single \code{if} with \code{\&\&}, print \code{"valid"} only
when a variable \code{hour} is between 0 and 23 inclusive.

\exercise Write a function \code{sign(n int) str} that returns \code{"negative"},
\code{"zero"}, or \code{"positive"}, and test it with three values.

\exercise Rewrite the grade example so the inline branches become multi-line
blocks that each also print a short message. Does the output change?

\exercise Translate the freezing/cold/warm/hot program into Persian.

\chapter{Loops}%
\label{ch:loops}

Computers are tireless. Where a person would balk at adding the numbers from 1 to
a million, a computer does it without complaint \emph{if} we can express the
work as a repetition. A \textbf{loop} runs a block of code over and over, and it
is what turns a few lines of source into millions of actions. Salam gives you
three looping tools: \code{until}, the counting \code{repeat}, and the
collection-walking \code{each}. We will meet all three.

\section{The \texttt{until} loop}

An \code{until} loop repeats its body \emph{as long as} a condition stays true. It
checks the condition, runs the body if true, then checks again round and round
until the condition becomes false.

\begin{salam}
func main:
    mut i : int = 1
    until i <= 3:
        print i
        print " "
        i += 1
    end
    println ""
end
\end{salam}

\begin{progout}
1 2 3
\end{progout}

Trace it: \code{i} starts at 1; \code{1 <= 3} is true, so we print \code{1} and
bump \code{i} to 2; and so on until \code{i} becomes 4, at which point
\code{4 <= 3} is false and the loop stops. The variable that drives the loop ---
here \code{i} must be \code{mut}, because the loop changes it.

\begin{warn}
Every \code{until} loop needs a way to \emph{end}. If the condition never becomes
false for example if you forget the \code{i += 1} above the loop runs
forever and your program hangs. Whenever you write an \code{until}, find the line
that makes progress toward the exit and check that it is really there.
\end{warn}

\code{until} shines when you do not know in advance how many times you will loop.
A classic example is Euclid's algorithm for the greatest common divisor, which
loops until a remainder reaches zero:

\begin{salam}
func gcd(a : int, b : int) : int:
    mut x : int = a
    mut y : int = b
    until y != 0:
        mut t : int = y
        y = x % y
        x = t
    end
    ret x
end

func main:
    println gcd(48, 36)
end
\end{salam}

\begin{progout}
12
\end{progout}

\section{The \texttt{repeat} loop}

When you \emph{do} know the pattern of counting a set number of times, a range
of numbers, or every $S$-th number in that range \code{repeat} is the tool.
It reads almost like plain English, and it is Salam's direct replacement for the
classic counting loop found in other languages. It comes in a few forms.

\subsection{Repeat a fixed number of times}

\code{repeat N:} runs its body \code{N} times. There is no loop variable to
manage when you only need the \emph{repetition}, this is the cleanest tool:

\begin{salam}
func main:
    repeat 3:
        println "Salam!"
    end
end
\end{salam}

\begin{progout}
Salam!
Salam!
Salam!
\end{progout}

If the count is zero or negative, the body simply does not run.

Sometimes you also want to know \emph{which} pass you are on, counting from
zero. Add \code{with} and a name to bind it:

\begin{salam}
func main:
    repeat 5 with i:
        print i
    end
    println ""
end
\end{salam}

\begin{progout}
01234
\end{progout}

\code{repeat 5 with i} binds \code{i} to \code{0}, \code{1}, \code{2}, \code{3},
then \code{4} one pass per value. This is the direct replacement for the
classic ``count from 0 up to $n - 1$'' loop that other languages spell with a
three-part \code{for} header.

\begin{warn}
Do not spell ``count from 0 up to $n - 1$'' as \code{repeat 0 to n - 1 with i}.
It looks reasonable, but when \code{n} is \code{0} the range \code{0 to -1} is
read as counting \emph{downward}, so the body still runs once when you meant it
to run zero times. The safe, idiomatic way to write that pattern is
\code{repeat n with i}, exactly as above; let \code{repeat} pick the count,
do not hand it a hand-built range that can point the wrong way.
\end{warn}

\subsection{Repeat over a range}

\code{repeat A to B:} runs the body once for every value from \code{A} up to
\code{B} \emph{inclusive} so \code{repeat 1 to 5} runs five times. When you
just need the repetition and not the value, that is all you need:

\begin{salam}
func main:
    mut sum : int = 0
    mut k : int = 0
    repeat 1 to 5:
        k += 1
        sum += k
    end
    println "1+2+3+4+5 =", sum
end
\end{salam}

\begin{progout}
1+2+3+4+5 = 15
\end{progout}

More often, though, you want to read the current value directly instead of
tracking your own counter. Add \code{with} and a name, just like with
\code{repeat N}:

\begin{salam}
func main:
    repeat 1 to 10 with n:
        println "7 x", n, "=", 7 * n
    end
end
\end{salam}

\begin{progout}
7 x 1 = 7
7 x 2 = 14
...
7 x 10 = 70
\end{progout}

\begin{note}
The upper bound is \emph{inclusive}: \code{repeat 1 to 5} covers
1,~2,~3,~4,~5 five iterations, not four. This matches how we count in everyday
speech (``count from 1 to 5''). Keep the distinction in mind whenever you set
up a range: the loop always includes both endpoints.
\end{note}

\subsection{Repeat with a step}

Add \code{by S} to stride by \code{S} instead of 1. Combine it with \code{with}
to read each value as you go:

\begin{salam}
func main:
    repeat 0 to 20 by 2 with i:
        print i
        print " "
    end
    println ""
end
\end{salam}

\begin{progout}
0 2 4 6 8 10 12 14 16 18 20
\end{progout}

If you only care how many times the loop ran and not the values themselves,
drop \code{with} just as with any other \code{repeat} range:

\begin{salam}
func main:
    mut count : int = 0
    repeat 0 to 20 by 2:
        count += 1
    end
    println "even numbers in 0..20:", count
end
\end{salam}

\begin{progout}
even numbers in 0..20: 11
\end{progout}

There are eleven of them (0,~2,~\dots,~20), and the loop counts each another
reminder that the upper bound is included.

\begin{tip}
\code{repeat} covers counting, but it is not the tool for walking the elements
of an array or other collection directly. That job belongs to \code{each}, which
we meet properly in Chapter~7: \code{each x in numbers:} hands you each element
in turn, with no index bookkeeping at all.
\end{tip}

\section{Breaking out and skipping ahead}

Two keywords give you finer control \emph{inside} any loop:

\begin{itemize}[leftmargin=1.4em]
  \item \code{break} stops the loop immediately and continues with whatever comes
        after it.
  \item \code{continue} abandons the current pass and jumps straight to the next
        one.
\end{itemize}

\begin{salam}
func main:
    repeat 6 with n:
        if n == 3: continue end      // skip 3
        if n == 5: break end         // stop before 5
        print n
    end
    println ""
end
\end{salam}

\begin{progout}
0124
\end{progout}

The loop prints 0, 1, and 2 normally; at \code{n == 3} the \code{continue} skips
the \code{print}, so 3 never appears; at \code{n == 5} the \code{break} ends the
loop before printing. Hence \code{0124}.

\begin{tip}
\code{break} and \code{continue} are handy, but a loop peppered with them can be
hard to follow. Often a clearer condition in the \code{until} or \code{repeat}
header expresses the same intent without jumping. Reserve \code{break} for cases
like ``search until found,'' where stopping early is the whole point.
\end{tip}

\section{Nested loops}

A loop body may itself contain a loop. Nested loops are how you work with grids,
tables, and shapes. Here a loop over rows contains a loop over columns to draw a
triangle of stars using the ``print then end the line'' trick from Chapter~1:

\begin{salam}
func main:
    mut row : int = 0
    repeat 1 to 5:
        row += 1
        repeat 1 to row:
            print "*"
        end
        println ""
    end
end
\end{salam}

\begin{progout}
*
**
***
****
*****
\end{progout}

The outer loop runs once per row; for each row, the inner loop prints as many
stars as the row number, and then \code{println ""} moves to the next line.
Nested loops multiply: an outer loop of 5 around an inner loop of up to 5 does up
to 25 iterations in total.

\begin{warn}
\code{break} and \code{continue} affect only the \emph{innermost} loop that
encloses them. A \code{break} inside the inner star loop would stop that row's
stars, not the whole triangle.
\end{warn}

\section{In Persian}

The loop keywords have Persian forms too: \code{\fa{تاوقتی}} is the Persian
spelling of \code{until}, \code{\fa{تکرار}} is the Persian spelling of
\code{repeat}, and \code{\fa{هر}} is the Persian spelling of \code{each}. Inside
a \code{repeat}, \code{\fa{تا}} spells the range word \code{to}, and
\code{\fa{گام}} spells the step word \code{by}. Salam has never had a C-style
\code{for}, so there is no Persian word tied to one \code{repeat} and
\code{each} cover that ground instead. Here is the star loop in Persian:

\begin{salamfa}
تابع اصلی:
    تکرار 1 تا 5:
        بنویس "*"
    تمام
    چاپ ""
تمام
\end{salamfa}

\section{Which loop should I use?}

\begin{itemize}[leftmargin=1.4em]
  \item Reach for \code{repeat N} when you just need to do something a set number
        of times, and \code{repeat A to B} (optionally with \code{by}) when you
        are counting over a clear range.
  \item Add \code{with i} to a \code{repeat} when you need to read the current
        number inside the body especially the classic ``count from 0 up to
        $n - 1$'' pattern, which is simply \code{repeat n with i}.
  \item Use \code{each} when you are walking the elements of an array or other
        collection directly, rather than counting (see Chapter~7).
  \item Use \code{until} when the number of iterations is not known up front and
        the loop ends when some condition changes.
\end{itemize}

\section{Chapter summary}

\begin{itemize}[leftmargin=1.4em]
  \item \code{until cond:} repeats while a boolean condition holds; make sure it
        can end.
  \item \code{repeat N}, \code{repeat N with i}, \code{repeat A to B}, and
        \code{repeat A to B by S} cover fixed counts, indices, and ranges; the
        upper bound of a range is \emph{inclusive}. Never spell ``count from 0''
        as \code{repeat 0 to n - 1 with i} use \code{repeat n with i}.
  \item \code{each x in collection:} walks the elements of a collection directly
        (Chapter~7).
  \item \code{break} ends the innermost loop; \code{continue} skips to its next
        pass.
  \item Loops nest, and that is how you handle grids and shapes.
\end{itemize}

\section{Exercises}

\exercise Use an \code{until} loop to print the numbers from 10 down to 1, then
print \code{"Lift-off!"}.

\exercise Use a \code{repeat} loop to compute and print the sum of all numbers
from 1 to 100.

\exercise Using \code{repeat ... by}, print every multiple of 5 from 5 to 50 on
one line, separated by spaces.

\exercise Use a loop with \code{continue} to print the numbers from 1 to 20,
skipping every multiple of 3.

\exercise Draw an upside-down triangle of stars (five stars on the first row, one
on the last) using nested loops.

\exercise Rewrite the times-table example so it prints the table for a number
stored in a variable, in Persian.

% ===================== Chapter 6 =====================
\chapter{Functions}
\label{ch:functions}

We have used \code{main} since the first page, and a few helper functions along
the way. Now we study functions properly, because they are the single most
important tool for organising a program. A \textbf{function} is a named,
reusable piece of code that takes some inputs, does some work, and (often) hands
back a result. Functions let you write a piece of logic once, give it a name,
and call it as many times as you like.

\section{Defining and calling a function}

You have already seen the shape. A function definition is the keyword
\code{func}, a name, a parenthesised list of parameters, an optional return
type, a colon, the body, and \code{end}:

\begin{salam}
func square(n i32) i32:
    ret n * n
end

func main:
    println square(5)
    println square(9)
end
\end{salam}

\begin{progout}
25
81
\end{progout}

\code{square} takes one \textbf{parameter}, \code{n} of type \code{i32}, and its
\textbf{return type} (written after the parameter list) is also \code{i32}. The
body computes \code{n * n} and hands it back with \code{ret}. In \code{main} we
\textbf{call} the function by writing its name and an \textbf{argument} in
parentheses: \code{square(5)}. The argument \code{5} becomes the parameter
\code{n} for that call.

\begin{note}
\emph{Parameter} and \emph{argument} are two ends of the same connection. The
parameter is the name in the definition (\code{n}); the argument is the value you
pass at the call site (\code{5}). The argument is copied into the parameter when
the call begins.
\end{note}

\section{Returning a value}

The \code{ret} keyword does two things at once: it ends the function immediately
and sends a value back to the caller. The type of that value must match the
function's declared return type.

A function can have several \code{ret} statements useful when different cases
produce different answers. The first one reached wins:

\begin{salam}
func abs(n i32) i32:
    if n < 0:
        ret -n
    end
    ret n
end

func main:
    println abs(-7), abs(7)
end
\end{salam}

\begin{progout}
7 7
\end{progout}

If \code{n} is negative, the first \code{ret} fires and the function is done; the
final \code{ret n} is reached only for non-negative inputs. Using an early
\code{ret} to dispatch a special case, then letting the main path continue
un-indented, is a clean and common style.

\section{Functions that return nothing}

Some functions \emph{do} something rather than \emph{compute} something they
print, they update, they save a file. Such a function has no return type: you
simply omit it after the parameter list. Its (implicit) result type is
\code{void}, the ``no value'' type from Chapter~2.

\begin{salam}
func greet(name str):
    println "Hello,", name
    println "Welcome to Salam."
end

func main:
    greet("Sara")
    greet("Ali")
end
\end{salam}

\begin{progout}
Hello, Sara
Welcome to Salam.
Hello, Ali
Welcome to Salam.
\end{progout}

Inside a \code{void} function you may still use a bare \code{ret} (with no value)
to return early for instance to bail out when an input is invalid but you
never return a value from one.

\section{Multiple parameters}

A function may take any number of parameters, separated by commas. Each needs its
own type:

\begin{salam}
func add(a i32, b i32) i32:
    ret a + b
end

func max3(a i32, b i32, c i32) i32:
    mut m i32 = a
    if b > m: m = b end
    if c > m: m = c end
    ret m
end

func main:
    println add(3, 4)
    println max3(7, 2, 9)
end
\end{salam}

\begin{progout}
7
9
\end{progout}

Arguments are matched to parameters \emph{by position}: in \code{add(3, 4)},
\code{3} becomes \code{a} and \code{4} becomes \code{b}.

\begin{warn}
The current version of Salam does not support \emph{default} parameter values, so
every parameter must be supplied at every call. Calling \code{add(3)} is an
error. (To offer a ``default,'' write a second overload see below.)
\end{warn}

\section{Arguments are passed by value}

When you pass an argument, the function receives a \emph{copy} of it. Changing a
parameter inside the function therefore does \emph{not} change the caller's
variable:

\begin{salam}
func tryToChange(x i32):
    mut local i32 = x
    local = 999            // only the copy changes
end

func main:
    n i32 = 5
    tryToChange(n)
    println n              // still 5
end
\end{salam}

\begin{progout}
5
\end{progout}

This ``pass by value'' rule makes functions easy to reason about: a function
cannot secretly reach back and alter the variables you passed in. When you want a
function to produce a new value, the clean approach is to \emph{return} it and let
the caller decide what to do:

\begin{salam}
func incremented(x i32) i32:
    ret x + 1
end

func main:
    mut n i32 = 5
    n = incremented(n)     // the caller chooses to update n
    println n
end
\end{salam}

\begin{progout}
6
\end{progout}

\begin{note}
Salam deliberately has no ``pass by reference'' for ordinary variables you
cannot take the address of a local with \code{\&} and hand it to a function. This
keeps the flow of data visible: values go \emph{in} as arguments and come
\emph{out} as return values. (In Chapter~9 you will see that a struct's own
methods can update the struct they belong to; that is the sanctioned way to
modify data in place.)
\end{note}

\section{Recursion}

A function is allowed to call \emph{itself}. This is called \textbf{recursion},
and it is the natural way to express anything defined in terms of a smaller
version of itself. The classic example is the factorial, where
$n! = n \times (n-1)!$:

\begin{salam}
func fact(n i32) i32:
    if n <= 1:
        ret 1
    end
    ret n * fact(n - 1)
end

func main:
    println fact(5)
end
\end{salam}

\begin{progout}
120
\end{progout}

Every recursive function needs two things to be correct: a \textbf{base case}
that stops the recursion (here, \code{n <= 1} returns \code{1} without calling
again), and a \textbf{recursive step} that moves \emph{toward} the base case
(here, calling \code{fact(n - 1)} with a smaller argument). Trace \code{fact(5)}:
it becomes \code{5 * fact(4)}, which becomes \code{5 * 4 * fact(3)}, and so on
down to \code{fact(1)}, which is \code{1} then the multiplications unwind:
$5\times4\times3\times2\times1 = 120$.

\begin{warn}
Forget the base case or fail to shrink the argument and a recursive
function calls itself forever, until the program runs out of stack space and
crashes. Always ask: ``what stops this?'' and ``does each call get closer to
stopping?''
\end{warn}

\begin{deep}[In Depth: recursion versus iteration]
Anything you can compute with recursion you can also compute with a loop, and
vice versa. Recursion often expresses a problem more clearly (tree walks, divide
and conquer), while a loop avoids the small cost of repeated function calls. For
factorials a loop is perfectly natural:
\begin{salam}
func fact(n i32) i32:
    mut result i32 = 1
    for mut i = 2; i <= n; i = i + 1:
        result = result * i
    end
    ret result
end
\end{salam}
Choose whichever makes the intent clearest; both are idiomatic Salam.
\end{deep}

\section{Definition order does not matter}

You may call a function that is defined \emph{later} in the file. Salam reads all
the top-level definitions before it runs anything, so \code{main} can freely call
helpers written below it:

\begin{salam}
func main:
    println later()
end

func later() str:
    ret "defined after main"
end
\end{salam}

\begin{progout}
defined after main
\end{progout}

This means you never need ``forward declarations,'' and you are free to put
\code{main} at the top of the file where a reader expects to find it.

\section{Function overloading}

Two functions may share a name as long as their parameters differ in number or
in type. This is called \textbf{overloading}, and Salam picks the right one based
on the arguments at each call:

\begin{salam}
func add(a i32, b i32) i32:
    ret a + b
end

func add(a f64, b f64) f64:
    ret a + b
end

func main:
    println add(2, 3)        // calls the i32 version
    println add(1.5, 2.5)    // calls the f64 version
end
\end{salam}

\begin{progout}
5
4
\end{progout}

The two \code{add}s are genuinely different functions that happen to share a name.
Overloading lets a family of related operations present one consistent name
instead of \code{addInt}, \code{addFloat}, and so on. (The second result prints as
\code{4} because it is the floating value \code{4.0}.)

\section{Use every parameter}

Salam is tidy about parameters: if you declare one and never use it, the compiler
reports an error rather than letting the dead weight slip by. If you genuinely
need a parameter you do not use perhaps to match an expected shape prefix
its name with an underscore to say ``unused on purpose'':

\begin{salam}
func ignore_second(a i32, _b i32) i32:
    ret a
end
\end{salam}

\begin{tip}
This rule nudges you toward honest signatures: a function's parameter list should
list exactly what it needs. When the compiler flags an unused parameter, it is
often pointing at a leftover from an earlier version of the code that you can
simply delete.
\end{tip}

\section{Designing good functions}

Functions are how you tame complexity. A few guidelines that pay off:

\begin{itemize}[leftmargin=1.4em]
  \item \textbf{One job each.} A function should do a single, nameable thing. If
        you struggle to name it, it is probably doing too much.
  \item \textbf{Short is good.} If a function grows past a screenful, look for a
        piece you can lift out into its own well-named helper.
  \item \textbf{Name for the result.} \code{is\_leap(year)} and
        \code{average(scores)} tell the reader what comes back. A good name turns
        a call site into a readable sentence.
  \item \textbf{Prefer returning over reaching out.} Compute a value and return
        it, rather than relying on shared mutable state.
\end{itemize}

\section{Chapter summary}

\begin{itemize}[leftmargin=1.4em]
  \item A function is \code{func name(params) RetType:} \dots\ \code{end};
        \code{ret} returns a value and ends the function.
  \item Omit the return type for a function that returns nothing (\code{void});
        a bare \code{ret} returns early from one.
  \item Parameters are matched by position and passed \emph{by value} (the
        function gets a copy). There are no default parameters and no reference
        parameters.
  \item A function may call itself (recursion); give it a base case and make each
        call smaller.
  \item Definition order is irrelevant you can call functions defined later.
  \item Functions may be \emph{overloaded} by differing parameters.
  \item Every declared parameter must be used, or prefixed with \code{\_}.
\end{itemize}

\section{Exercises}

\exercise Write a function \code{cube(n i32) i32} that returns $n^3$ using
multiplication, and print the cubes of 1 through 5 with a loop.

\exercise Write a \code{void} function \code{banner(text str)} that prints the
text surrounded by a line of dashes above and below it.

\exercise Write \code{min2(a i32, b i32) i32} and use it to build
\code{min3(a, b, c)} by calling \code{min2} twice.

\exercise Write a recursive function \code{sum\_to(n i32) i32} that returns
$1 + 2 + \dots + n$. Identify its base case and recursive step in comments.

\exercise Overload a function \code{describe} so that \code{describe(42)} prints
``an integer'' and \code{describe(3.14)} prints ``a float''. Test both.

\exercise Rewrite the recursive \code{sum\_to} from a previous exercise as a loop,
and confirm both give the same answers.

% ===================== Chapter 7 =====================
\chapter{Arrays and Vectors}
\label{ch:arrays-vectors}

A single variable holds one value. But programs constantly deal with
\emph{collections} a row of scores, a list of names, a grid of cells. Salam
gives you two tools for holding many values under one name: the \textbf{array},
whose size is fixed when you write it, and the \textbf{vector}, which grows and
shrinks as your program runs. This chapter covers both.

\section{Fixed-size arrays}

An \textbf{array} is a numbered row of values, all of the same type. You declare
one by writing the element type followed by the length in square brackets, and
you can give the initial contents in a bracketed list:

\begin{salam}
func main:
    scores i32[5] = [90, 85, 70, 95, 60]
    println scores[0]      // the first element
    println scores[4]      // the last element
end
\end{salam}

\begin{progout}
90
60
\end{progout}

The type \code{i32[5]} reads ``an array of five \code{i32}s.'' The values are
stored in \emph{slots} numbered from \code{0}, so a length-5 array has slots
\code{0}, \code{1}, \code{2}, \code{3}, \code{4}. You reach a slot with square
brackets: \code{scores[0]} is the first, \code{scores[4]} the last.

\begin{warn}
Indices start at \emph{zero}, not one. The first element is \code{scores[0]} and
the \emph{last} of a length-$n$ array is \code{scores[n-1]} here
\code{scores[4]}, not \code{scores[5]}. Off-by-one slips around this rule are
among the most common bugs in all of programming; the zero-based habit becomes
second nature with practice.
\end{warn}

\subsection{Changing elements}

To modify an element, assign to its slot. As with any modification, the array
must be \code{mut}:

\begin{salam}
func main:
    mut a i32[3] = [1, 2, 3]
    a[1] = 20
    println a[0], a[1], a[2]
end
\end{salam}

\begin{progout}
1 20 3
\end{progout}

\subsection{Walking an array with a loop}

The \code{for} loop and arrays are made for each other: the loop variable steps
through the valid indices, and the body does something with each element. Here we
sum an array and find its largest value in a single pass:

\begin{salam}
func main:
    a i32[7] = [5, 9, 1, 7, 3, 8, 2]
    mut sum i32 = 0
    mut hi i32 = a[0]
    for mut i = 0; i < 7; i = i + 1:
        sum += a[i]
        if a[i] > hi: hi = a[i] end
    end
    println "sum =", sum, "max =", hi
end
\end{salam}

\begin{progout}
sum = 35 max = 9
\end{progout}

Notice the loop condition \code{i < 7}: it stops \emph{before} index 7, visiting
exactly slots 0 through 6. This half-open pattern start at 0, continue while
\code{i < length} is the canonical way to walk an array, and it lines up
perfectly with zero-based indexing.

\section{Multi-dimensional arrays}

An array's elements can themselves be arrays, giving you a grid. \code{i32[2][3]}
is ``2 rows of 3 columns.'' You index with two pairs of brackets, row first:

\begin{salam}
func main:
    mut grid i32[2][3] = [[1, 2, 3], [4, 5, 6]]
    grid[1][2] = 60
    println grid[0][0], grid[1][2]
end
\end{salam}

\begin{progout}
1 60
\end{progout}

Nested loops walk a grid naturally: an outer loop over rows, an inner loop over
columns. This is the foundation for anything tabular game boards, matrices,
spreadsheets.

\section{Slices: a window onto an array}

Often you want to work with just \emph{part} of an array the middle four
elements, everything after the third, the first half without copying it into a
new array. A \textbf{slice} is a lightweight \emph{view} onto a stretch of an
array. It does not own any data and it does not allocate memory on the heap;
under the hood it is nothing more than a pointer to the first element plus a
length. Languages such as Go and Rust have the same idea.

You create a slice with a \textbf{range} inside the brackets, \code{a[lo:hi]}:

\begin{salam}
func main:
    nums i32[6] = [10, 20, 30, 40, 50, 60]
    middle auto = nums[1:4]
    println middle[0], middle[1], middle[2]
    println len(middle)
end
\end{salam}

\begin{progout}
20 30 40
3
\end{progout}

The range is \textbf{half-open}: \code{nums[1:4]} starts \emph{at} index~1 and
stops \emph{before} index~4, so it covers indices 1, 2 and 3 a length of
$4-1=3$. This is the same half-open convention as the loop condition
\code{i < length}, and it has a tidy consequence: the length of \code{a[lo:hi]}
is simply \code{hi - lo}, and \code{a[lo:hi]} followed by \code{a[hi:end]} covers
the whole range with no gap and no overlap.

\subsection{Leaving out a bound}

If you omit the low bound it defaults to \code{0}; omit the high bound and it runs
to the end of the array. Omit both, \code{a[:]}, and you get a view of the whole
array.

\begin{salam}
func main:
    nums i32[6] = [10, 20, 30, 40, 50, 60]
    println len(nums[:3])     // first three: 10, 20, 30
    println len(nums[3:])     // from index 3 on: 40, 50, 60
    println len(nums[:])      // the whole array
end
\end{salam}

\begin{progout}
3
3
6
\end{progout}

\subsection{The slice type \texttt{T[:]}}

A slice has its own type, written \code{T[:]} a slice of \code{T} as opposed to
\code{T[N]}, a fixed array of \code{N} items. You will most often see it as a
parameter type, because a function that takes \code{i32[:]} can be handed
\emph{any} run of \code{i32}s: a whole array, or just a piece of one.

\begin{salam}
func total(view i32[:]) i32:
    mut sum i32 = 0
    for mut i = 0; i < len(view); i = i + 1:
        sum = sum + view[i]
    end
    ret sum
end

func main:
    nums i32[6] = [10, 20, 30, 40, 50, 60]
    println total(nums[:])      // all six
    println total(nums[2:5])    // just 30 + 40 + 50
end
\end{salam}

\begin{progout}
210
120
\end{progout}

One \code{total} function now works on the whole array and on any slice of it. You
do not need a separate copy, and you do not pass the length as an extra argument
\code{len} reads it straight from the slice.

\subsection{A slice is a window, not a copy}

Because a slice borrows the array's storage rather than copying it, writing
through a slice changes the original array. This is the key difference from
copying out a sub-array, and it is exactly what makes slices cheap.

\begin{salam}
func main:
    mut nums i32[6] = [10, 20, 30, 40, 50, 60]
    mut tail auto = nums[4:]
    tail[0] = 500
    tail[1] = 600
    println nums[4], nums[5]
end
\end{salam}

\begin{progout}
500 600
\end{progout}

Assigning to \code{tail[0]} reaches through the window and overwrites
\code{nums[4]}. The slice stored no data of its own; it only ever pointed at the
array.

\begin{warn}
A slice borrows storage it does not own. Keep the underlying array alive for as
long as you use the slice, and remember that a slice has a \emph{fixed} length:
like a fixed array it does not grow. When you need something that grows, reach for
a vector (next).
\end{warn}

\section{The limit of fixed arrays}

A fixed array's length is part of its type and is decided when you write the
program. That is perfect when you know the count up front (a week has 7 days, a
board has 64 squares). But often you do \emph{not} know how many items there will
be you are reading lines from a file, or collecting whatever the user types.
For that you need a collection that can grow. Enter the vector.

\section{Vectors: arrays that grow}

A \textbf{vector}, written \code{Vector<T>}, is a sequence of \code{T} values
whose length can change at run time. You start with an empty one and add to it as
you go. Create an empty vector with \code{Vector \{\}}:

\begin{salam}
func main:
    mut v Vector<i32> = Vector {}
    v.push(10)
    v.push(20)
    v.push(30)
    println "length:", v.len()
end
\end{salam}

\begin{progout}
length: 3
\end{progout}

The angle brackets in \code{Vector<i32>} say what the vector holds here,
\code{i32}s. (This is your first glimpse of \emph{generics}, the subject of
Chapter~12: \code{Vector} is one type that works for any element type.) The
\code{.push(x)} method appends a value to the end, and \code{.len()} reports how
many values are currently stored.

\subsection{Reading and writing elements}

A vector's \code{.get(i)} method returns a \emph{pointer} to the element at index
\code{i}; to read the value itself, follow the pointer with \code{[0]}. To
overwrite an element, use \code{.set(i, value)}:

\begin{salam}
func main:
    mut v Vector<i32> = Vector {}
    v.push(10)
    v.push(20)
    v.push(30)
    println "first:", v.get(0)[0]
    v.set(0, 100)
    println "first now:", v.get(0)[0]
end
\end{salam}

\begin{progout}
first: 10
first now: 100
\end{progout}

\begin{note}
Why \code{v.get(i)[0]} rather than just \code{v[i]}? \code{.get} returns a
\emph{pointer} into the vector's storage, and \code{[0]} reads through that
pointer to the value. The pointer form is what makes vectors efficient no
copying and you will understand it fully after the chapter on memory
(Chapter~18). For now, treat \code{v.get(i)[0]} as the idiom for ``the value at
index \code{i}.''
\end{note}

\subsection{Walking a vector}

Because \code{.len()} can change, you typically walk a vector with a \code{while}
loop bounded by its current length:

\begin{salam}
func main:
    mut v Vector<i32> = Vector {}
    v.push(5)
    v.push(8)
    v.push(2)
    mut total i32 = 0
    mut i i32 = 0
    while i < v.len():
        total += v.get(i)[0]
        i += 1
    end
    println "total:", total
end
\end{salam}

\begin{progout}
total: 15
\end{progout}

\subsection{Removing from the end}

\code{.pop()} removes the last element and returns it turning a vector into a
\emph{stack}, where the last thing in is the first thing out:

\begin{salam}
func main:
    mut v Vector<i32> = Vector {}
    v.push(1)
    v.push(2)
    v.push(3)
    println "pop:", v.pop()      // 3
    println "pop:", v.pop()      // 2
    println "left:", v.len()     // 1
end
\end{salam}

\begin{progout}
pop: 3
pop: 2
left: 1
\end{progout}

\subsection{The common vector operations}

The methods you will reach for most often:

\begin{center}
\begin{tabular}{@{}ll@{}}
\toprule
\textbf{Method} & \textbf{What it does} \\
\midrule
\code{v.push(x)}     & append \code{x} to the end \\
\code{v.pop()}       & remove and return the last element \\
\code{v.len()}       & the number of elements \\
\code{v.get(i)[0]}   & read the element at index \code{i} \\
\code{v.set(i, x)}   & overwrite the element at index \code{i} \\
\code{v.is\_empty()} & \code{true} if the vector has no elements \\
\code{v.cap()}       & current capacity (see the In Depth box) \\
\code{v.free()}      & release the vector's memory \\
\bottomrule
\end{tabular}
\end{center}

\begin{warn}
When you are finished with a vector, call \code{v.free()} to release the memory it
allocated. Unlike a fixed array (whose storage is reclaimed automatically when it
goes out of scope), a vector manages a growable buffer that you are responsible
for freeing. We discuss this fully in Chapter~18; for now, make \code{v.free()}
the last thing you do with a vector.
\end{warn}

\begin{deep}[In Depth: capacity versus length, and why push is cheap]
A vector keeps two numbers: its \emph{length} (how many elements you have stored)
and its \emph{capacity} (how many it could store before needing a bigger buffer).
When you push past the capacity, the vector quietly allocates a larger buffer ---
typically doubling and copies the elements across. Because the capacity grows
geometrically, the \emph{average} cost of a \code{push} is constant, even though
an occasional one triggers a copy. That is why building a vector element by
element is efficient, and why \code{v.cap()} is usually larger than
\code{v.len()}: pushing two items into a fresh vector may already reserve room
for four.
\end{deep}

\section{Arrays or vectors?}

\begin{itemize}[leftmargin=1.4em]
  \item Use a \textbf{fixed array} when you know the number of elements when you
        write the code, and it will not change coordinates, a fixed lookup
        table, the cells of a board.
  \item Use a \textbf{vector} when the count is decided at run time or changes as
        the program works reading input, accumulating results, anything
        ``however many there turn out to be.''
\end{itemize}

\section{A worked example: collecting even numbers}

This program walks the numbers 1 to 10, collects the even ones into a vector, and
reports what it found a pattern (filter into a growing collection) you will use
constantly:

\begin{salam}
func main:
    mut evens Vector<i32> = Vector {}
    for mut n = 1; n <= 10; n = n + 1:
        if n % 2 == 0:
            evens.push(n)
        end
    end
    print "found "
    print evens.len()
    println " even numbers"
    mut i i32 = 0
    while i < evens.len():
        print evens.get(i)[0]
        print " "
        i += 1
    end
    println ""
    evens.free()
end
\end{salam}

\begin{progout}
found 5 even numbers
2 4 6 8 10
\end{progout}

\section{Chapter summary}

\begin{itemize}[leftmargin=1.4em]
  \item A fixed array \code{T[N]} holds \code{N} values of type \code{T};
        elements are indexed from \code{0} to \code{N-1}.
  \item Walk an array with \code{for mut i = 0; i < N; i = i + 1}.
  \item Arrays can nest (\code{T[R][C]}) to form grids, indexed row-first.
  \item A \code{Vector<T>} grows at run time; create it with \code{Vector \{\}}.
  \item Key vector methods: \code{push}, \code{pop}, \code{len}, \code{get(i)[0]}
        to read, \code{set(i, x)} to write, \code{is\_empty}, \code{free}.
  \item Call \code{v.free()} when done with a vector.
  \item Choose arrays for fixed counts, vectors for counts that vary.
\end{itemize}

\section{Exercises}

\exercise Create an array of the five vowels (\code{char} values) and print each
on its own line with a loop.

\exercise Given \code{a i32[6] = [4, 8, 15, 16, 23, 42]}, compute and print the
average as an \code{f64}.

\exercise Build a vector of the squares $1, 4, 9, \dots, 100$ using a loop, print
its length, then print every element.

\exercise Write a program that reads a fixed array of integers and prints the
minimum and maximum, using a single loop.

\exercise Fill a \code{i32[3][3]} grid so that \code{grid[r][c] = r * 3 + c}, then
print it as three rows of three numbers.

\exercise Use a vector as a stack: push the numbers 1 to 5, then pop and print
them all, observing that they come out in reverse order. Remember to
\code{free} the vector.

% ===================== Chapter 8 =====================
\chapter{Strings and Characters}
\label{ch:strings}

Text is everywhere in programming: names, messages, file contents, user input.
In Salam a piece of text is a \code{str}, and a single character is a \code{char}.
We have used both since Chapter~1; this chapter gathers everything about them in
one place how to write them, join them, search them, and transform them with
the standard library.

\section{String literals}

A \textbf{string literal} is text between double quotes:

\begin{salam}
func main:
    greeting str = "Hello, Salam!"
    empty str = ""
    println greeting
    println "length:", len(greeting)
end
\end{salam}

\begin{progout}
Hello, Salam!
length: 13
\end{progout}

The built-in \code{len} function reports the number of bytes in a string ---
\code{13} for \code{"Hello, Salam!"}. A string with nothing between the quotes,
\code{""}, is the \emph{empty string}; it is perfectly valid and has length 0.

\section{Strings are immutable}

A \code{str} cannot be changed in place. Every operation that seems to ``modify''
a string actually builds a \emph{new} one and leaves the original untouched. This
is why you can share strings freely without fear that someone will alter one
behind your back, and it is why joining strings always produces a fresh result
rather than growing an existing one.

\section{Joining strings with \texttt{+}}

The \code{+} operator concatenates glues together two strings, producing a
new string:

\begin{salam}
func main:
    first str = "Salam"
    last str = "World"
    full auto = first + ", " + last + "!"
    println full
end
\end{salam}

\begin{progout}
Salam, World!
\end{progout}

Each \code{+} builds a larger string from its two operands. Because the result is
itself a string, you can chain as many as you like, as the example does.

\begin{warn}
\code{+} joins a string to a string. It does \emph{not} automatically turn a
number into text: \code{"count: " + 5} is a type error, because \code{5} is an
\code{i32}, not a \code{str}. To attach a number, convert it first the
standard library's \code{str.FromInt} does exactly that, as we will see shortly.
(When you simply want to \emph{print} mixed values, \code{println} already
accepts a comma-separated mix of any types no conversion needed.)
\end{warn}

\section{Escape sequences}

Some characters are awkward to type directly inside a string a newline, a tab,
or a double quote (which would otherwise end the string). For these you use an
\textbf{escape sequence}: a backslash followed by a code.

\begin{center}
\begin{tabular}{@{}cl@{}}
\toprule
\textbf{Escape} & \textbf{Means} \\
\midrule
\code{\textbackslash n} & newline (move to the next line) \\
\code{\textbackslash t} & tab \\
\code{\textbackslash\textbackslash} & a literal backslash \\
\code{\textbackslash"} & a literal double quote \\
\bottomrule
\end{tabular}
\end{center}

\begin{salam}
func main:
    println "line one\nline two"
    println "name:\tSalam"
    println "she said \"hi\""
    println "a backslash: \\"
end
\end{salam}

\begin{progout}
line one
line two
name:	Salam
she said "hi"
a backslash: \
\end{progout}

The \code{\textbackslash n} split the first string across two lines; the
\code{\textbackslash t} inserted a tab; \code{\textbackslash"} put real quotes
\emph{inside} the text; and \code{\textbackslash\textbackslash} produced a single
backslash.

\section{Triple-quoted strings}

When a string itself contains many double quotes, escaping each one is tedious. A
\textbf{triple-quoted} string, wrapped in \code{"""}, lets you write double
quotes directly without escaping:

\begin{salam}
func main:
    line str = """She said "hi" to me"""
    println line
end
\end{salam}

\begin{progout}
She said "hi" to me
\end{progout}

Triple-quoted strings are handy for text with embedded quotes snippets of
other code, JSON, or dialogue.

\section{Characters}

A \code{char} is a single character, written between \emph{single} quotes. As we
saw in Chapter~2, a \code{char} is really a small number (its byte code), and you
can convert between the two with \code{as}:

\begin{salam}
func main:
    letter char = 'S'
    println letter              // S
    println letter as i32       // 83, the code for 'S'
    println 33 as char          // '!', the character with code 33
end
\end{salam}

\begin{progout}
S
83
!
\end{progout}

\begin{note}
\code{'S'} (single quotes) is a \code{char} one character. \code{"S"} (double
quotes) is a \code{str} text that happens to be one character long. They are
different types and are not interchangeable; keep an eye on which quotes you are
typing.
\end{note}

\section{Strings are not indexed with brackets}

You might expect \code{s[0]} to give the first character of a string, as it would
for an array. In Salam it does \emph{not} a \code{str} is not an array, and
indexing one is an error. To work with parts of a string, you use the functions
of the \code{str} standard-library module, which we turn to now.

\section{The \texttt{str} module}

The built-in \code{+} and \code{len} cover the basics, but most real string work
goes through the \code{str} module. Import it with \code{import "str"} and call
its functions with the \code{str.} prefix. Here are the ones you will use most.

\subsection{Case and trimming}

\begin{salam}
import "str"

func main:
    println str.Upper("salam")        // SALAM
    println str.Lower("SALAM")        // salam
    println str.Trim("   hi   ")      // "hi" (surrounding spaces removed)
    println str.Reverse("Salam")      // malaS
    println str.Repeat("ab", 3)       // ababab
end
\end{salam}

\begin{progout}
SALAM
salam
hi
malaS
ababab
\end{progout}

\subsection{Searching and testing}

\begin{salam}
import "str"

func main:
    println str.Find("hello world", "world")       // 6 (the index)
    println str.Contains("hello", "ell")           // true
    println str.StartsWith("hello.salam", "hello") // true
    println str.EndsWith("hello.salam", ".salam")  // true
end
\end{salam}

\begin{progout}
6
true
true
true
\end{progout}

\code{str.Find} returns the index at which the substring begins (here \code{6},
since \code{"world"} starts at position 6 in \code{"hello world"}), or a negative
value if it is not present.

\subsection{Substrings and replacement}

\begin{salam}
import "str"

func main:
    println str.Substr("programming", 0, 7)    // "program"
    println str.Replace("a-b-c", "-", "+")     // "a+b+c"
end
\end{salam}

\begin{progout}
program
a+b+c
\end{progout}

\code{str.Substr(s, start, n)} extracts \code{n} characters beginning at index
\code{start} so \code{Substr("programming", 0, 7)} takes seven characters from
the start. This is also how you read a single character: \code{str.Substr(s, i,
1)} gives the one-character string at index \code{i}.

\subsection{Converting between strings and numbers}

Turning numbers into text and back is constant work, and the \code{str} module
makes it direct:

\begin{salam}
import "str"

func main:
    n i32 = 2026
    text auto = str.FromInt(n as i64)     // "2026"
    println "Year: " + text               // now '+' works: both are strings

    back auto = str.ToInt("42")           // 42, an i32
    println back + 8                      // 50
end
\end{salam}

\begin{progout}
Year: 2026
50
\end{progout}

\code{str.FromInt} converts an integer to its text form (it takes an \code{i64},
so cast smaller integers with \code{as i64}); \code{str.ToInt} parses a string
back into an \code{i32}. With \code{FromInt} you can finally use \code{+} to
attach a number to a message, because both sides are now strings.

\subsection{Splitting and joining}

Two of the most useful string operations break a string apart and put one back
together. \code{str.Split} cuts a string at every occurrence of a separator,
returning a \code{Vector<str>}; \code{str.Join} does the reverse:

\begin{salam}
import "str"

func main:
    mut parts Vector<str> = str.Split("ada,grace,linus", ",")
    println "count:", parts.len()
    println str.Join(parts, " & ")
    parts.free()
end
\end{salam}

\begin{progout}
count: 3
ada & grace & linus
\end{progout}

Splitting \code{"ada,grace,linus"} on \code{","} yields a vector of three
strings; joining them with \code{" \& "} stitches them back into one. Because
\code{Split} returns a vector, remember to \code{free} it when you are done
(Chapter~7).

\section{A worked example: counting words}

Putting several of these together, here is a small program that counts the words
in a sentence and prints each on its own line:

\begin{salam}
import "str"

func main:
    sentence str = "the quick brown fox"
    mut words Vector<str> = str.Split(sentence, " ")
    println "word count:", words.len()
    mut i i32 = 0
    while i < words.len():
        println str.Upper(words.get(i)[0])
        i += 1
    end
    words.free()
end
\end{salam}

\begin{progout}
word count: 4
THE
QUICK
BROWN
FOX
\end{progout}

\begin{deep}[In Depth: bytes, characters, and Unicode]
A Salam \code{str} is stored as a sequence of bytes in UTF-8 encoding, and
\code{len} reports the number of \emph{bytes}. For plain English text, one byte
is one character, so the distinction never bites. But a character outside the
basic Latin range an accented letter, or a letter of another script may
occupy several bytes, so its byte length and its visible length differ. The
\code{str} module's higher-level functions (\code{Split}, \code{Find},
\code{Replace}, and friends) are the safe way to manipulate text; reach for raw
byte arithmetic only when you know your data is plain ASCII.
\end{deep}

\section{Chapter summary}

\begin{itemize}[leftmargin=1.4em]
  \item A \code{str} is immutable text in double quotes; \code{len(s)} gives its
        byte length; \code{+} concatenates.
  \item Escape sequences (\code{\textbackslash n}, \code{\textbackslash t},
        \code{\textbackslash"}, \code{\textbackslash\textbackslash}) embed
        special characters; \code{"""..."""} avoids escaping inner quotes.
  \item A \code{char} is a single character in single quotes and converts to/from
        its numeric code with \code{as}.
  \item Strings are \emph{not} indexed with \code{[]}; use the \code{str} module.
  \item \code{str} provides \code{Upper}, \code{Lower}, \code{Trim},
        \code{Reverse}, \code{Repeat}, \code{Find}, \code{Contains},
        \code{StartsWith}, \code{EndsWith}, \code{Substr}, \code{Replace},
        \code{FromInt}, \code{ToInt}, \code{Split}, and \code{Join}.
\end{itemize}

\section{Exercises}

\exercise Ask for a \code{str} variable holding your full name and print its
length, its uppercase form, and its reverse.

\exercise Build the string \code{"Item No.7 costs \$42"} by converting the numbers
7 and 42 with \code{str.FromInt} and joining with \code{+}.

\exercise Given \code{"one,two,three,four"}, split on the comma, print how many
parts there are, and print them re-joined with \code{" - "}.

\exercise Write a function \code{is\_palindrome(s str) bool} that returns whether
a string reads the same forwards and backwards (hint: compare \code{s} with
\code{str.Reverse(s)}).

\exercise Use \code{str.Substr} to print each character of \code{"Salam"} on its
own line.

\exercise Write a program that, given a filename like \code{"report.txt"}, prints
just the extension (\code{"txt"}) by finding the last \code{"."}.

% ===================== Chapter 9 =====================
\chapter{Structs and Methods}
\label{ch:structs-methods}

So far our data has been loose: an \code{i32} here, a \code{str} there. But the
things we model are rarely single values. A bank account has an owner \emph{and} a
balance; a point has an $x$ \emph{and} a $y$. A \textbf{struct} lets you bundle
several values into one named type, and \textbf{methods} let you attach behaviour
to that type. Together they are how you build your own data types in Salam.

\section{Defining a struct}

A struct is declared with the \code{struct} keyword, a name, and a list of
\textbf{fields} each a name, a type, and a default value closed with
\code{end}:

\begin{salam}
struct Point:
    x i32 = 0
    y i32 = 0
end
\end{salam}

This defines a new type, \code{Point}, with two fields. Each field has a default
value (here \code{0}), which is used when you do not supply one. \code{Point} is
now a type you can use anywhere a type is expected, just like \code{i32} or
\code{str}.

\section{Creating instances}

To make a value of a struct type an \textbf{instance} write the type name
followed by braces listing the fields you want to set:

\begin{salam}
struct Point:
    x i32 = 0
    y i32 = 0
end

func main:
    origin Point = Point {}                 // both fields default to 0
    p Point = Point { x = 3, y = 4 }        // set both fields
    println origin.x, origin.y
    println p.x, p.y
end
\end{salam}

\begin{progout}
0 0
3 4
\end{progout}

Inside the braces you assign fields by name with \code{field = value}. Any field
you leave out keeps its default \code{Point \{\}} relies on the defaults for
both, giving \code{(0, 0)}. You reach a field of an instance with a dot:
\code{p.x} is \code{p}'s \code{x} field.

\begin{note}
Because every field has a default, \code{Point \{\}} is always a valid, fully
formed value there is no such thing as a half-built struct with missing
fields. This is why field defaults are required in the struct definition: they
guarantee that an instance is never left with an uninitialised field.
\end{note}

\section{Reading and changing fields}

Field access with the dot works both ways: you can read a field, and if the
instance is \code{mut} you can assign to one:

\begin{salam}
struct Point:
    x i32 = 0
    y i32 = 0
end

func main:
    mut p Point = Point { x = 1, y = 1 }
    p.x = 10
    p.y = p.y + 5
    println p.x, p.y
end
\end{salam}

\begin{progout}
10 6
\end{progout}

\section{Structs are copied by value}

A struct behaves like the simple values of Chapter~2: assigning it, passing it to
a function, or returning it makes a \emph{copy}. Changing the copy does not affect
the original:

\begin{salam}
struct Point:
    x i32 = 0
    y i32 = 0
end

func main:
    mut p Point = Point { x = 1, y = 2 }
    mut q Point = p        // q is a copy of p
    q.x = 99
    println p.x, q.x       // p is untouched
end
\end{salam}

\begin{progout}
1 99
\end{progout}

This ``value semantics'' makes structs predictable: a \code{Point} you hand to a
function cannot be quietly modified by that function. It is the same guarantee you
saw for ordinary arguments in Chapter~6, now extended to your own types.

\section{Methods: behaviour on a struct}

A struct can carry not just data but \textbf{methods} functions that belong to
the type and act on a particular instance. You write a method inside the
\code{struct} block, exactly like a function, and refer to the current instance
as \code{this}:

\begin{salam}
struct Counter:
    value i32 = 0
    func increment():
        this.value = this.value + 1
    end
    func get() i32:
        ret this.value
    end
end

func main:
    mut c Counter = Counter {}
    c.increment()
    c.increment()
    c.increment()
    println c.get()
end
\end{salam}

\begin{progout}
3
\end{progout}

You \emph{call} a method through an instance with the dot: \code{c.increment()}.
Inside the method, \code{this} refers to the instance the method was called on ---
so \code{this.value} is \emph{that} counter's \code{value}. Each \code{Counter}
keeps its own \code{value}, and a method always works on the one it was called
through.

\begin{note}
Unlike an ordinary argument, a method \emph{can} change the instance it is called
on: \code{increment} really does update \code{c}'s \code{value}. This is the
sanctioned way to modify data in place that we hinted at in Chapter~6 a
struct's own methods are trusted to change it, while outside functions get
copies.
\end{note}

\section{Methods that compute and methods that act}

Just like free functions, a method may return a value or return nothing. A method
that returns a value is written with a return type after its parameter list:

\begin{salam}
struct Rectangle:
    width i32 = 0
    height i32 = 0
    func area() i32:
        ret this.width * this.height
    end
    func is_square() bool:
        ret this.width == this.height
    end
end

func main:
    r Rectangle = Rectangle { width = 4, height = 4 }
    println "area:", r.area()
    println "square?", r.is_square()
end
\end{salam}

\begin{progout}
area: 16
square? true
\end{progout}

Here \code{area} and \code{is\_square} only \emph{read} the rectangle (through
\code{this}) and compute an answer; they do not change anything. A method may also
take parameters of its own, in addition to the implicit \code{this}:

\begin{salam}
struct Rectangle:
    width i32 = 0
    height i32 = 0
    func scale(factor i32):
        this.width = this.width * factor
        this.height = this.height * factor
    end
end

func main:
    mut r Rectangle = Rectangle { width = 2, height = 3 }
    r.scale(5)
    println r.width, r.height
end
\end{salam}

\begin{progout}
10 15
\end{progout}

\section{A fuller example: a bank account}

Let us combine fields and methods into something realistic a bank account that
guards its own balance. The \code{withdraw} method returns a \code{bool} reporting
whether the withdrawal succeeded:

\begin{salam}
import "str"

struct Account:
    owner str = ""
    balance i32 = 0
    func deposit(amount i32):
        this.balance = this.balance + amount
    end
    func withdraw(amount i32) bool:
        if amount > this.balance:
            ret false
        end
        this.balance = this.balance - amount
        ret true
    end
    func describe() str:
        ret this.owner + " has " + str.FromInt(this.balance as i64)
    end
end

func main:
    mut acc Account = Account { owner = "Sara", balance = 100 }
    acc.deposit(50)
    println acc.describe()
    println "withdraw 30?", acc.withdraw(30)
    println "withdraw 1000?", acc.withdraw(1000)
    println acc.describe()
end
\end{salam}

\begin{progout}
Sara has 150
withdraw 30? true
withdraw 1000? false
Sara has 120
\end{progout}

This is the heart of good design: the \code{Account} \emph{owns} its balance, and
the only way to change it is through \code{deposit} and \code{withdraw}, which
enforce the rules (you cannot withdraw more than you have). The data and the
operations that protect it live together in one place.

\section{Structs within structs}

A field may itself be a struct, letting you build larger structures from smaller
ones. You reach a nested field by chaining dots:

\begin{salam}
struct Point:
    x i32 = 0
    y i32 = 0
end

struct Line:
    start Point = Point {}
    finish Point = Point {}
end

func main:
    ln Line = Line {
        start = Point { x = 0, y = 0 },
        finish = Point { x = 5, y = 5 }
    }
    println ln.start.x, ln.finish.y
end
\end{salam}

\begin{progout}
0 5
\end{progout}

\code{ln.finish.y} reads the \code{y} field of the \code{finish} field of
\code{ln}. Composing structs this way is how you model anything with structure ---
a line from two points, a person with an address, a game piece with a position
and a colour.

\begin{deep}[In Depth: structs become C structs]
Each Salam struct compiles to a plain C \code{struct} with the same fields in the
same order, and a method compiles to an ordinary C function that takes a pointer
to the struct as a hidden first argument which is exactly what \code{this} is.
Because the layout is a direct, flat C struct, accessing a field is a single
memory read with no indirection, and passing a struct by value is a simple
memory copy. This is why structs are both convenient \emph{and} fast, and why
they map so cleanly onto C libraries through the foreign function interface
(Chapter~19).
\end{deep}

\section{Chapter summary}

\begin{itemize}[leftmargin=1.4em]
  \item A \code{struct} bundles named fields, each with a default value, into a
        new type.
  \item Create an instance with \code{Type \{ field = value, ... \}}; omitted
        fields use their defaults.
  \item Read and write fields with the dot: \code{p.x}.
  \item Structs are copied by value on assignment, argument passing, and return.
  \item Methods live inside the \code{struct} block, refer to the instance as
        \code{this}, and are called with the dot: \code{c.increment()}.
  \item A method may read or modify \code{this}, take its own parameters, and
        return a value or nothing.
  \item Struct fields may themselves be structs; chain dots to reach nested
        fields.
\end{itemize}

\section{Exercises}

\exercise Define a \code{struct Circle} with a \code{radius f64} field and a
method \code{area() f64} that returns $\pi r^2$ (use \code{3.14159} for $\pi$).
Print the area of a circle of radius 2.

\exercise Add a \code{describe()} method to \code{Point} that returns a string
like \code{"(3, 4)"}, and print it for a couple of points.

\exercise Build a \code{struct Stack} around a fixed \code{i32[16]} array and a
\code{count}, with \code{push} and \code{pop} methods. Push three values and pop
them back.

\exercise Extend the \code{Account} example with a \code{transfer(other, amount)}
idea: withdraw from one account and deposit into another, only if the withdrawal
succeeds.

\exercise Define a \code{struct Rectangle} with a \code{perimeter()} method as well
as \code{area()}, and print both for a $3 \times 5$ rectangle.

\exercise Model a \code{struct Person} containing a \code{name} and an
\code{Address} struct (street, city). Create one and print the city via chained
field access.

\chapter{Enums and Pattern Matching}
\label{ch:enums-matching}

Many values in a program are not really numbers or text they are a choice from
a small, fixed set of possibilities. A day is one of seven; a traffic light is
red, amber, or green; an order is pending, shipped, or delivered. An
\textbf{enum} (short for \emph{enumeration}) lets you name such a set and use those
names directly, instead of remembering that ``3 means Thursday.''

\section{Defining an enum}

An enum is declared with the \code{enum} keyword, a name, and a list of
\textbf{members} the allowed values closed with \code{end}. The members may
be written on one line, separated by commas, or one per line:

\begin{salam}
enum Day: Mon, Tue, Wed, Thu, Fri, Sat, Sun end
\end{salam}

This defines a new type, \code{Day}, whose only values are the seven named days.
A variable of type \code{Day} can hold exactly one of them nothing else. That
is the whole point: the type makes invalid values \emph{unrepresentable}.

\section{Using enum values}

You refer to a member by qualifying it with the enum's name and a dot:
\code{Day.Mon}, \code{Day.Sat}. This makes it clear which enum the value belongs
to and avoids name clashes between enums that happen to share a member name.

\begin{salam}
enum Day: Mon, Tue, Wed, Thu, Fri, Sat, Sun end

func main:
    today Day = Day.Sat
    if today == Day.Sat || today == Day.Sun:
        println "It's the weekend!"
    else:
        println "A working day."
    end
end
\end{salam}

\begin{progout}
It's the weekend!
\end{progout}

You compare enum values with \code{==} and \code{!=}, exactly as you would
numbers. The variable \code{today} can only ever be one of the seven days, so the
compiler and the next person to read your code knows precisely what values
are in play.

\section{Enums are numbers underneath}

Each enum member has an integer value. By default they count up from zero in
order: in \code{Day}, \code{Mon} is 0, \code{Tue} is 1, and so on to \code{Sun}
at 6. You can see a member's number by converting it to an integer with \code{as}:

\begin{salam}
enum Day: Mon, Tue, Wed, Thu, Fri, Sat, Sun end

func main:
    println Day.Mon as i32      // 0
    println Day.Sat as i32      // 5
end
\end{salam}

\begin{progout}
0
5
\end{progout}

\subsection{Choosing the numbers}

When the underlying numbers matter because they correspond to codes, flags, or
external values you can assign them explicitly. Any member you do not give a
number to continues counting up from the previous one:

\begin{salam}
enum Color:
    Red            // 0
    Green = 5      // 5
    Blue           // 6 (continues from 5)
end

func main:
    println Color.Red as i32, Color.Green as i32, Color.Blue as i32
end
\end{salam}

\begin{progout}
0 5 6
\end{progout}

\code{Red} keeps the default 0, \code{Green} is pinned to 5, and \code{Blue}
follows on at 6. This is handy when an enum mirrors values defined elsewhere, such
as HTTP status codes or hardware register values.

\begin{warn}
The conversion goes one way freely: an enum value turns into its integer with
\code{as}. Going the other way treating an arbitrary integer as an enum is
something to do only with care, because a number might not correspond to any
member. Prefer to keep values in their enum form and convert to integers only at
the edges of your program, where you must.
\end{warn}

\section{Pattern matching with \texttt{if}/\texttt{else if}}

Some languages provide a dedicated \code{match} or \code{switch} statement for
choosing an action based on which enum value you have. Salam keeps things simple:
you ``match'' on an enum with an ordinary \code{if}/\code{else if} chain, which
you already know well from Chapter~4.

\begin{salam}
enum Light: Red, Amber, Green end

func action(light Light) str:
    if light == Light.Red:
        ret "Stop"
    else if light == Light.Amber:
        ret "Slow down"
    else:
        ret "Go"
    end
end

func main:
    println action(Light.Red)
    println action(Light.Amber)
    println action(Light.Green)
end
\end{salam}

\begin{progout}
Stop
Slow down
Go
\end{progout}

The function takes a \code{Light} and returns the right instruction for it. The
\code{else if} chain considers each possibility in turn; the final \code{else}
catches the last case (\code{Green}). This pattern one branch per enum
member is the Salam way of ``matching'' on an enum.

\begin{tip}
When you write such a chain, cover \emph{every} member, and let the final
\code{else} handle the last one (or an unexpected value) rather than leaving a
case out. If you later add a member to the enum, revisit these chains to handle
it a forgotten case is a classic source of bugs in any language.
\end{tip}

\section{Enums as state machines}

Enums shine when a value moves through a sequence of well-defined states. Consider
an order that progresses from \code{Pending} to \code{Shipped} to
\code{Delivered}. An enum names the states, and a function decides the next one:

\begin{salam}
enum Status: Pending, Shipped, Delivered end

func next(s Status) Status:
    if s == Status.Pending:
        ret Status.Shipped
    else if s == Status.Shipped:
        ret Status.Delivered
    else:
        ret Status.Delivered      // already final; stays put
    end
end

func label(s Status) str:
    if s == Status.Pending:   ret "pending"
    else if s == Status.Shipped: ret "shipped"
    else:                     ret "delivered"
    end
end

func main:
    mut s Status = Status.Pending
    println label(s)
    s = next(s)
    println label(s)
    s = next(s)
    println label(s)
end
\end{salam}

\begin{progout}
pending
shipped
delivered
\end{progout}

The enum guarantees that \code{s} is always a valid status, and the \code{next}
function encodes the legal transitions in one place. This is a small example of a
\textbf{state machine} a powerful and common way to model anything that moves
through stages, from a game's screens to a network connection's lifecycle.

\begin{deep}[In Depth: why an enum beats a bare integer]
You could represent a status with a plain \code{i32} 0 for pending, 1 for
shipped, and so on and many programs do. But then nothing stops you from
assigning \code{99}, or from adding a \code{Status} to a \code{Day} by accident,
or from forgetting which number meant what. An enum makes those mistakes
impossible: the type system knows the only valid values, the names document
themselves, and the compiler checks that you use them consistently. The cost is
nothing an enum compiles down to exactly the integer you would have written by
hand so you get the safety for free.
\end{deep}

\section{Chapter summary}

\begin{itemize}[leftmargin=1.4em]
  \item An \code{enum} names a fixed set of values; a variable of that type holds
        exactly one member.
  \item Refer to members as \code{EnumName.Member}; compare them with \code{==}
        and \code{!=}.
  \item Members have integer values, counting from 0 by default; convert with
        \code{as i32}. You may assign explicit values, and later members continue
        counting up.
  \item Salam has no \code{match}/\code{switch}; ``match'' on an enum with an
        \code{if}/\code{else if} chain that covers every member.
  \item Enums are ideal for modelling states and transitions (state machines),
        and are safer than bare integers at no run-time cost.
\end{itemize}

\section{Exercises}

\exercise Define \code{enum Suit: Hearts, Diamonds, Clubs, Spades end} and print
the integer value of each member.

\exercise Write a function \code{is\_red(s Suit) bool} that returns \code{true} for
\code{Hearts} and \code{Diamonds}, and test it on all four suits.

\exercise Define \code{enum Direction: North, East, South, West end} and a function
\code{turn\_right(d Direction) Direction} that rotates clockwise (North becomes
East, and so on, West wrapping back to North).

\exercise Model a traffic light that cycles \code{Green} $\to$ \code{Amber} $\to$
\code{Red} $\to$ \code{Green}. Write a \code{next} function and run it through two
full cycles.

\exercise Define \code{enum HttpStatus} with \code{Ok = 200},
\code{NotFound = 404}, and \code{ServerError = 500}, and print each member's
number.

% ===================== Chapter 11 =====================
\chapter{Interfaces and Polymorphism}
\label{ch:interfaces}

A \code{Circle} and a \code{Rect} are different structs, with different fields,
yet both have an \emph{area}. Often you want to write code that works with
``anything that has an area'' without caring exactly which shape it is. An
\textbf{interface} captures that idea: it describes a set of methods a type must
provide, and lets you write code against the interface rather than any one
concrete type. This ability to treat many types through one shared description is
called \textbf{polymorphism}.

\section{Defining an interface}

An interface is declared with the \code{interface} keyword and a list of method
\emph{signatures} method names with their parameters and return types, but no
bodies:

\begin{salam}
interface Shape:
    func area() f64
    func name() str
end
\end{salam}

This says: ``a \code{Shape} is anything that has an \code{area()} returning an
\code{f64} and a \code{name()} returning a \code{str}.'' The interface promises
what methods exist; it does not say how they work that is each type's job.

\section{Implementing an interface}

Here is the pleasant part: in Salam a struct implements an interface simply by
\emph{having the right methods}. There is no \code{implements} clause to write and
nothing to declare if a struct has an \code{area()} and a \code{name()} with
matching signatures, it already \emph{is} a \code{Shape}. This is called
\textbf{structural conformance}.

\begin{salam}
interface Shape:
    func area() f64
    func name() str
end

struct Circle:
    r f64 = 0.0
    func area() f64: ret 3.14159 * this.r * this.r end
    func name() str: ret "circle" end
end

struct Rect:
    w f64 = 0.0
    h f64 = 0.0
    func area() f64: ret this.w * this.h end
    func name() str: ret "rect" end
end
\end{salam}

Both \code{Circle} and \code{Rect} provide \code{area()} and \code{name()}, so
both conform to \code{Shape} automatically, without either struct mentioning
the interface. Now we can write code that works with any \code{Shape}.

\section{Static polymorphism: interface bounds}

The first and most efficient way to use an interface is as a \textbf{bound} on a
generic function. The notation \code{<T: Shape>} after a function name means ``for
any type \code{T} that conforms to \code{Shape}.'' Inside the function you may
call any of the interface's methods on a value of type \code{T}:

\begin{salam}
func describe<T: Shape>(s T):
    println s.name(), "area =", s.area()
end

func main:
    describe(Circle { r = 2.0 })
    describe(Rect { w = 3.0, h = 4.0 })
end
\end{salam}

\begin{progout}
circle area = 12.5664
rect area = 12
\end{progout}

The one \code{describe} function serves both shapes. You did not have to write it
twice, and you did not have to tell it which type you are passing the compiler
works that out from the argument. Notice there are no angle brackets at the call
site: you just write \code{describe(Circle \{ r = 2.0 \})}, and Salam infers that
\code{T} is \code{Circle}.

\begin{deep}[In Depth: monomorphization, or ``zero-cost'' generics]
When you call \code{describe(Circle \{...\})}, the compiler generates a
specialised copy of \code{describe} in which \code{T} is exactly \code{Circle},
and another copy for \code{Rect}. Each copy calls the methods directly, with no
run-time lookup it is as fast as if you had written two separate functions by
hand. This technique, called \emph{monomorphization}, is why interface bounds are
described as ``zero-cost'': you get the convenience of writing one generic
function and the speed of hand-specialised code. The price is a little extra
compiled size, one copy per type you actually use.
\end{deep}

\section{Dynamic polymorphism: \texttt{dyn}}

Interface bounds are resolved at compile time, which means each call site knows
the concrete type. But sometimes you want a \emph{single} collection holding a
mixture of shapes some circles, some rectangles and to call \code{area()}
on each without knowing which is which until the program runs. For that, Salam
offers \code{dyn}.

A value of type \code{dyn Shape} is ``some shape, decided at run time.'' You can
store different concrete types in one place and call interface methods on them,
and Salam dispatches each call to the right type's method dynamically:

\begin{salam}
func describe(s dyn Shape):
    println s.name(), "area =", s.area()
end

func main:
    describe(Circle { r = 1.0 })
    describe(Rect { w = 2.0, h = 3.0 })
end
\end{salam}

\begin{progout}
circle area = 3.14159
rect area = 6
\end{progout}

The difference from the bounded version is what happens under the hood: a
\code{dyn Shape} carries, alongside the data, a hidden table of which methods to
call, so the right \code{area()} is chosen at run time. This is exactly what lets
a single list hold a heterogeneous mix of shapes.

\begin{note}
\textbf{When to use which.} Reach for an interface \emph{bound} (\code{<T:
Shape>}) by default it is faster and is the right tool when each call works
with one known type. Reach for \code{dyn} when you genuinely need to mix different
concrete types in one collection or behind one variable, and decide the method at
run time. Bounds are about writing one function for many types; \code{dyn} is
about holding many types in one container.
\end{note}

\section{Designing with interfaces}

Interfaces let you program to a \emph{capability} rather than a concrete type,
which keeps your code flexible. A function that sorts ``anything comparable,'' a
logger that writes to ``anything writable,'' a renderer that draws ``any
shape'' each is expressed by naming the interface it needs and ignoring
everything else about its inputs. When you later add a new type that conforms, the
existing code works with it unchanged, because it only ever depended on the
methods, never the specific struct.

\begin{tip}
Keep interfaces small. An interface with one or two methods (``has an area,''
``can be drawn'') is easy for many types to satisfy and easy to reason about. A
large interface demanding a dozen methods is a burden to implement and tends to
couple your code to one specific design. Prefer several small interfaces over one
big one.
\end{tip}

\section{Chapter summary}

\begin{itemize}[leftmargin=1.4em]
  \item An \code{interface} lists method signatures (no bodies) that a conforming
        type must provide.
  \item Conformance is \emph{structural}: a struct implements an interface just by
        having the matching methods no \code{implements} declaration.
  \item An interface \emph{bound}, \code{func f<T: Iface>(x T)}, writes one
        function usable with any conforming type; the type is inferred at the call
        site and resolved at compile time (monomorphization, zero run-time cost).
  \item \code{dyn Iface} holds ``some conforming type, decided at run time,''
        enabling heterogeneous collections via dynamic dispatch.
  \item Use bounds by default; use \code{dyn} when you must mix concrete types in
        one place.
  \item Favour small interfaces that many types can easily satisfy.
\end{itemize}

\section{Exercises}

\exercise Define an interface \code{Greeter} with a method \code{greet() str}, and
two structs (say \code{Formal} and \code{Casual}) that conform. Write a bounded
function that prints any greeter's greeting.

\exercise Add a third shape, \code{Triangle} (base and height), to the
\code{Shape} example. Confirm the existing \code{describe} works on it with no
changes.

\exercise Give \code{Shape} a third method \code{perimeter() f64} and implement it
for \code{Circle} and \code{Rect}. What error do you get if you forget one?

\exercise Write a bounded function \code{bigger<T: Shape>(a T, b T) str} that
returns the name of whichever shape has the larger area. (Both arguments are the
same type here why?)

\exercise Explain, in a short comment, when you would choose \code{dyn Shape} over
\code{<T: Shape>}, using a list of mixed shapes as your example.

\chapter{Generics}
\label{ch:generics}

You met \code{Vector<int>} back in Chapter~7 and the bound \code{<T: Shape>} in
Chapter~11. Both rely on \textbf{generics}: the ability to write code that works
for \emph{any} type, with the type filled in later. Generics let you write a
function or a data structure once and reuse it for integers, strings, your own
structs anything without copying and pasting. This chapter brings the idea
into full focus.

\section{The problem generics solve}

Suppose you want a function that returns the larger of two values. For integers:

\begin{salam}
func max_int(a : int, b : int) : int:
    if a > b: ret a end
    ret b
end
\end{salam}

For floating-point numbers you would write \emph{the same logic again} with
\code{f64} in place of \code{int}, and again for any other comparable type. That
is wasteful and error-prone a bug fixed in one copy lingers in the others.
Generics let you write the logic a single time.

\section{Generic functions}

A \textbf{type parameter} in angle brackets after the function name turns a
function generic. The parameter conventionally \code{T} stands for ``some
type, chosen by the caller,'' and you use it like any other type within the
function:

\begin{salam}
func Max<T>(a : T, b : T) : T:
    if a > b: ret a end
    ret b
end

func main:
    println Max(10, 20)        // T is int
    println Max(2.5, 1.5)      // T is f64
end
\end{salam}

\begin{progout}
20
2.5
\end{progout}

One \code{Max} now serves every type. The caller writes \code{Max(10, 20)} with
no angle brackets; Salam looks at the arguments and \emph{infers} that \code{T} is
\code{int} for the first call and \code{f64} for the second. The two parameters
\code{a} and \code{b} share the type \code{T}, so they must be the \emph{same}
type at each call you cannot write \code{Max(10, 2.5)}, because \code{T} cannot
be both \code{int} and \code{f64} at once.

\section{Generic structs}

Data structures benefit even more from generics. A ``box that holds one value''
should not care whether that value is an integer or a string. Declare a type
parameter after the struct's name and use it for fields:

\begin{salam}
struct Box<T>:
    value : T
end

func main:
    mut a : Box<int> = Box { value = 42 }
    mut b : Box<str> = Box { value = "hi" }
    println a.value
    println b.value
end
\end{salam}

\begin{progout}
42
hi
\end{progout}

\code{Box<int>} is a box of integers; \code{Box<str>} is a box of strings. You
spell out the concrete type in the variable's declaration \code{Box<int>} ---
and from then on the field \code{value} has that type. This is precisely how
\code{Vector<T>} works: it is a generic struct holding a growable buffer of
\code{T}.

\begin{note}
At the \emph{call site} of a generic function, you usually omit the type ---
\code{Max(10, 20)} because Salam infers it from the arguments. But when you
declare a variable of a generic \emph{struct} type, you write the concrete type
explicitly: \code{Box<int>}, \code{Vector<str>}. The struct's type parameter
cannot always be inferred from an empty \code{Box \{\}}, so you state it.
\end{note}

\section{A generic container: a simple stack}

Let us build something real: a generic stack, a last-in-first-out store, on top of
a fixed array. Because it is generic, the same code gives you a stack of
integers, a stack of strings, or a stack of anything:

\begin{salam}
struct Stack<T>:
    items : T[16] = [0,0,0,0,0,0,0,0,0,0,0,0,0,0,0,0]
    count : int = 0
    func push(x : T):
        this.items[this.count] = x
        this.count = this.count + 1
    end
    func pop() : T:
        this.count = this.count - 1
        ret this.items[this.count]
    end
    func is_empty() : bool:
        ret this.count == 0
    end
end

func main:
    mut s : Stack<int> = Stack {}
    s.push(3)
    s.push(7)
    s.push(1)
    println s.pop()
    println s.pop()
    println s.is_empty()
end
\end{salam}

\begin{progout}
1
7
false
\end{progout}

The methods \code{push} and \code{pop} are written entirely in terms of \code{T}.
Declare \code{Stack<int>} and you get an integer stack; declare \code{Stack<str>}
and the very same definition gives you a string stack. This is how the standard
library's collections are built, and how you build your own.

\begin{note}
The \code{items} field carries a default (sixteen zeros) so that \code{Stack
\{\}} is a complete value with every field initialized the same rule from
Chapter~9 that every struct field needs a default if you want to construct the
struct with empty braces. For a real, growable container you would build on
\code{Vector<T>} (Chapter~7) instead of a fixed array; this small fixed-size
stack keeps the focus on the generics.
\end{note}

\section{Bounded generics revisited}

A bare \code{<T>} accepts \emph{any} type, but then you can only do to a \code{T}
what every type supports. The moment you call a method on a \code{T} ---
\code{s.area()}, say the compiler must know that \code{T} actually \emph{has}
that method. That is what an interface \textbf{bound} (Chapter~11) provides:
\code{<T: Shape>} means ``any \code{T}, as long as it conforms to \code{Shape},''
which unlocks the interface's methods inside the function.

\begin{salam}
interface Shape:
    func area() : f64
end

func total_area<T: Shape>(a : T, b : T) : f64:
    ret a.area() + b.area()
end
\end{salam}

So there is a spectrum: \code{<T>} for code that treats values opaquely (store,
move, return them), and \code{<T: SomeInterface>} for code that needs to
\emph{do} something specific with them. You reach for the bound exactly when you
need to call the interface's methods.

\section{How generics work: monomorphization}

\begin{deep}[In Depth: one template, many concrete copies]
Salam implements generics by \emph{monomorphization}, the same mechanism behind
interface bounds. The generic definition \code{Max<T>} or \code{Stack<T>} ---
is a \emph{template}; the compiler never emits it directly. Instead, each time you
use it with a concrete type, the compiler stamps out a specialised copy with the
type filled in: \code{Max} becomes a real \code{Max\_i32} and a real
\code{Max\_f64} (the compiler names the copies with the explicit type names);
\code{Stack<int>} becomes a concrete C struct with an \code{int[16]} field. Those
copies are ordinary, fully type-checked code, as fast
as anything hand-written there is no boxing, no run-time type tag, no
indirection. The trade-off is compiled size: one copy exists for each distinct
type you instantiate. In practice that is a small price for code that is both
reusable and fast.

To make this possible, the compiler needs to know an element's size at compile
time. The built-in \code{sizeof(T)} gives it: the standard library's
\code{Vector<T>}, for instance, uses \code{sizeof(T)} to allocate exactly enough
bytes for each element when it grows its buffer. You rarely call \code{sizeof}
yourself, but it is the quiet machinery that lets one generic container hold
elements of any size.
\end{deep}

\section{When to make something generic}

Generics are powerful, but not everything needs them. A good rule:

\begin{itemize}[leftmargin=1.4em]
  \item Make a \emph{container} generic almost always a stack, queue, list, or
        tree is naturally indifferent to what it stores.
  \item Make a \emph{function} generic when you find yourself about to copy it for
        a second type. The first copy is fine; the second is the signal.
  \item Do \emph{not} reach for generics when a single concrete type is all you
        will ever need the extra abstraction earns its keep only when there is
        real reuse.
\end{itemize}

\section{Chapter summary}

\begin{itemize}[leftmargin=1.4em]
  \item Generics write code once for \emph{any} type, using a type parameter
        \code{<T>}.
  \item A generic function infers \code{T} from its arguments at the call site;
        parameters sharing \code{T} must be the same type.
  \item A generic struct (\code{Box<T>}, \code{Stack<T>}, \code{Vector<T>}) is
        declared with its concrete type, e.g. \code{Stack<int>}.
  \item A bare \code{<T>} treats values opaquely; an interface bound \code{<T:
        Iface>} lets you call the interface's methods on a \code{T}.
  \item Generics compile by \emph{monomorphization} into specialised, fast copies;
        \code{sizeof(T)} is the machinery that lets a container size its elements.
\end{itemize}

\section{Exercises}

\exercise Write a generic function \code{Min<T>(a T, b T) T} mirroring \code{Max},
and test it on integers and floats.

\exercise Write a generic function \code{first\_of<T>(a T, b T) T} that simply
returns its first argument, and explain why it needs no bound.

\exercise Define a generic \code{struct Pair<T>} holding two values of type
\code{T} and a method \code{swap()} that exchanges them. Test it with a pair of
integers.

\exercise Extend the generic \code{Stack} with a \code{peek()} method that returns
the top element without removing it. Test it with a \code{Stack<str>}.

\exercise Explain in a comment why \code{Max(10, 2.5)} fails to compile, in terms
of the type parameter \code{T}.

% ===================== Chapter 13 =====================
\chapter{Closures and Lambdas}
\label{ch:closures-lambdas}

So far functions have been things we \emph{define} and \emph{call}. In this
chapter functions become \emph{values}: things we can store in a variable, pass to
another function, and create on the spot without giving them a name. A small,
unnamed function is called a \textbf{lambda}, and when one captures variables from
around it, it becomes a \textbf{closure}. These ideas unlock a flexible,
expressive style of programming.

\section{Functions as values}

A \textbf{lambda} is a function written as an expression, with no name. Its syntax
is a parameter list, the arrow \code{=>}, and a body. For a one-expression body,
just write the expression:

\begin{salam}
func main:
    double auto = (x i32) => x * 2
    println double(21)
end
\end{salam}

\begin{progout}
42
\end{progout}

Here \code{(x i32) => x * 2} is a function value: it takes an \code{i32} called
\code{x} and produces \code{x * 2}. We store it in the variable \code{double}
(letting \code{auto} infer its type) and then \emph{call} it just like any
function: \code{double(21)}. The lambda has no name of its own the name
\code{double} belongs to the variable, not the function.

\subsection{The type of a function value}

The type of \code{double} is \code{func(i32) i32}: ``a function taking an
\code{i32} and returning an \code{i32}.'' In general a function-value type is
written \code{func(}\textit{argument types}\code{) }\textit{return type}. A
function value taking nothing and returning an \code{i32} has type
\code{func() i32}. You will see these types most often where one function accepts
another as a parameter.

\subsection{Multi-statement bodies}

When a lambda needs more than one statement, give it a block body: a colon after
the arrow, the statements, and \code{end} exactly like a named function, just
without the name:

\begin{salam}
(x i32) => :
    mut y i32 = x * x
    ret y + 1
end
\end{salam}

The single-expression form (\code{=> x * 2}) is shorthand for the common case; the
block form (\code{=> : ... end}) is there when you need it.

\section{Higher-order functions}

A function that takes another function as a parameter or returns one is
called a \textbf{higher-order function}. Because a function value has a type like
\code{func(i32) i32}, you can declare a parameter of that type and call it inside:

\begin{salam}
func apply(f func(i32) i32, x i32) i32:
    ret f(x)
end

func main:
    println apply((x i32) => x + 100, 5)
end
\end{salam}

\begin{progout}
105
\end{progout}

\code{apply} does not know or care what \code{f} does; it just calls \code{f(x)}.
We hand it a lambda created right at the call site, \code{(x i32) => x + 100}, and
\code{apply} runs it on \code{5}. This is the essence of higher-order
programming: behaviour passed around as data. The same \code{apply} could square,
negate, or double its argument, depending only on the lambda you give it.

\section{Closures: capturing the surrounding state}

Here is where function values become truly powerful. A lambda can \emph{use
variables from the scope in which it was created} not just its own parameters.
When it does, it ``closes over'' those variables, and we call it a
\textbf{closure}. The captured values travel with the function, even after the
scope that created them has returned.

The classic example is a function that \emph{builds} another function:

\begin{salam}
func make_adder(n i32) func(i32) i32:
    ret (x i32) => x + n
end

func main:
    add10 auto = make_adder(10)
    add100 auto = make_adder(100)
    println add10(5)
    println add100(5)
end
\end{salam}

\begin{progout}
15
105
\end{progout}

\code{make\_adder} returns a lambda that adds \code{n} but \code{n} is
\code{make\_adder}'s parameter, not the lambda's. The returned closure
\emph{captures} the value of \code{n}. So \code{make\_adder(10)} produces a
function that always adds 10, and \code{make\_adder(100)} produces a separate one
that always adds 100. Each closure remembers its own captured \code{n}, long after
\code{make\_adder} has finished running.

\section{Closures that remember and change state}

Because a closure carries its captured variables with it, it can also \emph{update}
them between calls giving a function that has memory. A counter is the
canonical example:

\begin{salam}
func make_counter() func() i32:
    mut count i32 = 0
    ret () => :
        count = count + 1
        ret count
    end
end

func main:
    next func() i32 = make_counter()
    println next()
    println next()
    println next()
end
\end{salam}

\begin{progout}
1
2
3
\end{progout}

Each call to \code{next()} increments and returns the captured \code{count}. The
variable \code{count} lives on inside the closure even though \code{make\_counter}
returned long ago the closure has kept it alive. Call \code{make\_counter()} a
second time and you get a fresh, independent counter starting again at zero.

\begin{deep}[In Depth: how closures are implemented]
A closure is more than a pointer to code it must also carry the variables it
captured. Salam implements this by allocating a small \emph{environment} on the
heap holding the captured values, and pairing it with the lambda's code. The
function value you pass around is really this code-plus-environment bundle. When
you call it, the captured variables are read from (and written to) that
environment, which is why \code{count} survives between calls and why two counters
do not interfere. The compiler performs this ``lifting'' of the lambda and its
captures automatically; you simply write the closure and it works.
\end{deep}

\section{Why closures matter}

Closures let you parameterise \emph{behaviour}, not just values. A sort routine
can take a closure describing how to compare two elements; a retry helper can take
a closure describing the work to attempt; an event system can take closures to run
when something happens. In each case you hand over a small piece of custom logic,
created exactly where it is needed and carrying whatever surrounding context it
requires. This keeps general machinery (the sorter, the retrier) cleanly separated
from the specific behaviour you plug into it.

\begin{tip}
Lambdas are at their best when they are small a comparison, a transformation, a
quick test. If a lambda grows into a substantial block of logic, give it a real
name with \code{func} instead; a named function is easier to read, reuse, and
discuss. Reach for a lambda when naming it would add noise, not clarity.
\end{tip}

\section{Chapter summary}

\begin{itemize}[leftmargin=1.4em]
  \item A \textbf{lambda} is an unnamed function value:
        \code{(params) => expression}, or \code{(params) => : ... end} for a
        block body.
  \item A function value has a type like \code{func(i32) i32} or \code{func()
        i32}.
  \item A \textbf{higher-order function} takes or returns a function value; call a
        function-typed parameter just like any function.
  \item A \textbf{closure} is a lambda that captures variables from its
        surrounding scope; the captured values live on with the closure.
  \item Closures can carry mutable state (e.g. a counter), and each closure
        instance has its own independent captured variables.
  \item Use lambdas for small, local behaviour; promote anything substantial to a
        named function.
\end{itemize}

\section{Exercises}

\exercise Create a lambda \code{square} that returns its argument times itself,
store it with \code{auto}, and print \code{square(7)}.

\exercise Write a higher-order function \code{apply\_twice(f func(i32) i32, x
i32) i32} that returns \code{f(f(x))}, and test it with a lambda that adds 3.

\exercise Write \code{make\_multiplier(factor i32) func(i32) i32} that returns a
closure multiplying by \code{factor}. Build a ``times 3'' and a ``times 10''
function from it.

\exercise Using \code{make\_counter}, create two independent counters and show by
interleaving their calls that they do not share state.

\exercise Write a higher-order function \code{transform\_all} that takes a
\code{Vector<i32>} and a \code{func(i32) i32}, and returns a new vector with the
function applied to every element. Test it with a doubling lambda.

\chapter{Error Handling and Deferred Execution}
\label{ch:error-defer}

Real programs deal with things that can go wrong: a withdrawal larger than the
balance, a number that will not parse, a file that is not there. Equally, real
programs acquire resources open files, allocated memory that must be
released no matter how the work turns out. This chapter covers Salam's
straightforward approach to both: reporting failure through return values, and
guaranteeing cleanup with \code{defer}.

\section{Salam's philosophy: errors are values}

Salam has no exceptions, no \code{try}/\code{catch}, no invisible control flow
that whisks you off to some distant handler. Instead, a function that can fail
\emph{says so in what it returns}, and the caller checks. This keeps the flow of
your program on the page: every place that can fail is visible, and every failure
is handled where it happens. There are a few common patterns, from simplest to
most expressive.

\section{Pattern 1: a boolean success flag}

The simplest signal is a \code{bool}: \code{true} for success, \code{false} for
failure. You saw this in Chapter~9 with the bank account, whose \code{withdraw}
returned whether it succeeded:

\begin{salam}
struct Account:
    balance : int = 0
    func withdraw(amount : int) : bool:
        if amount > this.balance:
            ret false              // failure: not enough funds
        end
        this.balance = this.balance - amount
        ret true                   // success
    end
end

func main:
    mut acc : Account = Account { balance = 100 }
    if acc.withdraw(30):
        println "withdrew 30; balance is", acc.balance
    else:
        println "withdrawal failed"
    end
    if acc.withdraw(1000):
        println "withdrew 1000"
    else:
        println "withdrawal failed: insufficient funds"
    end
end
\end{salam}

\begin{progout}
withdrew 30; balance is 70
withdrawal failed: insufficient funds
\end{progout}

The caller cannot ignore the outcome by accident: it is right there in the
\code{if}. This pattern is perfect when the only thing you need to know is
``did it work?''

\section{Pattern 2: a sentinel return value}

When a function returns a value but might have nothing meaningful to give back, a
\textbf{sentinel} a special value that means ``no result'' is a common
choice. A search that returns the index of an item can return \code{-1} to mean
``not found,'' since a real index is never negative:

\begin{salam}
func index_of(values : int[5], target : int) : int:
    for mut i = 0, i < 5, i = i + 1:
        if values[i] == target:
            ret i
        end
    end
    ret -1                         // sentinel: not found
end

func main:
    nums : int[5] = [10, 20, 30, 40, 50]
    found : auto = index_of(nums, 30)
    if found >= 0:
        println "found at index", found
    else:
        println "not found"
    end
    println "search for 99:", index_of(nums, 99)
end
\end{salam}

\begin{progout}
found at index 2
search for 99: -1
\end{progout}

\begin{warn}
A sentinel only works when there is a value the real results can never take. For
an index, \code{-1} is safe. But ``zero means failure'' is dangerous if zero is a
legitimate result, and ``empty string means error'' fails if the empty string is
valid data. When no value can be safely reserved, prefer the \code{Option}
approach below.
\end{warn}

\section{Pattern 3: \texttt{Option} a value that may be absent}

The standard library's \code{Option} type makes ``maybe a value, maybe nothing''
explicit and safe, without stealing a value from the result's range. An
\code{Option<T>} either holds a \code{T} (it is \emph{some} value) or holds
nothing (it is \emph{none}). Import it with \code{import "option"}.

\begin{salam}
import "option"

func safe_divide(a : int, b : int) : option.Option<int>:
    if b == 0:
        ret Option { value = 0, present = false }    // none
    end
    ret option.Some(a / b)                            // some result
end

func main:
    r1 : auto = safe_divide(10, 2)
    r2 : auto = safe_divide(10, 0)
    println "10/2 is some?", option.IsSome(r1)
    println "10/0 is some?", option.IsSome(r2)
    println "10/2 =", option.UnwrapOr(r1, -1)
    println "10/0 =", option.UnwrapOr(r2, -1)    // falls back to -1
end
\end{salam}

\begin{progout}
10/2 is some? true
10/0 is some? false
10/2 = 0
10/0 = -1
\end{progout}

The \code{option} module gives you the tools to work with these values safely:

\begin{center}
\begin{tabular}{@{}ll@{}}
\toprule
\textbf{Function} & \textbf{What it does} \\
\midrule
\code{option.Some(x)}        & wrap \code{x} as a present value \\
\code{option.IsSome(o)}      & \code{true} if \code{o} holds a value \\
\code{option.IsNone(o)}      & \code{true} if \code{o} holds nothing \\
\code{option.UnwrapOr(o, d)} & the held value, or \code{d} if none \\
\bottomrule
\end{tabular}
\end{center}

\code{UnwrapOr} is especially handy: it gets the value out, supplying a fallback
for the ``none'' case in one step, so you never accidentally use a value that is
not there. To build a ``none,'' you construct the \code{Option} with
\code{present = false}, as the example shows.

\begin{note}
\code{Option} is the type to reach for whenever ``no answer'' is a legitimate
outcome: a lookup that may miss, a parse that may fail, the first element of a
possibly-empty collection. It forces the caller to acknowledge the empty case,
turning a whole class of ``forgot to check'' bugs into something the type system
helps you remember.
\end{note}

\section{Deferred execution with \texttt{defer}}

Now to the second half of the chapter: making sure cleanup happens. The
\code{defer} keyword schedules a statement to run \emph{when the current function
exits} no matter how it exits. You write the cleanup right next to the thing it
cleans up, and Salam guarantees it runs later:

\begin{salam}
func main:
    defer println "cleanup runs last"
    println "doing work"
    println "more work"
end
\end{salam}

\begin{progout}
doing work
more work
cleanup runs last
\end{progout}

The \code{defer}'d line was \emph{written} first but \emph{ran} last, at the
moment \code{main} finished. The point is that the cleanup is registered the
instant you set up whatever needs cleaning, so you can never forget it later.

\subsection{Defers run in reverse order}

If you register several defers, they run in \textbf{last-in, first-out} order ---
the most recently deferred runs first. This mirrors how resources nest: the last
thing you opened is the first thing you should close.

\begin{salam}
func main:
    defer println "cleanup 1"
    defer println "cleanup 2"
    println "start"
    println "end"
end
\end{salam}

\begin{progout}
start
end
cleanup 2
cleanup 1
\end{progout}

\subsection{Defer runs even on an early return}

The real power of \code{defer} is that it runs however the function ends ---
including an early \code{ret} from somewhere in the middle. This is what makes it
reliable: the cleanup cannot be skipped by a code path you forgot about.

\begin{salam}
func greet(name : str):
    defer println "done greeting"
    println "hello"
    if name == "":
        println "no name given"
        ret                        // early return -- defer still runs
    end
    println name
end

func main:
    greet("Ali")
    println "---"
    greet("")
end
\end{salam}

\begin{progout}
hello
Ali
done greeting
---
hello
no name given
done greeting
\end{progout}

Both calls print \code{"done greeting"} the normal path through \code{Ali} and
the early-return path through the empty name. You write the cleanup once, and it
fires on every exit.

\section{Putting it together: acquire, defer, use}

The idiomatic pattern is: acquire a resource, immediately \code{defer} its
release, then use it freely without worrying about the exits. Whether the function
returns normally, returns early on an error, or falls through, the resource is
released exactly once. We will see this with real files and memory in Chapters~17
and~18; the shape is always the same:

\begin{salam}
func process():
    // acquire something here
    defer println "release the resource"

    println "use the resource"
    if true:
        println "early exit on some condition"
        ret               // the defer above still runs
    end
    println "this line is skipped"
end

func main:
    process()
end
\end{salam}

\begin{progout}
use the resource
early exit on some condition
release the resource
\end{progout}

\begin{tip}
Pair every acquisition with a \code{defer} on the very next line. Reading top to
bottom, ``open file; defer close file'' is impossible to get wrong the close is
guaranteed and sits right where you can see it. This single habit eliminates the
most common kind of resource leak.
\end{tip}

\section{Chapter summary}

\begin{itemize}[leftmargin=1.4em]
  \item Salam reports errors through \emph{return values}, not exceptions every
        failure is visible and handled at the call site.
  \item Return a \code{bool} when the caller only needs ``did it work?''.
  \item Return a \textbf{sentinel} (like \code{-1}) when a value's range has room
        for a ``no result'' marker but only then.
  \item Use \code{option.Option<T>} (\code{Some}, \code{IsSome}, \code{IsNone},
        \code{UnwrapOr}) when ``no value'' must be represented safely.
  \item \code{defer} schedules a statement to run when the function exits, in
        last-in-first-out order, \emph{including} on early returns.
  \item The reliable pattern is: acquire, immediately \code{defer} the release,
        then use.
\end{itemize}

\section{Exercises}

\exercise Write \code{safe\_sqrt(x f64) option.Option<f64>} that returns
\code{none} for negative input and \code{Some} of the square root otherwise (use
the \code{math} module from the next chapter, or \code{x ** 0.5}).

\exercise Write a function \code{find\_first\_negative(a int[5]) int} that returns
the index of the first negative element or \code{-1} if there is none.

\exercise Add two \code{defer} statements and a normal statement to a function,
and predict the output order before running it.

\exercise Write a function with an early \code{ret} guarded by a condition and a
\code{defer} that prints ``finished''. Verify ``finished'' prints on both paths.

\exercise Rewrite the \code{safe\_divide} example so the caller prints a friendly
message in the ``none'' case instead of using \code{UnwrapOr}.

\chapter{Modules and Packages}
\label{ch:modules-packages}

A program that outgrows a single file needs structure. \textbf{Modules} and
\textbf{packages} are how Salam splits a large program into pieces, controls what
each piece exposes to the others, and pulls in the standard library. You have been
quietly using this system every time you wrote \code{import "str"}; now we look at
it properly.

\section{Importing the standard library}

The standard library is organised into modules, each a self-contained bundle of
related functions: \code{math} for numbers, \code{str} for strings, \code{io} for
input and output, and many more. To use one, \code{import} it by name and call its
functions with the module name as a prefix:

\begin{salam}
import "math"

func main:
    println math.Sqrt(144.0)
    println math.MaxI(7, 3)
end
\end{salam}

\begin{progout}
12
7
\end{progout}

The prefix \code{math.} is called \textbf{qualification}, and it does two
useful things. It tells the reader exactly where \code{Sqrt} comes from, and it
keeps names from colliding: \code{math.Max} and your own \code{max} can coexist
without confusion.

\begin{note}
A handful of operations need no import at all: \code{print}, \code{println}, and
\code{len} are built into the language and are always available. Everything else
in the standard library lives in a module you import.
\end{note}

\section{Importing several modules}

You can write one \code{import} per module, or list several module paths after a
single \code{import} no commas, no parentheses, just the paths side by side:

\begin{salam}
import "math" "str" "io"

func main:
    io.Println(str.Upper("salam"))
    println math.Gcd(12, 18)
end
\end{salam}

\begin{progout}
SALAM
6
\end{progout}

\section{Aliasing an import}

Sometimes a module's name is long, or clashes with a name you want to use
yourself. You can give an import a shorter \textbf{alias} by writing the alias
before the module name:

\begin{salam}
import strutil "str" "io"

func main:
    io.Println(strutil.Upper("salam"))
end
\end{salam}

\begin{progout}
SALAM
\end{progout}

Now the \code{str} module is reached as \code{strutil}. Aliases are a matter of
convenience and clarity; use them when they make the code read better, and leave
them off otherwise.

\section{How imports are resolved}

When you write \code{import "math"}, Salam looks for the module \code{math} in the
standard library specifically the file \code{std/math/math.salam}. A nested
module path works the same way: \code{import "db/sqlite"} resolves to
\code{std/db/sqlite/sqlite.salam}. An import whose name ends in \code{.salam},
such as \code{import "helpers.salam"}, is instead treated as a \emph{file relative
to the importing file} which is how you bring in your \emph{own} modules, the
subject of the next section.

\section{Writing your own modules}

A module is just a \code{.salam} file with definitions in it. To split your
program across files, put related functions and types in their own file and import
it where you need them. Two pieces of machinery make this work: the \code{package}
declaration and the \code{pub} keyword.

\subsection{Declaring a package}

Every file belongs to a \textbf{package}, named at the top with the \code{package}
keyword. Your program's entry point lives in package \code{main}:

\begin{salam}
package main

func main:
    println "Hello from package main"
end
\end{salam}

\begin{progout}
Hello from package main
\end{progout}

A file with no \code{package} line belongs to \code{main} by default which is
why the small single-file programs earlier in this book never needed one. As soon
as you split code across files, naming the package makes the structure explicit.

\subsection{Exposing names with \texttt{pub}}

By default, a top-level function, struct, or type is \textbf{private} to its
package: other files cannot see it. To make something visible to code that imports
your module, mark it \code{pub} (short for \emph{public}):

\begin{salam}
// file: geometry.salam
package geometry

pub func area_of_circle(r : f64) : f64:
    ret 3.14159 * r * r
end

func helper() : f64:        // no 'pub' -- private to this file
    ret 0.0
end
\end{salam}

Another file can then import \code{geometry.salam} and use its public functions,
while \code{helper} stays hidden:

\begin{salam}
// file: main.salam
package main
import "geometry.salam"

func main:
    println geometry.area_of_circle(2.0)
end
\end{salam}

This is \textbf{encapsulation}: a module presents a small, deliberate public
surface the things it promises to support while keeping its internal
helpers private and free to change. The same \code{pub} keyword marks the public
parts of the standard library itself; open any \code{std} module and you will see
\code{pub func} and \code{pub struct} on exactly the names you are meant to use.

\begin{tip}
Make things \code{pub} sparingly. The smaller a module's public surface, the
easier it is to use correctly and the freer you are to rewrite its insides without
breaking anyone. Start everything private; promote a name to \code{pub} only when
another file genuinely needs it.
\end{tip}

\section{A tour of the standard library's layout}

It helps to know roughly what is available before you need it. The standard
library is a collection of modules under \code{std}, each in its own folder. A
selection:

\begin{center}
\begin{tabular}{@{}ll@{}}
\toprule
\textbf{Module} & \textbf{Purpose} \\
\midrule
\code{io}, \code{console} & printing and reading input, files \\
\code{fmt}    & formatting values into strings \\
\code{math}   & numeric functions and constants \\
\code{str}    & string operations \\
\code{rand}   & random numbers \\
\code{time}   & dates, timestamps, sleeping \\
\code{collections} & generic containers (\code{Vector}, \code{HashMap}, \dots) \\
\code{os}, \code{fs} & the operating system and file system \\
\code{json}, \code{csv} & data formats \\
\code{regex}  & regular expressions \\
\code{option} & the \code{Option} type \\
\code{mem}    & manual memory management \\
\code{sync}, \code{thread} & concurrency primitives \\
\code{testing} & assertions for unit tests \\
\code{net}, \code{tcp}, \code{http} & networking \\
\bottomrule
\end{tabular}
\end{center}

The next two chapters take a proper tour: Chapter~16 covers the everyday modules
(\code{io}, \code{fmt}, \code{math}, \code{str}, \code{rand}, \code{time}), and
Chapter~17 covers data structures and the world outside your program
(\code{collections}, \code{os}, \code{fs}, \code{json}, and friends).

\begin{deep}[In Depth: qualified names and the generated C]
Qualification is not only for readability it also guarantees that names from
different modules never collide in the compiled program. When \code{math.Sqrt} is
compiled, it becomes a uniquely named C function (something like
\code{\_Salam\_math\_Sqrt\_f64}) that cannot clash with a \code{Sqrt} from any
other module or with your own code. The package system, in other words, gives you
human-friendly short names at the source level and collision-proof unique names
underneath the best of both worlds, with no run-time cost.
\end{deep}

\section{Chapter summary}

\begin{itemize}[leftmargin=1.4em]
  \item \code{import "name"} brings in a standard-library module; call its
        functions qualified, as \code{name.func(...)}.
  \item Import several modules at once by listing their paths after one
        \code{import}; give any import an \code{alias} by writing the alias
        before its path.
  \item \code{import "math"} resolves to \code{std/math/math.salam}; an import
        ending in \code{.salam} is a file relative to the importer (your own
        modules).
  \item Every file has a \code{package}; the entry point is in \code{main} (the
        default when omitted).
  \item Top-level names are private to their package unless marked \code{pub}.
  \item \code{print}, \code{println}, and \code{len} are built in and need no
        import.
\end{itemize}

\section{Exercises}

\exercise Write a program that imports \code{math} and \code{str} on a single
\code{import} line and prints the square root of 2 alongside the word
\code{"root"} in uppercase.

\exercise Import \code{str} under the alias \code{s} and use it to reverse a word.

\exercise Sketch (on paper) a two-file program: a \code{shapes.salam} module with
a \code{pub} function and a private helper, and a \code{main.salam} that imports
it. Mark which names are visible where.

\exercise Open one of the real \code{std} modules in the Salam source and list
five \code{pub} functions it offers. Try calling two of them.

\exercise Explain in a comment why qualification (\code{math.Max}) lets you keep
your own function called \code{max} without conflict.

\chapter{The Standard Library, Part 1}
\label{ch:stdlib-1}

A language is only as productive as the library that comes with it. Salam ships a
substantial standard library so that you rarely have to reinvent the basics. This
chapter tours the everyday modules input and output, formatting, mathematics,
randomness, and time with runnable examples of the functions you will use most.
The next chapter covers data structures and the world outside your program.

\section{Input and output: \texttt{io}}

You already know \code{println} and \code{print}, which are built in. The
\code{io} module offers the same operations as named functions (handy when you
import it for file work) plus the all-important \code{Input}, which reads a line
of text the user types:

\begin{salam}
import "io"

func main:
    io.Print("What is your name? ")
    name str = io.Input()
    io.Println("Hello, " + name + "!")
end
\end{salam}

\code{io.Input()} reads one line from the keyboard and returns it as a \code{str}
(without the trailing newline). Combined with the \code{str} conversions from
Chapter~8, this is how you build interactive programs:

\begin{salam}
import (
    "io"
    "str"
)

func main:
    io.Print("Enter a number: ")
    n auto = str.ToInt(io.Input())
    io.Println("Its square is " + str.FromInt((n * n) as i64))
end
\end{salam}

\begin{note}
Because \code{io.Input} returns text, you must convert it to a number with
\code{str.ToInt} (or \code{str.ToFloat}) before doing arithmetic. Reading input as
text and converting explicitly is safer than guessing types you decide exactly
how each line should be interpreted.
\end{note}

\section{Formatting values: \texttt{fmt}}

The \code{fmt} module turns values into strings and lays them out neatly. The
conversions are the bread and butter:

\begin{salam}
import "fmt"

func main:
    println fmt.Int(255 as i64)        // "255"
    println fmt.Float(3.14159)         // "3.14159"
    println fmt.Bool(true)             // "true"
end
\end{salam}

\begin{progout}
255
3.14159
true
\end{progout}

For aligning text into columns tables, menus, reports \code{fmt} provides
padding and repetition:

\begin{salam}
import "fmt"

func main:
    println fmt.PadLeft("hi", 5)       // "   hi"
    println fmt.PadRight("hi", 5) + "|" // "hi   |"
    println fmt.Repeat("=", 10)        // "=========="
end
\end{salam}

\begin{progout}
   hi
hi   |
==========
\end{progout}

\code{PadLeft} and \code{PadRight} widen a string to a given width by adding
spaces on the left or right, so numbers and labels line up in columns.
\code{Repeat} builds a string of repeated copies perfect for rules and
separators.

\section{Mathematics: \texttt{math}}

The \code{math} module covers the numeric functions you would expect, working on
\code{f64} values (with a few integer helpers). Here are the most useful, grouped
by purpose.

\subsection{Rounding and sign}

\begin{salam}
import "math"

func main:
    println math.Floor(3.7)    // 3  (round down)
    println math.Ceil(3.2)     // 4  (round up)
    println math.Round(3.5)    // 4  (nearest)
    println math.Abs(-4.0)     // 4  (absolute value)
end
\end{salam}

\begin{progout}
3
4
4
4
\end{progout}

\subsection{Roots, powers, and logarithms}

\begin{salam}
import "math"

func main:
    println math.Sqrt(144.0)       // 12
    println math.Pow(2.0, 10.0)    // 1024
    println math.Min(2.0, 9.0), math.Max(2.0, 9.0)
end
\end{salam}

\begin{progout}
12
1024
2 9
\end{progout}

\subsection{A reference of common \texttt{math} functions}

\begin{center}
\begin{tabularx}{\linewidth}{@{}lX@{}}
\toprule
\textbf{Function} & \textbf{Returns} \\
\midrule
\code{math.Sqrt(x)} & square root of \code{x} \\
\code{math.Pow(b, e)} & \code{b} raised to the power \code{e} \\
\code{math.Abs(x)} & absolute value (float) \\
\code{math.Floor(x)} / \code{math.Ceil(x)} & round down / up \\
\code{math.Round(x)} & round to nearest \\
\code{math.Min(a, b)} / \code{math.Max(a, b)} & smaller / larger (float) \\
\code{math.MinI(a, b)} / \code{math.MaxI(a, b)} & smaller / larger (integer) \\
\code{math.Gcd(a, b)} / \code{math.Lcm(a, b)} & greatest common divisor / least common multiple \\
\code{math.Sin(x)}, \code{math.Cos(x)}, \code{math.Tan(x)} & trigonometric functions \\
\code{math.Log(x)}, \code{math.Log10(x)}, \code{math.Exp(x)} & logarithms and the exponential \\
\bottomrule
\end{tabularx}
\end{center}

The module also defines the constants \code{math.PI}, \code{math.E}, and
\code{math.TAU}:

\begin{salam}
import "math"

func main:
    r f64 = 2.0
    area auto = math.PI * r * r
    println "area of a circle of radius 2:", area
end
\end{salam}

\begin{progout}
area of a circle of radius 2: 12.5664
\end{progout}

\subsection{Integer helpers}

A handful of functions work directly on integers, so you do not have to cast back
and forth. \code{math.MaxI} and \code{math.MinI} compare integers, and
\code{math.Gcd} and \code{math.Lcm} are exact integer operations:

\begin{salam}
import "math"

func main:
    println math.MaxI(7, 3)    // 7
    println math.Gcd(12, 18)   // 6
    println math.Lcm(4, 6)     // 12
end
\end{salam}

\begin{progout}
7
6
12
\end{progout}

\section{Randomness: \texttt{rand}}

The \code{rand} module produces random values for games, simulations, sampling,
and test data. Start by seeding the generator (use \code{rand.SeedAuto()} for an
unpredictable seed from the clock, or a fixed seed for repeatable runs), then draw
values:

\begin{salam}
import "rand"

func main:
    rand.Seed(42 as i64)                       // fixed seed: repeatable
    dice auto = rand.IntRange(1 as i64, 6 as i64)
    println "you rolled:", dice
    println "coin flip heads?", rand.Bool()
    println "random fraction:", rand.Float()   // in [0, 1)
end
\end{salam}

The most useful functions:

\begin{center}
\begin{tabular}{@{}ll@{}}
\toprule
\textbf{Function} & \textbf{Returns} \\
\midrule
\code{rand.Seed(s)} / \code{rand.SeedAuto()} & set a fixed / clock-based seed \\
\code{rand.IntRange(lo, hi)} & a random integer in \code{[lo, hi]} \\
\code{rand.Float()} & a random \code{f64} in \code{[0, 1)} \\
\code{rand.Bool()} & a random \code{true}/\code{false} \\
\code{rand.UUID()} & a random unique identifier string \\
\bottomrule
\end{tabular}
\end{center}

\begin{tip}
Seed with a \emph{fixed} number while you are developing and testing, so your
program behaves the same on every run and bugs are reproducible. Switch to
\code{rand.SeedAuto()} for the real, shipped program, where you want genuine
unpredictability. Forgetting to seed at all gives you the same ``random'' sequence
every time occasionally useful, usually a surprise.
\end{tip}

\section{Time: \texttt{time}}

The \code{time} module reads the clock, formats timestamps, and pauses your
program. A timestamp is a count of seconds, which you can format into a readable
date or pick apart into fields:

\begin{salam}
import "time"

func main:
    now auto = time.Now()                 // seconds since the epoch
    println "formatted:", time.Format(now)
    println "year:", time.Year(now)
    println "seconds in a day:", time.SecondsPerDay
end
\end{salam}

\code{time.Now()} returns the current time as an integer count of seconds;
\code{time.Format} renders it as \code{"YYYY-MM-DD HH:MM:SS"}, and helpers like
\code{time.Year}, \code{time.Month}, and \code{time.Day} extract individual
fields. The module also defines the constants \code{time.SecondsPerMinute},
\code{time.SecondsPerHour}, and \code{time.SecondsPerDay}.

To pause, use \code{time.Sleep}, which takes a number of milliseconds:

\begin{salam}
import "time"

func main:
    println "wait for it..."
    time.Sleep(500)        // pause half a second
    println "done"
end
\end{salam}

\section{Strings: a reminder}

The \code{str} module \code{Upper}, \code{Lower}, \code{Split}, \code{Join},
\code{Find}, \code{Replace}, \code{FromInt}, \code{ToInt}, and the rest is the
other workhorse of everyday programming, and we covered it thoroughly in
Chapter~8. Keep it in mind alongside the modules here; together, \code{io},
\code{fmt}, \code{math}, \code{str}, \code{rand}, and \code{time} handle the vast
majority of small-program needs.

\section{A worked example: a guessing game}

Here is a complete little program that ties several of these modules together a
number-guessing game (shown with a fixed secret so the output is predictable):

\begin{salam}
import (
    "io"
    "str"
    "rand"
)

func main:
    rand.Seed(7 as i64)
    secret auto = rand.IntRange(1 as i64, 100 as i64)
    io.Print("Guess my number (1-100): ")
    guess auto = str.ToInt(io.Input())
    if guess == secret:
        io.Println("Exactly right!")
    else guess < secret:
        io.Println("Too low.")
    else:
        io.Println("Too high.")
    end
end
\end{salam}

It seeds the generator, picks a secret in $[1, 100]$, reads a guess as text,
converts it to a number, and compares drawing on \code{rand}, \code{io}, and
\code{str} at once. That is the standard library doing exactly what it is for:
letting you assemble a real program from ready-made parts.

\section{Chapter summary}

\begin{itemize}[leftmargin=1.4em]
  \item \code{io} provides \code{Print}, \code{Println}, and \code{Input} (read a
        line as a \code{str}).
  \item \code{fmt} converts values (\code{Int}, \code{Float}, \code{Bool}) and
        aligns text (\code{PadLeft}, \code{PadRight}, \code{Repeat}).
  \item \code{math} offers rounding, roots, powers, trig, logs, integer helpers
        (\code{MaxI}, \code{Gcd}, \code{Lcm}), and the constants \code{PI},
        \code{E}, \code{TAU}.
  \item \code{rand} generates random integers, floats, booleans, and ids; seed it
        for repeatability or with \code{SeedAuto} for unpredictability.
  \item \code{time} reads and formats the clock, extracts date fields, exposes
        second-count constants, and sleeps with \code{Sleep}.
\end{itemize}

\section{Exercises}

\exercise Write a program that asks for the user's name and age and prints a
sentence using both (convert the age with \code{str.ToInt}).

\exercise Print a small multiplication table aligned in columns using
\code{fmt.PadLeft}.

\exercise Use \code{math} to print the hypotenuse of a right triangle with legs 3
and 4 (\code{Sqrt} of the sum of squares).

\exercise Roll two dice with \code{rand.IntRange} and print their sum; seed with a
fixed value so you can predict the result.

\exercise Print the current date with \code{time.Format(time.Now())}, then sleep
one second and print it again.

\exercise Combine \code{rand} and \code{fmt} to print five random integers between
1 and 99, each padded to width 2, on one line.

% ===================== Chapter 17 =====================
\chapter{The Standard Library, Part 2}
\label{ch:stdlib-2}

The first standard-library chapter covered computation: numbers, text, time. This
one covers \emph{structure and the outside world} the containers that organise
your data, and the modules that read files, parse data formats, and match
patterns. With these you can write programs that do real work on real data.

\section{Collections beyond the vector}

You met \code{Vector<T>} in Chapter~7. The \code{collections} module which is
available without an explicit import for the core containers offers several more
generic data structures, each suited to a particular access pattern.

\subsection{HashMap: key-value lookup}

A \code{HashMap<K, V>} associates \textbf{keys} of type \code{K} with
\textbf{values} of type \code{V}, and looks a value up by its key almost
instantly. It is the right tool whenever you think ``given \emph{this}, find
\emph{that}'' a username to a score, a word to its count, a product code to a
price.

\begin{salam}
func main:
    mut prices HashMap<str, i32> = HashMap {}
    prices.insert("apple", 3)
    prices.insert("banana", 5)
    println "size:", prices.size()
    println "apple costs:", prices.get("apple")
    println "have banana?", prices.has("banana")
    println "have cherry?", prices.has("cherry")
    prices.remove("apple")
    println "size now:", prices.size()
    prices.free()
end
\end{salam}

\begin{progout}
size: 2
apple costs: 3
have banana? true
have cherry? false
size now: 1
\end{progout}

The key methods are \code{insert(k, v)} to add or update an entry, \code{get(k)}
to retrieve a value, \code{has(k)} to test for a key, \code{remove(k)} to delete
one, \code{size()} for the count, and \code{free()} when you are done.

\begin{warn}
Call \code{get} only for a key you know is present check first with \code{has}
if you are not sure. As with vectors, a \code{HashMap} manages its own memory, so
remember \code{free()} when you have finished with it.
\end{warn}

\subsection{Stack and Queue}

When you care about the \emph{order} in which items come back out, two classic
containers help. A \code{Stack<T>} is last-in-first-out the most recently pushed
item pops first, like a stack of plates. A \code{Queue<T>} is first-in-first-out ---
items leave in the order they arrived, like a line at a counter.

\begin{salam}
func main:
    mut s Stack<i32> = Stack {}
    s.push(1)
    s.push(2)
    println "top:", s.peek()        // 2
    println "pop:", s.pop()         // 2 (last in, first out)
    s.free()

    mut q Queue<str> = Queue {}
    q.enqueue("first")
    q.enqueue("second")
    println q.dequeue()             // "first" (first in, first out)
    println q.dequeue()             // "second"
    q.free()
end
\end{salam}

\begin{progout}
top: 2
pop: 2
first
second
\end{progout}

A \code{Stack} uses \code{push}/\code{pop}/\code{peek}; a \code{Queue} uses
\code{enqueue}/\code{dequeue}/\code{peek}. Both report \code{size()} and
\code{is\_empty()}, and both should be \code{free}d. Reach for a stack when you
need to undo or backtrack (the most recent thing first), and a queue when you must
process things fairly in arrival order.

\section{Reading and writing files}

The \code{io} module, which you met for console output, also reads and writes
files. The simplest operations handle a whole file at once:

\begin{salam}
import "io"

func main:
    io.WriteFile("notes.txt", "line one\nline two\n")
    content auto = io.ReadFile("notes.txt")
    print content
end
\end{salam}

\begin{progout}
line one
line two
\end{progout}

\code{io.WriteFile(path, text)} writes (creating or replacing the file), and
\code{io.ReadFile(path)} returns the whole file as a \code{str}. For
line-by-line processing, \code{io.Lines(path)} returns a \code{Vector<str>}, one
entry per line, which you can walk with the techniques from Chapter~7.

\section{The operating system: \texttt{os}}

The \code{os} module is your program's window onto the system it runs on: files
and directories, environment variables, command-line arguments, and the exit
code.

\begin{salam}
import "os"

func main:
    println "exists?", os.Exists("notes.txt")
    println "current dir:", os.Cwd()
    home auto = os.Env("HOME")
    println "home:", home
end
\end{salam}

Common \code{os} operations include \code{os.Exists(path)} to test for a file,
\code{os.Remove(path)} to delete one, \code{os.Mkdir(path)} to create a directory,
\code{os.ListDir(path)} to list a directory's contents as a \code{Vector<str>},
\code{os.Args()} for the command-line arguments, \code{os.Env(name)} for an
environment variable, and \code{os.Exit(code)} to end the program with a status
code.

\begin{salam}
import (
    "os"
    "io"
)

func main:
    // back up a file, then clean up after ourselves
    io.WriteFile("data.txt", "important\n")
    defer os.Remove("data.txt")
    io.WriteFile("data.bak", io.ReadFile("data.txt"))
    defer os.Remove("data.bak")
    println "backup made; files exist?", os.Exists("data.bak")
end
\end{salam}

Notice the \code{defer}s from Chapter~14 doing exactly what they were made for:
each temporary file is scheduled for removal the instant it is created, so the
program tidies up no matter how it ends.

\section{Working with JSON}

JSON is the lingua franca of data exchange, and the \code{json} module reads it
without ceremony. You can validate a document, pull out fields by key, and ask
whether a key is present:

\begin{salam}
import "json"

func main:
    doc str = "{ \"user\": \"ada\", \"age\": 36 }"
    println "valid?", json.Valid(doc)
    println "user:", json.Get(doc, "user")
    println "age:", json.GetInt(doc, "age")
    println "has email?", json.Has(doc, "email")
end
\end{salam}

\begin{progout}
valid? true
user: ada
age: 36
has email? false
\end{progout}

\code{json.Valid} checks that the text is well-formed JSON; \code{json.Get} returns
a field as a string, while \code{json.GetInt} and \code{json.GetFloat} return it as
a number; \code{json.Has} tests for a key. There is also \code{json.Escape} to make
a string safe to embed in JSON, and \code{json.Minify} to strip whitespace. This is
enough to consume the JSON that web services and config files speak.

\section{Pattern matching with regular expressions}

A \textbf{regular expression} (regex) is a compact pattern that describes a set of
strings ``one or more digits,'' ``a word followed by a space.'' The
\code{regex} module matches text against such patterns. For one-off use, the
\code{...Str} functions take the pattern as a plain string:

\begin{salam}
import "regex"

func main:
    println "has digits?", regex.MatchStr("[0-9]+", "abc123")
    println "digits start at:", regex.FindStr("[0-9]+", "abc123def")
    println regex.ReplaceStr("[0-9]+", "a1b22c333", "#")
end
\end{salam}

\begin{progout}
has digits? true
digits start at: 3
a#b22c333
\end{progout}

\code{regex.MatchStr(pattern, text)} reports whether the pattern occurs anywhere in
the text; \code{regex.FindStr} returns the index where the first match begins (or a
negative value if there is none); and \code{regex.ReplaceStr} substitutes a match
with replacement text.

\begin{deep}[In Depth: compile once, match many]
If you will apply the same pattern many times say, validating every line of a
large file compiling it once is more efficient than re-parsing the pattern on
each call. \code{regex.Compile(pattern)} returns a reusable \code{Regex} value, and
\code{regex.Match(re, text)}, \code{regex.Find(re, text)}, and
\code{regex.Replace(re, text, with)} operate on it. The \code{...Str} convenience
functions are exactly these with an implicit compile each time; prefer the
compiled form in a hot loop.
\end{deep}

\section{Other data-format modules}

Rounding out the picture, the standard library also includes a \code{csv} module
for reading and writing comma-separated values (the format spreadsheets export),
an \code{encoding} module for common encodings, and the networking modules
\code{net}, \code{tcp}, and \code{http} for talking to other machines. The shape
is always the same as what you have seen here: import the module, call its
qualified functions, and free anything that manages its own memory. Once you are
comfortable with the modules in these two chapters, the rest follow the same
patterns.

\section{A worked example: word frequencies}

Let us combine a \code{HashMap}, the \code{str} module, and a vector into a small
but genuinely useful program: counting how often each word appears in a sentence.

\begin{salam}
import "str"

func main:
    text str = "the cat sat on the mat the cat ran"
    mut counts HashMap<str, i32> = HashMap {}
    mut words Vector<str> = str.Split(text, " ")
    mut i i32 = 0
    while i < words.len():
        w auto = words.get(i)[0]
        if counts.has(w):
            counts.insert(w, counts.get(w) + 1)
        else:
            counts.insert(w, 1)
        end
        i += 1
    end
    println "distinct words:", counts.size()
    println "the appears", counts.get("the"), "times"
    println "cat appears", counts.get("cat"), "times"
    words.free()
    counts.free()
end
\end{salam}

\begin{progout}
distinct words: 6
the appears 3 times
cat appears 2 times
\end{progout}

Split the sentence into words, then for each word either start its count at 1 or
add one to its existing count, using \code{has} to tell the two cases apart. The
\code{HashMap} does the heavy lifting instant lookup by word and the result
is a tally of the text. This ``count occurrences'' pattern appears everywhere,
from analytics to spell-checkers.

\section{Chapter summary}

\begin{itemize}[leftmargin=1.4em]
  \item \code{HashMap<K, V>} maps keys to values: \code{insert}, \code{get},
        \code{has}, \code{remove}, \code{size}, \code{free}.
  \item \code{Stack<T>} (push/pop/peek) is last-in, first-out;
        \code{Queue<T>} (enqueue/dequeue) is first-in, first-out; both
        provide \code{free}.
  \item \code{io.WriteFile}, \code{io.ReadFile}, and \code{io.Lines} handle whole
        files; \code{os} provides \code{Exists}, \code{Remove}, \code{Mkdir},
        \code{ListDir}, \code{Args}, \code{Env}, \code{Exit}.
  \item \code{json} reads JSON: \code{Valid}, \code{Get}, \code{GetInt},
        \code{Has}, \code{Escape}, \code{Minify}.
  \item \code{regex} matches patterns: \code{MatchStr}, \code{FindStr},
        \code{ReplaceStr}, or the compiled \code{Compile}/\code{Match}/\code{Find}.
  \item \code{csv}, \code{encoding}, and the networking modules round out the
        library, all following the same import-and-qualify pattern.
\end{itemize}

\section{Exercises}

\exercise Build a \code{HashMap<str, i32>} of three people's ages, then look up and
print one age and test for a name that is not present.

\exercise Use a \code{Stack} to reverse the order of the numbers 1 to 5: push them
all, then pop them into a line of output.

\exercise Write a program that writes three lines to a file with
\code{io.WriteFile}, reads them back with \code{io.Lines}, and prints how many
lines there were.

\exercise Given the JSON \code{"\{ \textbackslash"title\textbackslash":
\textbackslash"Salam\textbackslash", \textbackslash"pages\textbackslash": 400
\}"}, print the title and the page count.

\exercise Use \code{regex.MatchStr} to test whether a string looks like it contains
a year (four digits in a row).

\exercise Extend the word-frequency example to also print the word ``the'' count
only if it appears more than twice.

\chapter{Memory Management}
\label{ch:memory}

Every value your program works with lives somewhere in the computer's memory.
Most of the time Salam handles that for you and you need not think about it. But
the growable containers vectors, hash maps and a few advanced techniques do
involve memory you are responsible for. This chapter explains Salam's memory model
plainly, so you know which memory takes care of itself and which you must tend.

\section{Two kinds of memory}

It helps to picture memory in two parts. The \textbf{stack} holds your local
variables the \code{int}s, the \code{bool}s, the fixed-size arrays and structs.
This memory is managed completely automatically: when a variable comes into scope
it is created, and when its scope ends it is reclaimed. You have been relying on
this since Chapter~2 without a second thought, and you can continue to.

The \textbf{heap} holds things whose size is not known until the program runs, or
that must outlive the function that created them the growable buffer behind a
vector, the entries of a hash map. Heap memory must be explicitly released when no
longer needed, which is why those containers have a \code{free} method.

\section{Values that manage themselves}

The good news first: the great majority of what you write needs no memory
management at all.

\begin{itemize}[leftmargin=1.4em]
  \item \textbf{Scalars} \code{int}, \code{f64}, \code{bool}, \code{char}, an
        enum live on the stack and vanish automatically.
  \item \textbf{Fixed-size arrays and structs} live inline, on the stack, and are
        likewise reclaimed automatically when they go out of scope. Recall from
        Chapter~9 that assigning or passing a struct \emph{copies} it there is
        no shared hidden allocation to worry about.
  \item \textbf{Strings} are special: although a \code{str} points at heap bytes,
        Salam manages those bytes for you (see the box below). You never
        \code{free} a string.
\end{itemize}

\begin{deep}[In Depth: the string arena]
Joining strings with \code{+}, or calling functions like \code{str.Upper}, must
allocate room for the new text on the heap yet you never free a string. How?
Salam keeps a per-program \emph{arena} for string data: all the bytes that string
operations produce are allocated from it, and the whole arena is released at once
when the program ends. The trade-off is deliberate: you get the convenience of
never managing string memory, at the cost of strings living until the program
exits. For the short-lived, string-heavy programs that are most common, this is
exactly the right bargain.
\end{deep}

\section{Containers you must free}

The exception to ``it manages itself'' is the growable containers. A
\code{Vector<T>} or a \code{HashMap<K, V>} owns a heap buffer that it resizes as
it grows, and that buffer must be returned to the system when you are finished. So
every growable container has a \code{free} method, and you should call it:

\begin{salam}
func main:
    mut v : Vector<int> = Vector {}
    v.push(1)
    v.push(2)
    v.push(3)
    println "sum of three pushes:", v.get(0)[0] + v.get(1)[0] + v.get(2)[0]
    v.free()              // release the buffer
end
\end{salam}

\begin{progout}
sum of three pushes: 6
\end{progout}

This applies to \code{Vector}, \code{HashMap}, \code{Stack}, \code{Queue}, and the
other growable collections. The fixed-size array \code{int[5]}, by contrast, needs
no freeing it lives on the stack.

\subsection{Pairing creation with \texttt{free} using \texttt{defer}}

The cleanest way never to forget a \code{free} is to schedule it the moment you
create the container, with the \code{defer} you learned in Chapter~14. Written
together on adjacent lines, the release is guaranteed and impossible to overlook:

\begin{salam}
func sum_first_n(n : int) : int:
    mut v : Vector<int> = Vector {}
    defer v.free()                 // guaranteed cleanup, however we return
    for mut i = 1, i <= n, i = i + 1:
        v.push(i)
    end
    mut total : int = 0
    mut j : int = 0
    while j < v.len():
        total += v.get(j)[0]
        j += 1
    end
    ret total
end

func main:
    println sum_first_n(5)
end
\end{salam}

\begin{progout}
15
\end{progout}

Because \code{defer v.free()} runs whenever \code{sum\_first\_n} returns through
the normal path or any early \code{ret} the vector's memory is always released.
This ``create, defer free, use'' pattern is the single most important habit for
leak-free Salam code.

\begin{warn}
Two mistakes to avoid. \textbf{Forgetting to free} a container leaks its
memory harmless in a tiny script that exits immediately, but a real problem in a
long-running program that creates containers in a loop. \textbf{Using a container
after freeing it} is worse: the buffer is gone, and reading it is undefined
behaviour. Free a container exactly once, as the last thing you do with it ---
which is precisely what \code{defer} at the point of creation guarantees.
\end{warn}

\section{No hidden ownership rules}

If you have heard of languages with elaborate ``ownership'' and ``borrowing''
systems, you can relax: Salam has none of that ceremony. There are no lifetime
annotations to write, no borrow checker to satisfy, no move semantics to learn.
The model is the simple one above stack values manage themselves, growable
containers get a \code{free} and the discipline is the single habit of pairing
each container with a \code{defer free}.

This simplicity is a deliberate design choice. It keeps the language approachable
and the code readable, at the cost of asking you to remember \code{free} for the
handful of containers that need it. In practice that is a small, learnable
responsibility, and \code{defer} reduces it to a reflex.

\section{Manual memory with \texttt{mem}}

For most programs the containers are all the heap memory you will ever touch. But
when you need raw control implementing your own data structure, or working at
the boundary with C (Chapter~19) the \code{mem} module gives you direct
allocation, the same primitives the standard library's containers are built on:

\begin{salam}
import "mem"

func main:
    p : void* = mem.Allocate(64 as u64)     // ask for 64 bytes
    if p == null:
        println "allocation failed"
        ret
    end
    println "got 64 bytes"
    mem.Free(p)                            // hand them back
    println "released"
end
\end{salam}

\begin{progout}
got 64 bytes
released
\end{progout}

\code{mem.Allocate(n)} requests \code{n} bytes and returns a \code{void*} (a raw,
untyped pointer), or \code{null} if the request fails which you should check.
\code{mem.Free(p)} returns the memory. The module also offers
\code{mem.Reallocate} to resize a block, \code{mem.Set} to fill bytes, and
\code{mem.Copy} to copy them.

\begin{warn}
Manual memory is powerful and unforgiving: a missing \code{Free} leaks, a double
\code{Free} corrupts, and using freed memory crashes. You will rarely need
\code{mem} directly the built-in containers cover almost every case. Reach for
it only when you are building something the standard library does not provide, and
when you do, guard every allocation with a \code{defer mem.Free} just as you would
a container.
\end{warn}

\begin{deep}[In Depth: tracking down leaks]
Because Salam compiles to C, the usual C tooling works on Salam programs. The
compiler can build with AddressSanitizer (the \code{--asan} flag, or the
\code{salam memcheck} command) to catch use-after-free and out-of-bounds errors
at run time, and the \code{mem} module exposes counters allocations made, frees
done, bytes still live that let a program audit its own memory use. If you ever
suspect a leak, these turn ``something is wrong'' into a precise location.
\end{deep}

\section{Chapter summary}

\begin{itemize}[leftmargin=1.4em]
  \item Scalars, fixed arrays, and structs live on the stack and are reclaimed
        automatically; you never manage them.
  \item Strings are heap-backed but managed by a per-program arena never
        \code{free} a string.
  \item Growable containers (\code{Vector}, \code{HashMap}, \code{Stack},
        \code{Queue}) own heap buffers and must be \code{free}d.
  \item Pair every container creation with \code{defer x.free()} to guarantee
        cleanup on every exit path.
  \item Salam has no ownership/borrowing system the model is deliberately
        simple.
  \item \code{mem} (\code{Allocate}, \code{Free}, \code{Reallocate}, \code{Set},
        \code{Copy}) provides raw manual memory for advanced use; guard it
        carefully.
\end{itemize}

\section{Exercises}

\exercise Write a function that builds a \code{Vector<str>} of three greetings,
prints them, and frees the vector using \code{defer}.

\exercise Take a function that creates a \code{HashMap}, returns early on some
condition, and uses \code{defer} to ensure the map is freed on both paths.

\exercise Explain in a comment why a fixed array \code{int[10]} needs no
\code{free} but a \code{Vector<int>} does.

\exercise Use \code{mem.Allocate} and \code{mem.Free} to grab and release 128
bytes, checking for a \code{null} result, and guard the free with \code{defer}.

\exercise Describe, in your own words, what the string arena does and one
situation where its ``strings live until the program ends'' behaviour could matter.

% ===================== Chapter 19 =====================
\chapter{The Foreign Function Interface}
\label{ch:ffi}

Salam compiles to C, which gives it a superpower: it can call C functions
directly, with no glue code and no performance penalty. The decades of C
libraries the world runs on math, compression, databases, graphics, the
operating system itself are all reachable from Salam through the
\textbf{foreign function interface}, or FFI. This chapter shows how to declare a C
function and call it as if it were your own.

\section{Declaring a C function}

To call a C function you first tell Salam about it: its name, its parameters, and
its return type. You do this in an \code{extern "C"} block, which lists the
\emph{signatures} of C functions without bodies the bodies live in C, after
all:

\begin{salam}
extern "C":
    func printf(format str, ...) i32
    func strlen(s str) u64
end

func main:
    printf("Hello from C's printf!\n")
    msg str = "external"
    printf("the length of \"%s\" is %d\n", msg, strlen(msg))
end
\end{salam}

\begin{progout}
Hello from C's printf!
the length of "external" is 8
\end{progout}

Inside the \code{extern "C"} block, each \code{func} line is a \emph{declaration}:
it says ``a C function with this name and shape exists; let me call it.'' There is
no body. When you call \code{printf} or \code{strlen} from Salam, the compiler
emits a direct call to the real C function exactly as a C program would. The
trailing \code{...} on \code{printf} marks it \textbf{variadic} (it accepts a
variable number of arguments), which is how \code{printf} takes its format
arguments.

\section{How Salam types map to C}

For an FFI call to work, Salam's types must line up with C's. The mapping is
direct and is why the integration is seamless:

\begin{center}
\begin{tabular}{@{}lll@{}}
\toprule
\textbf{Salam} & \textbf{C} & \textbf{Notes} \\
\midrule
\code{i32} & \code{int32\_t} & a 32-bit signed integer \\
\code{i64} & \code{int64\_t} & a 64-bit signed integer \\
\code{u64} & \code{uint64\_t} & a 64-bit unsigned integer \\
\code{f64} & \code{double} & a 64-bit float \\
\code{f32} & \code{float} & a 32-bit float \\
\code{bool} & one byte & \\
\code{char} & one byte & \\
\code{str} & \code{const char*} & a NUL-terminated C string \\
\code{void*} & \code{void*} & an untyped pointer \\
(no return type) & \code{void} & a procedure \\
\bottomrule
\end{tabular}
\end{center}

When you declare an \code{extern} function, choose the Salam types that match the C
function's real signature according to this table. \code{strlen} in C takes a
\code{const char*} and returns a \code{size\_t}, so we declared it as
\code{func strlen(s str) u64}.

\section{Calling the C math library}

A practical example: C's math library (\code{libm}) has fast, accurate
implementations of functions like \code{sqrt}, \code{pow}, and \code{hypot}.
Declare the ones you want and call them:

\begin{salam}
extern "C":
    func printf(format str, ...) i32
    func sqrt(x f64) f64
    func pow(base f64, exp f64) f64
    func hypot(x f64, y f64) f64
end

func main:
    printf("sqrt(2)    = %.6f\n", sqrt(2.0))
    printf("2^10       = %.0f\n", pow(2.0, 10.0))
    printf("hypot(3,4) = %.1f\n", hypot(3.0, 4.0))
end
\end{salam}

\begin{progout}
sqrt(2)    = 1.414214
2^10       = 1024
hypot(3,4) = 5.0
\end{progout}

Each declared function is called like any Salam function, and the result comes
straight back from the C implementation. (Of course, for these particular
functions you would normally just use Salam's own \code{math} module which is
itself partly built on \code{libm} but the same technique works for \emph{any}
C library, including ones with no Salam wrapper.)

\section{Linking external libraries}

The functions above \code{printf}, \code{sqrt} come from the C standard
library, which is always linked in, so declaring them is enough. To use a function
from a \emph{separate} library, you must also tell the compiler to link that
library, with a \code{link} directive before the \code{extern} block:

\begin{salam}
link "sqlite3"

extern "C":
    func sqlite3_libversion() str
    func printf(format str, ...) i32
end

func main:
    printf("SQLite version: %s\n", sqlite3_libversion())
end
\end{salam}

\code{link "sqlite3"} instructs the compiler to link against the SQLite library
(the equivalent of passing \code{-lsqlite3} to a C compiler), making its functions
available to your \code{extern} declarations. You can link several libraries with
several \code{link} lines.

\begin{deep}[In Depth: platform-conditional linking]
Library names sometimes differ across operating systems. Salam's preprocessor lets
you choose the right one at build time with \code{@if}:
\begin{salam}
@if SALAM_OS_WINDOWS
link "winsqlite3"
@else
link "sqlite3"
@end
\end{salam}
The symbols \code{SALAM\_OS\_WINDOWS}, \code{SALAM\_OS\_LINUX}, and
\code{SALAM\_OS\_MACOS} let a single program link the correct library on each
platform the same technique C programmers use with \code{\#ifdef}, expressed in
Salam.
\end{deep}

\section{Working with pointers across the boundary}

Many C functions deal in pointers \code{malloc} returns one, \code{free} takes
one. Salam's \code{void*} type matches C's, so you can call them directly:

\begin{salam}
extern "C":
    func malloc(size u64) void*
    func free(ptr void*)
    func puts(s str) i32
end

func main:
    mut p void* = malloc(64 as u64)
    if p == null:
        puts("malloc failed")
    else:
        puts("malloc(64) gave a non-null pointer")
        free(p)
    end
end
\end{salam}

\begin{progout}
malloc(64) gave a non-null pointer
\end{progout}

Notice that the safe habits from Chapter~18 carry straight over the boundary:
check the pointer against \code{null}, and free what you allocate. The FFI does not
add safety of its own you are calling C, with C's rules so the discipline is
yours to keep.

\section{When to use the FFI}

The FFI is a bridge to be crossed deliberately, not a place to live. Reach for it
when:

\begin{itemize}[leftmargin=1.4em]
  \item you need a capability that exists as a mature C library a database
        driver, a compression codec, an image format and rewriting it in Salam
        would be wasteful;
  \item you must call an operating-system function the standard library does not
        wrap;
  \item you are integrating Salam into an existing C codebase.
\end{itemize}

For everything else, prefer Salam's own standard library and types: they are safer
(the compiler checks them), more convenient, and bilingual. The FFI is the escape
hatch for when you genuinely need the wider C world powerful exactly because it
is there when you need it and ignorable when you do not.

\begin{warn}
Code reached through the FFI is outside Salam's safety net. The compiler trusts
your \code{extern} declarations: if you declare a function with the wrong types,
or pass a bad pointer, the result is a crash or corruption, not a friendly
compile-time error. Declare each C function to match its real signature exactly,
and treat values that cross the boundary with the same care a C programmer would.
\end{warn}

\section{Chapter summary}

\begin{itemize}[leftmargin=1.4em]
  \item An \code{extern "C":} block declares C functions (name, parameters,
        return type, no body) so Salam can call them directly.
  \item Salam types map one-to-one onto C types (\code{i32}$\to$\code{int32\_t},
        \code{str}$\to$\code{const char*}, \code{void*}$\to$\code{void*}, \dots);
        declare each function to match.
  \item A trailing \code{...} marks a variadic function such as \code{printf}.
  \item Standard-library C functions need only declaration; others need a
        \code{link "lib"} directive (optionally guarded by \code{@if} for
        per-platform names).
  \item The FFI adds no safety check pointers, free what you allocate, and
        declare signatures exactly.
  \item Use the FFI to reach mature C libraries and OS functions; prefer Salam's
        own library otherwise.
\end{itemize}

\section{Exercises}

\exercise Declare C's \code{puts} (\code{func puts(s str) i32}) and use it to print
two lines.

\exercise Declare \code{sqrt} and \code{pow} from the math library and print the
square root of 2 and 2 to the 16th power.

\exercise Declare \code{getenv(name str) str} and print the value of an environment
variable of your choice.

\exercise Declare \code{malloc} and \code{free}, allocate 256 bytes, check for
\code{null}, and free them guarding the free with \code{defer}.

\exercise Explain in a comment why an incorrect \code{extern} declaration cannot be
caught by the Salam compiler, and what discipline protects you instead.

\chapter{Concurrency}
\label{ch:concurrency}

So far our programs have done one thing at a time, in order. \textbf{Concurrency}
is doing several things at once splitting work across the multiple cores every
modern computer has, so a big job finishes sooner. Salam expresses concurrency
with \emph{threads}, and protects shared data with \emph{synchronization}. This
chapter introduces both, along with the care they demand.

\section{Threads: running code in parallel}

A \textbf{thread} is an independent strand of execution. Your program already has
one the one running \code{main}. With \code{spawn} you can start another, which
runs a function of your choosing \emph{alongside} the rest of the program:

\begin{salam}
func greet():
    println "hello from a thread"
end

func main:
    handle : i64 = spawn(greet)
    println "hello from main"
    join(handle)
    println "the thread has finished"
end
\end{salam}

\code{spawn(greet)} starts \code{greet} running on a new thread and immediately
returns a \textbf{handle} an \code{i64} that identifies the thread while
\code{main} carries on. Because the two run concurrently, the order of the first
two lines of output is not guaranteed. \code{join(handle)} then \emph{waits} for
the thread to finish before continuing, so the final line always prints last.

\begin{note}
\code{spawn} takes a function with no parameters. To give a thread something to
work on, have it read shared variables (carefully see below) or use a closure
that captures the data it needs. \code{join} is how you know a thread is done; a
program should \code{join} every thread it spawns before relying on that thread's
results.
\end{note}

\section{The danger of shared data}

Threads become useful and dangerous the moment they touch the \emph{same}
data. Suppose two threads each add to a shared counter a thousand times. You would
expect the total to climb by two thousand. But incrementing a counter is really
three steps under the hood read the value, add one, write it back and if two
threads interleave those steps, they can read the same old value and each write
back the same new one, losing an increment. This is a \textbf{race condition}: the
result depends on the unpredictable timing of the threads, and it is one of the
hardest kinds of bug to find, because it appears only sometimes.

\section{Mutexes: one at a time}

The fix is to ensure that only one thread touches the shared data at a time. A
\textbf{mutex} (from ``mutual exclusion'') is a lock that exactly one thread can
hold. A thread \emph{locks} it before touching the shared data and \emph{unlocks}
it afterwards; any other thread that tries to lock it must wait until it is free.
The \code{sync} module provides them:

\begin{salam}
import "sync"

mut counter : int = 0
mut lock : sync.Mutex = sync.NewMutex()

func work():
    repeat 1000:
        sync.Lock(lock)
        counter = counter + 1     // protected: only one thread here at a time
        sync.Unlock(lock)
    end
end

func main:
    t1 : i64 = spawn(work)
    t2 : i64 = spawn(work)
    join(t1)
    join(t2)
    println "counter:", counter
    sync.Destroy(lock)
end
\end{salam}

\begin{progout}
counter: 2000
\end{progout}

Now the answer is reliably \code{2000}. Between \code{sync.Lock(lock)} and
\code{sync.Unlock(lock)} the \emph{critical section} only one thread ever
runs, so the read-add-write happens without interference. Create a mutex with
\code{sync.NewMutex()}, and call \code{sync.Destroy} on it when all the threads are
done.

\begin{warn}
A mutex protects data only if \emph{every} access to that data goes through it.
One thread that touches the shared counter without locking re-introduces the race.
And always pair a \code{Lock} with an \code{Unlock} on every path a lock that is
never released leaves other threads waiting forever (a \emph{deadlock}). Keep
critical sections short and make the lock/unlock pairing obvious.
\end{warn}

\section{Waiting for many threads}

When you spawn several threads and need to wait for all of them, you can collect
their handles and \code{join} each. For larger fan-outs, the \code{sync} module
also offers a \textbf{wait group} a counter you raise as you spawn work and
that threads decrement as they finish, letting the main thread block until the
count reaches zero:

\begin{salam}
import "sync"

mut wg : sync.WaitGroup = sync.NewWaitGroup()

func task():
    println "working"
    sync.Done(wg)         // signal completion
end

func main:
    sync.Add(wg, 3 as i64)
    spawn(task)
    spawn(task)
    spawn(task)
    sync.Wait(wg)         // block until all three call Done
    println "all tasks complete"
    sync.DestroyWaitGroup(wg)
end
\end{salam}

\code{sync.Add(wg, n)} announces \code{n} tasks; each task calls \code{sync.Done}
when finished; \code{sync.Wait} blocks until every \code{Done} has arrived. A wait
group is convenient when the number of workers is known but you do not want to
track each handle by hand.

\section{When concurrency pays and when it costs}

Concurrency is not free. Threads have a setup cost, locks add overhead, and ---
most of all concurrent code is harder to reason about and harder to debug than
sequential code. It pays off when there is genuine parallel work to do: a big
computation you can split into independent chunks, or many slow operations (such
as network requests) that can wait at the same time instead of one after another.

\begin{tip}
Reach for concurrency only when a measured need justifies it, and keep shared,
mutable state to an absolute minimum the less data threads share, the fewer
locks you need and the fewer races can occur. The simplest concurrent designs give
each thread its own data to work on and combine the results at the end, touching
shared state only briefly and under a lock. Start sequential; add threads when the
workload truly calls for them.
\end{tip}

\begin{deep}[In Depth: how threads map to the system]
Salam's \code{spawn} creates a real operating-system thread (a \code{pthread} on
Unix-like systems, a Windows thread on Windows), so threads run on separate cores
and your program can genuinely use a multi-core machine. The \code{sync.Mutex} is
a thin wrapper over the platform's native mutex. Because these are true OS
threads, the costs and the rules are the same as in C which is why the
discipline around shared data matters, and why the standard library gives you the
classic tools (mutexes, wait groups) to manage it.
\end{deep}

\section{Chapter summary}

\begin{itemize}[leftmargin=1.4em]
  \item \code{spawn(f)} runs the parameterless function \code{f} on a new thread
        and returns an \code{i64} handle; \code{join(handle)} waits for it to
        finish.
  \item Sharing mutable data between threads risks a \emph{race condition}, where
        the result depends on timing.
  \item A \code{sync.Mutex} (\code{NewMutex}, \code{Lock}, \code{Unlock},
        \code{Destroy}) ensures only one thread is in a critical section at a time.
  \item Every access to shared data must go through the lock, and every
        \code{Lock} must be matched by an \code{Unlock}.
  \item A \code{sync.WaitGroup} (\code{Add}, \code{Done}, \code{Wait}) waits for a
        known number of threads to complete.
  \item Use concurrency only when real parallel work justifies it, and minimise
        shared mutable state.
\end{itemize}

\section{Exercises}

\exercise Spawn a thread that prints the numbers 1 to 5 while \code{main} prints
the letters A to E. Join the thread, and observe that the interleaving varies
between runs.

\exercise Modify the shared-counter example to use \emph{four} threads of 1000
increments each and confirm the total is 4000.

\exercise Deliberately remove the \code{Lock}/\code{Unlock} from the counter
example, run it several times, and describe what you observe about the final total.

\exercise Use a \code{WaitGroup} to run five tasks that each print their own
number, and wait for all of them before printing ``done''.

\exercise Explain in a comment why keeping the critical section (the code between
\code{Lock} and \code{Unlock}) short is good practice.

\chapter{Testing and Debugging}
\label{ch:testing-debugging}

Code that is never tested is code you only \emph{hope} works. Tests turn hope into
evidence: small programs that check your code does what you claim, and keep
checking as you change it. And when something is wrong, debugging is how you find
out why. This chapter covers both writing tests with the \code{testing} module,
and the tools and habits that make finding bugs faster.

\section{Why test?}

A test is a tiny experiment: ``if I call this function with these inputs, I should
get this result.'' Writing tests has three payoffs. They catch mistakes
immediately, while the code is fresh in your mind. They guard against
\emph{regressions} a change that fixes one thing while quietly breaking
another because you can re-run every test after every change. And they document
how your code is meant to be used, in the most honest form possible: working
examples.

\section{Assertions with the \texttt{testing} module}

The \code{testing} module provides \textbf{assertions} statements that check a
condition and record a pass or a failure. You write a normal program that calls a
function and asserts the expected result:

\begin{salam}
import "testing"

func add(a : int, b : int) : int:
    ret a + b
end

func main:
    testing.AssertEqInt(add(2, 3), 5)
    testing.AssertEqInt(add(-1, 1), 0)
    testing.AssertTrue(add(2, 2) == 4)
    testing.Summary()
end
\end{salam}

\begin{progout}
tests: 3 passed, 0 failed
\end{progout}

Each \code{Assert} checks its condition and tallies the result; \code{testing.\-
Summary()} prints the totals at the end. When everything passes, you get a clean
report like the one above quiet confidence that your code does what you said.

\subsection{The assertion functions}

\begin{center}
\begin{tabular}{@{}ll@{}}
\toprule
\textbf{Assertion} & \textbf{Passes when} \\
\midrule
\code{testing.AssertTrue(c)}        & \code{c} is \code{true} \\
\code{testing.AssertFalse(c)}       & \code{c} is \code{false} \\
\code{testing.AssertEqInt(got, want)} & two integers are equal \\
\code{testing.AssertEqStr(got, want)} & two strings are equal \\
\code{testing.AssertEqBool(got, want)} & two booleans are equal \\
\code{testing.Assert(c)}            & \code{c} is \code{true} (general) \\
\bottomrule
\end{tabular}
\end{center}

The \code{...Eq...} forms are worth preferring over a plain \code{AssertTrue(a ==
b)} because, when they fail, they tell you \emph{both} values what you got and
what you wanted which is exactly what you need to diagnose the problem.

\subsection{Reading a failure}

A failing assertion does not stop the program; it records the failure, prints a
message, and lets the remaining checks run, so one run tells you about \emph{all}
the problems. Watch what happens when an expected value is wrong:

\begin{salam}
import "testing"

func add(a : int, b : int) : int:
    ret a + b
end

func main:
    testing.AssertEqInt(add(2, 3), 5)
    testing.AssertEqInt(add(2, 2), 5)    // wrong on purpose
    testing.Summary()
end
\end{salam}

\begin{progout}
FAIL: ints differ: got 4 want 5
tests: 1 passed, 1 failed
\end{progout}

The message names exactly what went wrong \code{got 4 want 5} and the
summary counts one pass and one failure. In a real test file you would fix the
code (or the expectation) until the summary reads all passes.

\section{Organising tests}

Group related checks into named functions, then call them from \code{main}. This
keeps each function focused on one piece of behaviour and makes the output easier
to follow:

\begin{salam}
import "testing"

func is_even(n : int) : bool:
    ret n % 2 == 0
end

func test_is_even():
    testing.AssertTrue(is_even(4))
    testing.AssertFalse(is_even(7))
    testing.AssertTrue(is_even(0))
end

func main:
    test_is_even()
    testing.Summary()
end
\end{salam}

\begin{progout}
tests: 3 passed, 0 failed
\end{progout}

As your program grows, keep test functions next to the code they exercise, give
them clear \code{test\_} names, and call them all from a single \code{main} that
ends in \code{testing.Summary()}. The exit code reflects the number of failures,
so a test program also slots neatly into an automated build that should stop when
a test fails.

\begin{tip}
Write a test the moment you find a bug, \emph{before} you fix it: a test that
reproduces the bug. Then fix the code until the test passes. Now you have not only
solved the problem but guaranteed it can never silently come back the test will
catch it if it does. This habit, fixing bugs ``test first,'' steadily builds a net
that makes your whole program safer to change.
\end{tip}

\section{Debugging: finding out why}

When a program misbehaves, the goal is to replace guessing with \emph{seeing}. The
most powerful debugging tool is also the humblest.

\subsection{Print debugging}

A well-placed \code{println} that shows a variable's value at a key moment will
solve the large majority of bugs. When a result is wrong, print the inputs and
intermediate values along the path that produces it, and compare what you see with
what you expected. The discrepancy is almost always where the bug lives:

\begin{salam}
func average(a : int[3]) : f64:
    mut sum : int = 0
    each x in a:
        sum += x
        println "after adding", x, "sum is", sum    // a temporary probe
    end
    ret (sum as f64) / 3.0
end

func main:
    data : int[3] = [10, 20, 30]
    println "average:", average(data)
end
\end{salam}

\begin{progout}
after adding 10 sum is 10
after adding 20 sum is 30
after adding 30 sum is 60
average: 20
\end{progout}

The probe lines confirm the running total is building correctly. Once you have
found and fixed the problem, remove the temporary prints they were scaffolding,
not part of the finished program.

\subsection{Let the compiler help}

Many bugs never reach run time at all: Salam's static checks catch them first. A
type mismatch, an immutable variable you tried to reassign, an unused parameter, an
undefined name each is reported at compile time with a clear message and the
line it occurred on. Read those messages carefully; they are the cheapest bug
reports you will ever get, arriving before the program even runs. Much of what
would be a debugging session in a dynamic language is, in Salam, simply a compiler
error you fix on the spot.

\subsection{Heavier tools}

For the stubborn cases, the compiler offers more. Building with debug information
(\code{salam debug}) lets you step through your program in a debugger, inspecting
variables as it runs. And because Salam compiles to C, the memory checker
(\code{salam memcheck}, which builds with AddressSanitizer) pinpoints
use-after-free and out-of-bounds errors invaluable when you have been working
with manual memory or the FFI. You will not need these often, but they are there
when a print statement is not enough.

\begin{deep}[In Depth: a debugging method]
When a bug resists, work like a scientist rather than thrashing. \textbf{Reproduce}
it reliably find the smallest input that triggers it. \textbf{Localise} it by
bisection: confirm the data is right at the start and wrong at the end, then check
the middle, narrowing the region until you find the line where right becomes wrong.
\textbf{Form a hypothesis} about the cause, and \textbf{test} it with a targeted
print or a change. The smallest reproducing case plus bisection will corner almost
any bug, far faster than staring at the code hoping for inspiration.
\end{deep}

\section{Chapter summary}

\begin{itemize}[leftmargin=1.4em]
  \item Tests turn ``I hope it works'' into evidence and guard against
        regressions.
  \item The \code{testing} module asserts conditions (\code{AssertTrue},
        \code{AssertEqInt}, \code{AssertEqStr}, \dots) and reports totals with
        \code{Summary}; failures print both values and do not halt the run.
  \item Group checks into \code{test\_}-named functions called from \code{main}.
  \item Write a failing test that reproduces a bug before fixing it.
  \item \code{println} probes solve most bugs; remove them once done.
  \item Salam's compile-time checks catch many bugs before run time read the
        messages.
  \item \code{salam debug} and \code{salam memcheck} handle the hard cases.
\end{itemize}

\section{Exercises}

\exercise Write a function \code{factorial(n int) int} and a \code{test\_factorial}
that asserts \code{factorial(0) == 1}, \code{factorial(5) == 120}, and one more
case of your choosing.

\exercise Add a deliberately wrong assertion to a passing test file and study the
failure message and the summary line.

\exercise Take a buggy function (for example, a loop with an off-by-one error) and
use \code{println} probes to locate the exact line where the values go wrong.

\exercise Write tests for a \code{clamp(x, lo, hi)} function covering a value below
the range, inside it, and above it.

\exercise Describe, in a comment, the difference between a bug Salam catches at
compile time and one you would only find by running tests, with an example of
each.

% ===================== Chapter 22 =====================
\chapter{Bilingual Programming in Salam}
\label{ch:bilingual}

We have mentioned throughout this book that Salam can be written in Persian as
well as English. Now we give that feature the chapter it deserves. Salam's
bilingualism is not a translation layer bolted on top --- it is built into the
language: the very same compiler reads English and Persian keywords, and produces
identical programs from either. This chapter shows how, and why it matters.

\section{One language, two vocabularies}

The key idea is simple. A programming language is made of \emph{keywords} (the
fixed words like \code{func}, \code{if}, \code{while}) and \emph{your} names (the
variables and functions you invent). Salam keeps the grammar --- the structure,
the rules, the meaning --- exactly the same in both languages, and changes only the
spelling of the keywords. \code{func} becomes \code{\fa{تابع}}; \code{if} becomes
\code{\fa{اگر}}; everything else about the program is unchanged.

Here is the same program twice. First in English:

\begin{salam}
func add(a i32, b i32) i32:
    ret a + b
end

func main:
    mut total i32 = 0
    for mut i = 1; i <= 5; i = i + 1:
        total = total + i
    end
    println "sum 1..5 =", total
end
\end{salam}

And now, exactly the same program, in Persian:

\begin{salamfa}
تابع جمع(الف صحیح۳۲, ب صحیح۳۲) صحیح۳۲:
    بازگشت الف + ب
پایان

تابع اصلی:
    متغیر مجموع صحیح۳۲ = 0
    برای متغیر i = 1; i <= 5; i = i + 1:
        مجموع = مجموع + i
    پایان
    چاپ "مجموع ۱ تا ۵ =", مجموع
پایان
\end{salamfa}

\begin{progout}
مجموع ۱ تا ۵ = 15
\end{progout}

Look closely and you will see the structures are identical, line for line. Only
the keywords and the type names changed. A Persian-speaking student reads the
second version as naturally as you read the first --- and they are learning exactly
the same concepts, because they \emph{are} the same concepts.

\section{The keyword equivalents}

Every keyword you have learned has a Persian twin. Here are the most common; the
complete list is in Appendix~A.

\begin{center}
\begin{tabular}{@{}llll@{}}
\toprule
\textbf{English} & \textbf{Persian} & \textbf{English} & \textbf{Persian} \\
\midrule
\code{func}   & \code{\fa{تابع}}    & \code{ret}     & \code{\fa{بازگشت}} \\
\code{if}     & \code{\fa{اگر}}     & \code{else}    & \code{\fa{وگرنه}} \\
\code{while}  & \code{\fa{تاوقتی}}  & \code{for}     & \code{\fa{برای}} \\
\code{repeat} & \code{\fa{تکرار}}   & \code{to}      & \code{\fa{تا}} \\
\code{mut}    & \code{\fa{متغیر}}   & \code{const}   & \code{\fa{ثابت}} \\
\code{struct} & \code{\fa{ساختار}}  & \code{enum}    & \code{\fa{شمارش}} \\
\code{interface} & \code{\fa{واسط}} & \code{end}     & \code{\fa{پایان}} \\
\code{import} & \code{\fa{واردکردن}} & \code{pub}    & \code{\fa{عمومی}} \\
\code{true}   & \code{\fa{درست}}    & \code{false}   & \code{\fa{نادرست}} \\
\code{break}  & \code{\fa{بشکن}}    & \code{continue} & \code{\fa{ادامه}} \\
\code{println} & \code{\fa{چاپ}}    & \code{print}   & \code{\fa{بنویس}} \\
\code{main}   & \code{\fa{اصلی}}    & \code{this}    & \code{\fa{این}} \\
\bottomrule
\end{tabular}
\end{center}

The type names have Persian forms too, which you first saw in Chapter~2:

\begin{center}
\begin{tabular}{@{}llll@{}}
\toprule
\textbf{English} & \textbf{Persian} & \textbf{English} & \textbf{Persian} \\
\midrule
\code{i32} & \code{\fa{صحیح۳۲}} & \code{i64} & \code{\fa{صحیح۶۴}} \\
\code{f64} & \code{\fa{اعشاری۶۴}} & \code{bool} & \code{\fa{منطقی}} \\
\code{str} & \code{\fa{رشته}}   & \code{char} & \code{\fa{کاراکتر}} \\
\code{auto} & \code{\fa{خودکار}} & & \\
\bottomrule
\end{tabular}
\end{center}

\begin{note}
The entry point itself is part of the vocabulary: it is \code{main} in English and
\code{\fa{اصلی}} in Persian. When you run a Persian program, the compiler looks for
\code{\fa{اصلی}} as the starting point, just as it looks for \code{main} in an
English one.
\end{note}

\section{Compiling a Persian program}

A Salam source file is written in \emph{one} language, and you tell the compiler
which with the \code{--lang} flag. English is the default; for a Persian file, pass
\code{--lang=fa}:

\begin{shell}
salam run hello.salam --lang=fa
salam build hello.salam --lang=fa --output=hello
\end{shell}

Everything else about the workflow --- \code{run}, \code{build}, \code{exec},
\code{fmt} --- is identical. The flag simply tells the compiler which keyword
vocabulary to expect.

\begin{warn}
A single file uses a single language: you cannot mix English and Persian
\emph{keywords} in the same file. A file is either an English program or a Persian
one, compiled with the matching \code{--lang}. This is deliberate --- it keeps each
file coherent and unambiguous. What you \emph{can} freely mix is covered next.
\end{warn}

\section{What you can mix}

The keyword rule is about keywords only. The \emph{names you choose} --- variables,
functions, struct fields --- may be in whatever language you like, in either kind of
file, because Salam identifiers may contain Unicode letters. So a Persian program
can use English identifiers (as the loop variable \code{i} did above), and an
English program can use Persian identifiers:

\begin{salam}
func main:
    naam str = "Sara"        // a Latin-letter Persian word as a name
    println "Hello,", naam
end
\end{salam}

\begin{progout}
Hello, Sara
\end{progout}

Likewise, \emph{strings} and \emph{comments} are free text and may be in any
language regardless of the file's keyword language --- which is why the Persian
program above could print the Persian message \code{"\fa{مجموع ۱ تا ۵ =}"} and an
English program can hold Persian comments. Keywords pick a side; your content does
not have to.

\begin{note}
Notice the digits. Inside the type name \code{\fa{صحیح۳۲}} the digits are Persian
(\fa{۳۲}), but arithmetic values such as \code{15} or loop bounds such as \code{5}
are written with the ordinary digits \code{0}--\code{9}. The Persian digits are part
of the \emph{spelling} of certain keywords; computation uses the standard digits.
\end{note}

\section{Why bilingualism matters}

This feature is more than a novelty. For a learner whose first language is
Persian, writing \code{\fa{اگر}} instead of \code{if} removes a layer of
translation between thought and code --- the program reads the way they think. A
classroom can teach loops and functions without first teaching English. And
because the structure is identical, the day that learner needs to work in English
--- to use an English library, join an international team, read English
documentation --- the transition is trivial: they already know every concept, and
only the spellings change.

\begin{tip}
If you are learning, start in whichever language is more comfortable, and once a
concept is solid, try rewriting a program in the other language. Because only the
keywords differ, this is quick --- and it cements both the idea \emph{and} its
English vocabulary at once, which is exactly the bridge Salam was built to be.
\end{tip}

\begin{deep}[In Depth: how the compiler reads two languages]
Internally, Salam keeps a \emph{language pack}: a table mapping each keyword to its
spellings in each supported language. When the lexer (the first stage of the
compiler, from Chapter~1) reads your source, it consults the pack for the language
you selected and recognises \code{\fa{تابع}} or \code{func} as the very same
internal token. From that point on --- parsing, type-checking, code
generation --- the compiler has no idea which language the source was written in,
because the difference has already been erased. That is why the two versions
produce byte-for-byte identical programs: after the first stage, there is only one
language left, the language of meaning.
\end{deep}

\section{Chapter summary}

\begin{itemize}[leftmargin=1.4em]
  \item Salam's grammar is identical in English and Persian; only the spelling of
        keywords and type names differs.
  \item Every keyword has a Persian twin (\code{func}/\code{\fa{تابع}},
        \code{if}/\code{\fa{اگر}}, \dots); the entry point is \code{main} /
        \code{\fa{اصلی}}. The full table is in Appendix~A.
  \item Compile a Persian file with \code{--lang=fa}; English is the default. The
        rest of the workflow is unchanged.
  \item A file uses one keyword language, but identifiers, strings, and comments
        may be in any language.
  \item Persian digits appear inside certain keyword spellings; computation uses
        ordinary digits.
  \item Bilingualism lets learners code in their own language and cross to English
        effortlessly, because the concepts are identical.
\end{itemize}

\section{Exercises}

\exercise Take any English program from an earlier chapter and rewrite it in
Persian, then run it with \code{--lang=fa}. Confirm the output matches.

\exercise Write a Persian program with a function (\code{\fa{تابع}}) that returns
(\code{\fa{بازگشت}}) the larger of two numbers, using \code{\fa{اگر}}.

\exercise In an English program, use a Persian word (in Latin letters, like
\code{naam}) as a variable name and a Persian sentence as a printed string.

\exercise Write the same ``count to ten'' loop twice --- once in English, once in
Persian --- and place them side by side to see how only the keywords change.

\exercise Explain in a comment why two versions of a program, one English and one
Persian, compile to identical results.

% ===================== Chapter 23 =====================
\chapter{Advanced Topics and Best Practices}
\label{ch:advanced}

You now know the whole language: types, control flow, functions, data structures,
generics, interfaces, closures, the standard library, concurrency, and the FFI.
This final chapter steps back from features to \emph{judgement} --- how to use what
you know well. It gathers advice on performance, structure, style, and the
pitfalls worth knowing, so that the programs you write next are not just correct
but good.

\section{Performance: fast by default, faster with care}

Salam programs are fast because they compile to C and run as native code, with
static types and no hidden boxing. For most programs that is the end of the
performance story --- you get speed for free. When you do need to squeeze more, a
few principles matter far more than micro-tricks:

\begin{itemize}[leftmargin=1.4em]
  \item \textbf{Choose the right data structure.} An algorithm that looks up by
        key thousands of times wants a \code{HashMap}, not a linear scan through a
        \code{Vector}. The biggest speed wins come from structure, not from
        shaving instructions.
  \item \textbf{Avoid needless allocation.} Building a string by concatenating in
        a tight loop allocates on every step. When you can, compute sizes up front
        or reuse a buffer. Reserve a vector's capacity if you know roughly how many
        elements it will hold.
  \item \textbf{Measure before optimising.} Intuition about what is slow is
        usually wrong. Time the program, find the part that actually dominates,
        and improve that --- not the part that merely \emph{looks} expensive.
  \item \textbf{Prefer the compiled path.} \code{salam exec} (the interpreter) is
        perfect for learning and quick runs, but a \code{salam build} executable
        is far faster for real work.
\end{itemize}

\begin{deep}[In Depth: optimization levels and the LLVM backend]
Because the default backend emits C, the C compiler's own optimizer does a great
deal of work on your behalf --- inlining small functions, unrolling loops,
eliminating dead code. Salam also has an experimental LLVM backend
(\code{salam llvm}) with explicit optimization levels (\code{-O0} through
\code{-O3}, plus \code{-Os}/\code{-Oz} for size), which can target native code or
WebAssembly. For everyday work the defaults are excellent; reach for the
optimization flags only when a measurement tells you to.
\end{deep}

\section{Structuring a program}

Small programs fit in one file; real ones do not. As a program grows, structure is
what keeps it understandable:

\begin{itemize}[leftmargin=1.4em]
  \item \textbf{Split by responsibility.} Put related types and functions together
        in a module (a file with a \code{package}), and keep each module focused on
        one area --- parsing, storage, the user interface. A reader should be able
        to guess what is in a file from its name.
  \item \textbf{Keep public surfaces small.} Mark only what other modules truly
        need with \code{pub}; keep helpers private. The less a module exposes, the
        freer you are to change its insides (Chapter~15).
  \item \textbf{Let data own its rules.} Put the operations that protect a struct's
        invariants in its methods, as the \code{Account} did in Chapter~9, so the
        rules live with the data they guard.
  \item \textbf{Depend on interfaces, not concrete types,} where you want
        flexibility (Chapter~11) --- it lets you swap implementations without
        rewriting the code that uses them.
\end{itemize}

\section{Coding style}

Style is not decoration; it is how you make code readable to the next person ---
who is often you, months later. Salam even ships a formatter, \code{salam fmt},
that rewrites your code into the canonical layout automatically, so you never argue
about spacing. Beyond formatting, a few habits carry most of the weight:

\begin{itemize}[leftmargin=1.4em]
  \item \textbf{Name for meaning.} \code{user\_count} and \code{is\_valid} tell the
        reader what a value \emph{is} and what a function \emph{answers}.
        Single-letter names earn their place only as short-lived loop indices.
  \item \textbf{Prefer immutability.} Reach for \code{mut} only when a value must
        change (Chapter~2). Every binding that cannot change is one fewer thing to
        track.
  \item \textbf{Keep functions small and single-purpose.} If you cannot name a
        function for the one thing it does, it is doing too many things
        (Chapter~6).
  \item \textbf{Comment the \emph{why}.} The code already shows \emph{what} it
        does; a comment is most valuable when it records the reasoning the code
        cannot.
  \item \textbf{Handle the unhappy path.} Check the \code{bool} a function returns,
        the \code{Option} that might be empty, the pointer that might be
        \code{null}. The cases you ignore are the bugs you ship.
\end{itemize}

\begin{tip}
Run \code{salam fmt} before you save or share code, and let it settle every
question of layout. Consistent formatting makes differences between versions show
real changes, not reshuffled whitespace, and frees your attention for what the code
actually does.
\end{tip}

\section{Common pitfalls, revisited}

A consolidated list of the traps this book has flagged --- worth a glance whenever
something behaves unexpectedly:

\begin{itemize}[leftmargin=1.4em]
  \item \textbf{Integer division truncates.} \code{7 / 2} is \code{3}; use a float
        operand for a fractional result (Chapter~3).
  \item \textbf{Literals are \code{i32}.} Storing one in a smaller or unsigned type
        needs an explicit \code{as} cast (Chapter~2).
  \item \textbf{Off-by-one with indices.} Arrays run \code{0} to \code{n-1}, while
        \code{repeat A to B} \emph{includes} \code{B} (Chapters~7 and~5).
  \item \textbf{Forgetting \code{free}.} Growable containers and manual allocations
        leak unless freed; pair each with a \code{defer} at creation
        (Chapter~18).
  \item \textbf{Unprotected shared state.} Two threads touching the same data
        without a lock race (Chapter~20).
  \item \textbf{Ignoring a result.} A \code{bool} or \code{Option} that signals
        failure is only useful if you check it (Chapter~14).
  \item \textbf{An \code{=} where you meant \code{==}.} Assignment versus
        comparison (Chapter~3).
\end{itemize}

\section{Where to go next}

The best way to consolidate everything in this book is to build something real ---
something a little beyond what you are sure you can do. A few projects that exercise
the whole language:

\begin{itemize}[leftmargin=1.4em]
  \item A \textbf{command-line tool} that reads a file, processes it, and writes a
        result --- exercising \code{io}, \code{str}, and the collections.
  \item A small \textbf{interpreter or calculator} that parses expressions ---
        exercising enums, structs, recursion, and error handling.
  \item A \textbf{simulation or game} (a grid world, a card game) --- exercising
        structs, methods, arrays, and \code{rand}.
  \item A \textbf{data tool} that reads JSON or CSV, tallies something with a
        \code{HashMap}, and reports --- exercising the standard library end to end.
\end{itemize}

Read the standard library's own source, too: it is written in Salam, and seeing how
\code{Vector} and \code{HashMap} are built is one of the best ways to deepen your
understanding. And when you hit something the language does not yet do, remember
the FFI --- the whole world of C is one \code{extern} block away.

\section{A closing word}

You began this book, perhaps, never having written a line of code; you end it
knowing a complete, modern programming language --- one you can write in your own
words. Programming is a craft, and like every craft it rewards practice more than
reading. So close the book and open an editor. Write a program that does something
you care about. Make it work, then make it clear, then make it fast --- in that
order. Argue with the compiler until you both agree. That is the whole of the
craft, and you now have every tool you need to practise it.

\bigskip
\noindent\emph{Salam, and happy coding.}

\section{Chapter summary}

\begin{itemize}[leftmargin=1.4em]
  \item Salam is fast by default; for more, choose good data structures, avoid
        needless allocation, and \emph{measure} before optimising.
  \item Structure a program by responsibility, keep public surfaces small, and let
        data own its rules.
  \item Good style --- meaningful names, immutability, small functions, ``why''
        comments, handled error paths --- is how code stays readable; \code{salam
        fmt} handles layout.
  \item Keep the list of common pitfalls in mind when something surprises you.
  \item Consolidate by building real projects and reading the standard library's
        source.
\end{itemize}

\section{Exercises}

\exercise Take a program you wrote for an earlier chapter and improve it: run
\code{salam fmt} on it, rename anything unclear, and add a comment explaining one
non-obvious decision.

\exercise Find a function in one of your programs longer than a screenful and split
it into smaller, well-named helpers.

\exercise Pick one project from ``Where to go next'' and sketch its modules: what
files, what \code{pub} functions, what data structures.

\exercise Review one of your programs against the ``common pitfalls'' list and fix
any you find.

\exercise Open a standard-library module (such as \code{collections}) and read one
of its types end to end. Write a paragraph on how it works.


\appendix
\chapter{Keyword Reference (English / Persian)}
\label{app:keywords}

This appendix lists every Salam keyword in both languages, drawn directly from the
compiler's language pack. A program is written in one language and compiled with
the matching \code{--lang} flag (English is the default, \code{--lang=fa} for
Persian); only the spellings below differ between the two the meaning and
grammar are identical.

\section*{Keywords}

\begin{center}
\begin{longtable}{@{}lll@{}}
\toprule
\textbf{English} & \textbf{Persian} & \textbf{Meaning} \\
\midrule
\endhead
\code{func}       & \code{\fa{تابع}}        & define a function \\
\code{ret}        & \code{\fa{بازگشت}}      & return from a function \\
\code{if}         & \code{\fa{اگر}}         & conditional \\
\code{else}       & \code{\fa{وگرنه}}       & alternative branch \\
\code{while}      & \code{\fa{تاوقتی}}      & while loop \\
\code{for}        & \code{\fa{برای}}        & C-style for loop \\
\code{repeat}     & \code{\fa{تکرار}}       & repeat / range loop \\
\code{to}         & \code{\fa{تا}}          & range upper bound (with \code{repeat}) \\
\code{every}      & \code{\fa{هر}}          & loop step (with \code{repeat ... to}) \\
\code{break}      & \code{\fa{بشکن}}        & exit the innermost loop \\
\code{continue}   & \code{\fa{ادامه}}       & skip to the next iteration \\
\code{mut}        & \code{\fa{متغیر}}       & mutable binding \\
\code{const}      & \code{\fa{ثابت}}        & compile-time constant \\
\code{type}       & \code{\fa{نوع}}         & type alias \\
\code{struct}     & \code{\fa{ساختار}}      & define a struct \\
\code{enum}       & \code{\fa{شمارش}}       & define an enum \\
\code{interface}  & \code{\fa{واسط}}        & define an interface \\
\code{impl}       & \code{\fa{پیاده‌سازی}}  & implementation block \\
\code{end}        & \code{\fa{پایان}}       & close a block \\
\code{import}     & \code{\fa{واردکردن}}    & import a module \\
\code{package}    & \code{\fa{بسته}}        & declare a package \\
\code{pub}        & \code{\fa{عمومی}}       & make a name public \\
\code{as}         & \code{\fa{بعنوان}}      & type conversion \\
\code{this}       & \code{\fa{این}}         & the current instance (in a method) \\
\code{true}       & \code{\fa{درست}}        & boolean true \\
\code{false}      & \code{\fa{نادرست}}      & boolean false \\
\code{null}       & \code{\fa{پوچ}}         & null pointer \\
\code{defer}      & \code{\fa{تعلیق}}       & schedule cleanup at scope exit \\
\code{extern}     & \code{\fa{خارجی}}       & declare external (C) functions \\
\code{operator}   & \code{\fa{عملگر}}       & operator overload \\
\code{component}  & \code{\fa{مولفه}}       & layout component \\
\code{layout}     & \code{\fa{صفحه}}        & layout block (web DSL) \\
\code{print}      & \code{\fa{بنویس}}       & print without a newline \\
\code{println}    & \code{\fa{چاپ}}         & print with a newline \\
\code{printerr}   & \code{\fa{خطابنویس}}    & print to standard error \\
\code{printerrln} & \code{\fa{خطاچاپ}}      & print to standard error with a newline \\
\code{input}      & \code{\fa{بخوان}}       & read a line of input \\
\code{main}       & \code{\fa{اصلی}}        & the program's entry point \\
\bottomrule
\end{longtable}
\end{center}

\begin{note}
\code{print} and \code{println} have the short aliases \code{\_} and \code{\_\_}
in English. The entry point \code{main} / \code{\fa{اصلی}} is part of the language:
the compiler looks for it as the place to start.
\end{note}

\section*{Type names}

\begin{center}
\begin{tabularx}{\linewidth}{@{}llX@{}}
\toprule
\textbf{English} & \textbf{Persian} & \textbf{Meaning} \\
\midrule
\code{i8}, \code{i16}, \code{i32}, \code{i64} & \code{\fa{صحیح۸}}, \code{\fa{صحیح۱۶}}, \code{\fa{صحیح۳۲}}, \code{\fa{صحیح۶۴}} & signed integers \\
\code{u8}, \code{u16}, \code{u32}, \code{u64} & \code{\fa{طبیعی۸}}, \dots & unsigned integers \\
\code{f32}, \code{f64} & \code{\fa{اعشار۳۲}}, \code{\fa{اعشار۶۴}} & floating-point \\
\code{bool}  & \code{\fa{منطقی}}   & boolean \\
\code{char}  & \code{\fa{کاراکتر}} & single byte / character \\
\code{str}   & \code{\fa{رشته}}    & string \\
\code{auto}  & \code{\fa{خودکار}}  & inferred type \\
\bottomrule
\end{tabularx}
\end{center}

\noindent The short spellings \code{int} ($=$ \code{i32}, and \code{\fa{صحیح}} in
Persian), \code{uint} ($=$ \code{u32}) and \code{float} ($=$ \code{f32}) name the
very same types; they are what this book uses in ordinary code. Note that the
digits inside
the Persian type names (for example \fa{۳۲} in \code{\fa{صحیح۳۲}}) are Persian
digits, part of the keyword's spelling; numeric \emph{values} in a program use the
ordinary digits \code{0}--\code{9}.

\chapter{Grammar (Simplified)}
\label{app:grammar}

This is a simplified, informal grammar of the Salam general language, in the style
of EBNF, intended as a reader's reference rather than a precise specification.
Square brackets \code{[ ]} mark an optional part, a star \code{*} means ``zero or
more,'' a plus \code{+} means ``one or more,'' and a vertical bar \code{|} means
``or.'' It reflects the colon-and-\code{end} block style used throughout this book.

\section*{Program structure}

\begin{shell}
program     = [ package ] import* definition*

package     = "package" IDENT

import      = "import" STRING
            | "import" "(" ( [ IDENT ] STRING )* ")"

definition  = function | struct | enum | interface
            | typealias | const | extern
\end{shell}

\section*{Functions}

\begin{shell}
function    = "func" IDENT [ generics ] "(" params ")" [ ":" type ] ":"
                  statement*
              "end"

params      = [ param ( "," param )* ]
param       = IDENT ":" type

generics    = "<" typeparam ( "," typeparam )* ">"
typeparam   = IDENT [ ":" IDENT ]          // optional interface bound
\end{shell}

\section*{Types}

\begin{shell}
type        = basetype dimension* [ "*" ]
basetype    = "i8" | "i16" | "i32" | "i64"
            | "u8" | "u16" | "u32" | "u64"
            | "f32" | "f64"
            | "bool" | "char" | "str" | "void"
            | "int" | "uint" | "float"          // aliases
            | IDENT [ "<" type ( "," type )* ">" ]   // struct/enum/generic
            | "dyn" IDENT                        // dynamic interface
dimension   = "[" [ INT ] "]"                    // array
\end{shell}

\section*{Declarations}

\begin{shell}
vardecl     = [ "mut" ] IDENT ":" type [ "=" expression ]
            | [ "mut" ] IDENT ":" "auto" "=" expression
const       = "const" IDENT [ ":" type ] "=" expression
typealias   = "type" IDENT "=" type
\end{shell}

\section*{Statements}

\begin{shell}
statement   = vardecl | const
            | assignment
            | expression                 // e.g. a function or method call
            | if | while | for | repeat
            | "break" | "continue"
            | "ret" [ expression ]
            | "defer" statement
            | "println" args | "print" args

assignment  = lvalue ( "=" | "+=" | "-=" | "*=" | "/=" | "%=" | "**=" ) expression
lvalue      = IDENT ( "." IDENT | "[" expression "]" )*
args        = [ expression ( "," expression )* ]
\end{shell}

\section*{Control flow}

\begin{shell}
if          = "if" expression ":" statement*
              ( "else" expression ":" statement* )*
              [ "else" ":" statement* ]
              "end"

while       = "while" expression ":" statement* "end"

for         = "for" vardecl "," expression "," assignment ":"
                  statement*
              "end"

repeat      = "repeat" expression ":" statement* "end"
            | "repeat" expression "to" expression [ "every" expression ] ":"
                  statement*
              "end"
\end{shell}

\section*{Structs, enums, and interfaces}

\begin{shell}
struct      = "struct" IDENT [ generics ] ":"
                  field*
                  method*
              "end"
field       = IDENT ":" type [ "=" expression ]
method      = function                       // defined inside the struct body

enum        = "enum" IDENT ":" member ( "," member )* [ "," ] "end"
            | "enum" IDENT ":" ( member )* "end"
member      = IDENT [ "=" INT ]

interface   = "interface" IDENT ":"
                  signature*
              "end"
signature   = "func" IDENT "(" params ")" [ ":" type ]   // no body
\end{shell}

\section*{Expressions (by decreasing precedence)}

\begin{shell}
primary     = INT | FLOAT | STRING | CHAR | "true" | "false" | "null"
            | IDENT
            | IDENT "(" args ")"                 // call
            | IDENT "{" ( IDENT "=" expression ("," ...)* )? "}"  // struct literal
            | "[" ( expression ("," expression)* )? "]"           // array literal
            | "(" params ")" "=>" ( expression | ":" statement* "end" )  // lambda
            | "(" expression ")"

postfix     = primary ( "." IDENT [ "(" args ")" ]   // field / method
                      | "[" expression "]" )*        // index

unary       = [ "!" | "-" ] postfix
power       = unary [ "**" power ]
cast        = power [ "as" type ]
product     = cast ( ( "*" | "/" | "%" ) cast )*
sum         = product ( ( "+" | "-" ) product )*
compare     = sum ( ( "<" | ">" | "<=" | ">=" | "==" | "!=" ) sum )*
logic       = compare ( ( "&&" | "||" ) compare )*
expression  = logic
\end{shell}

\section*{Foreign functions}

\begin{shell}
extern      = [ link ] "extern" "\"C\"" ":"
                  ( "func" IDENT "(" externparams ")" [ ":" type ] )*
              "end"
externparams= [ param ( "," param )* [ "," "..." ] ]   // ... = variadic
link        = "link" STRING
\end{shell}

\section*{Lexical notes}

\begin{itemize}[leftmargin=1.4em]
  \item \code{IDENT} a name: a letter or underscore, then letters, digits, or
        underscores. Unicode (including Persian) letters are allowed.
  \item \code{INT} a whole number (type \code{int} by default).
  \item \code{FLOAT} a number with a decimal point or exponent (type
        \code{f64}).
  \item \code{STRING} text in double quotes, or triple quotes
        (\code{"""..."""}); supports escapes \code{\textbackslash n},
        \code{\textbackslash t}, \code{\textbackslash\textbackslash},
        \code{\textbackslash"}.
  \item \code{CHAR} a single character in single quotes.
  \item Comments are \code{//} to end of line, or \code{/* ... */} across lines.
  \item Blocks open with \code{:} and close with \code{end}; a single-statement
        branch or body may follow the colon on the same line.
\end{itemize}

\chapter{Cheat Sheet}
\label{app:cheatsheet}

A compact reference to Salam's syntax and the most-used standard-library calls.
Every snippet is valid Salam; refer back to the named chapter for the full story.

\section*{Hello, world (Chapter 1)}

\begin{salam}
func main:
    println "Hello, Salam!"
end
\end{salam}

\noindent Run it: \code{salam exec hello.salam}\quad Build it:
\code{salam build hello.salam --output=hello}

\section*{Variables and types (Chapters 2--3)}

\begin{salam}
name: str = "Sara"          // immutable
mut count: int = 0          // mutable
const MAX = 100             // compile-time constant
x: auto = 3.14              // inferred (f64)
y: int = 250                // int by default
b: u8 = 250 as int as u8    // explicit conversion with 'as'
\end{salam}

\noindent Types: \code{i8 i16 i32 i64}, \code{u8 u16 u32 u64}, \code{f32 f64},
\code{bool}, \code{char}, \code{str}; \code{int}, \code{uint} and \code{float}
are the short spellings of \code{i32}, \code{u32} and \code{f32}. \\
Operators: \code{+ - * / \%}, \code{**} (power, float), \code{== != < > <= >=},
\code{\&\& || !}, \code{+= -= *= /= \%= **=}. \\
Integer \code{/} truncates; no bitwise operators.

\section*{Control flow (Chapter 4)}

\begin{salam}
if x > 0:
    println "positive"
else x == 0:
    println "zero"
else:
    println "negative"
end
\end{salam}

\section*{Loops (Chapter 5)}

\begin{salam}
until cond:           ... end
repeat 3:             ... end          // 3 times
repeat n with i:      ... end          // i = 0..n-1
repeat 1 to 5:        ... end          // 1..5 inclusive
repeat 0 to 20 by 2:      ... end      // step by 2
each x in xs:         ... end          // walk a collection
// break  -- exit loop;  continue  -- next iteration
\end{salam}

\section*{Functions (Chapter 6)}

\begin{salam}
func add(a: int, b: int): int:
    ret a + b
end

func greet(name: str):       // no return type = void
    println "Hi,", name
end
\end{salam}

\noindent Pass by value; no default or reference parameters; overloading by
parameter types; definition order does not matter.

\section*{Arrays and vectors (Chapter 7)}

\begin{salam}
a: int[3] = [1, 2, 3]            // fixed array, indexed 0..2
g: int[2][3] = [[1,2,3],[4,5,6]] // 2-D

mut v: Vector<int> = Vector {}
v.push(10)        // append
v.get(0)[0]       // read element 0
v.set(0, 99)      // write
v.len()           // count
v.pop()           // remove last
v.free()          // release memory
\end{salam}

\section*{Strings (Chapter 8)}

\begin{salam}
import "str"
"a" + "b"                  // concatenation
len("hello")               // 5
str.Upper(s)  str.Lower(s)  str.Trim(s)  str.Reverse(s)
str.Find(s, sub)  str.Contains(s, sub)  str.Replace(s, a, b)
str.Substr(s, start, n)
str.Split(s, sep)  str.Join(vec, sep)
str.FromInt(n as i64)  str.ToInt(s)
\end{salam}

\section*{Structs and methods (Chapter 9)}

\begin{salam}
struct Point:
    x: int = 0
    y: int = 0
    func move_by(dx: int, dy: int):
        this.x = this.x + dx
        this.y = this.y + dy
    end
end

mut p: Point = Point { x = 1, y = 2 }
p.move_by(3, 4)
println p.x, p.y
\end{salam}

\section*{Enums (Chapter 10)}

\begin{salam}
enum Day: Mon, Tue, Wed, Thu, Fri, Sat, Sun end
d: Day = Day.Sat
println d as int           // 5
// "match" with if / else if (no match statement)
\end{salam}

\section*{Interfaces and generics (Chapters 11--12)}

\begin{salam}
interface Shape:
    func area(): f64
end

func describe<T: Shape>(s: T):     // static bound
    println s.area()
end

func draw(s: dyn Shape): ... end   // dynamic dispatch

func Max<T>(a: T, b: T): T:          // generic function
    if a > b: ret a end
    ret b
end

struct Box<T>: value: T end        // generic struct
\end{salam}

\section*{Closures (Chapter 13)}

\begin{salam}
double: auto = (x: int) => x * 2
apply((x: int) => x + 1, 5)            // lambda as argument
func make_adder(n: int): func(int) int:
    ret (x: int) => x + n               // closure capturing n
end
\end{salam}

\section*{Errors and defer (Chapter 14)}

\begin{salam}
func withdraw(n: int): bool: ... end    // bool success flag

import "option"
option.Some(x)  option.IsSome(o)  option.UnwrapOr(o, fallback)
Option { value = 0, present = false } // a "none"

defer v.free()                        // runs at scope exit (LIFO)
\end{salam}

\section*{Modules (Chapter 15)}

\begin{salam}
package main
import "math" strutil "str" "io"      // strutil is an alias for str

pub func public_fn(): ... end         // exported; default is private
\end{salam}

\section*{Standard library (Chapters 16--17)}

\begin{salam}
io.Print(s)  io.Println(s)  io.Input()  io.ReadFile(p)  io.WriteFile(p, t)
fmt.Int(n as i64)  fmt.Float(x)  fmt.PadLeft(s, w)
math.Sqrt  math.Pow  math.Floor  math.MaxI  math.Gcd  math.PI
rand.Seed(s as i64)  rand.IntRange(lo, hi)  rand.Float()  rand.Bool()
time.Now()  time.Format(t)  time.Sleep(ms)
HashMap: insert(k,v)  get(k)  has(k)  remove(k)  size()  free()
Stack: push/pop/peek;  Queue: enqueue/dequeue
json.Valid  json.Get  json.GetInt  json.Has
regex.MatchStr(pat, s)  regex.FindStr(pat, s)  regex.ReplaceStr(pat, s, w)
os.Exists  os.Remove  os.Mkdir  os.ListDir  os.Args  os.Env  os.Exit
\end{salam}

\section*{Memory and FFI (Chapters 18--19)}

\begin{salam}
import "mem"
mem.Allocate(n as u64)  mem.Free(p)  mem.Reallocate  mem.Copy

extern:
    func printf(format: str, ...): int
    func sqrt(x: f64): f64
end
link "sqlite3"                        // link an external library
\end{salam}

\section*{Concurrency (Chapter 20)}

\begin{salam}
import "sync"
mut lock: sync.Mutex = sync.NewMutex()
t: i64 = spawn(work)                   // start a thread
join(t)                               // wait for it
sync.Lock(lock) ... sync.Unlock(lock) // critical section
sync.Destroy(lock)
\end{salam}

\section*{Testing (Chapter 21)}

\begin{salam}
import "testing"
testing.AssertEqInt(got, want)
testing.AssertEqStr(got, want)
testing.AssertTrue(cond)
testing.Summary()                     // prints "tests: N passed, M failed"
\end{salam}

\section*{The \texttt{salam} commands}

\begin{center}
\begin{tabularx}{\linewidth}{@{}lX@{}}
\toprule
\textbf{Command} & \textbf{Does} \\
\midrule
\code{salam exec f.salam}   & run with the interpreter (no C toolchain) \\
\code{salam run f.salam}    & build and run, keep nothing \\
\code{salam build f.salam --output=f} & build a native executable \\
\code{salam fmt f.salam}    & reformat source in place \\
\code{salam new name}       & scaffold a new project \\
\code{salam debug f.salam}  & build with debug info, launch a debugger \\
\code{salam memcheck f.salam} & build with AddressSanitizer and run \\
\code{salam ... --lang=fa}  & treat the source as Persian \\
\bottomrule
\end{tabularx}
\end{center}

\chapter{Solutions to Selected Exercises}%
\label{app:solutions}

Here are worked solutions to a selection of exercises from across the book. They
are not the only correct answers programming rewards many paths to the same
result but each is complete, idiomatic, and runs as shown. Try the exercise
yourself before reading the solution; the struggle is where the learning happens.

\section*{Chapter 1 Getting Started}

\textbf{Exercise 1.2} (print a labelled calculation):

\begin{salamen}
func main:
    println "12 * 8 =", 12 * 8
end
\end{salamen}

\begin{progout}
12 * 8 = 96
\end{progout}

\section*{Chapter 2 Variables and Data Types}

\textbf{Exercise 2.3} (average as a float):

\begin{salamen}
func main:
    sum : int = 10 + 15 + 20
    avg : f64 = (sum as f64) / 3.0
    println "average =", avg
end
\end{salamen}

\begin{progout}
average = 15
\end{progout}

The casts to \code{f64} are what make this a real division; without them
\code{45 / 3} would (here, coincidentally) also give 15, but for non-divisible
totals integer division would be wrong.

\section*{Chapter 3 Operators}

\textbf{Exercise 3.3} (between 1 and 12 inclusive):

\begin{salamen}
func main:
    n : int = 7
    in_range : auto = n >= 1 && n <= 12
    println n, "in 1..12?", in_range
end
\end{salamen}

\begin{progout}
7 in 1..12? true
\end{progout}

\textbf{Exercise 3.6} (cube in place with \code{\textasciicircum\textasciicircum=}):

\begin{salamen}
func main:
    mut side : f64 = 3.0
    side ^^= 3
    println side
end
\end{salamen}

\begin{progout}
27
\end{progout}

Starting from \code{mut side : int = 3} instead gives a compile error:
\code{\textasciicircum\textasciicircum} always produces a \code{f64}, and an \code{f64} cannot be assigned
back into an \code{int} variable without an explicit \code{as int} cast, so
\code{side \textasciicircum\textasciicircum= 3} is rejected the same way \code{side = side \textasciicircum\textasciicircum 3} would be.

\section*{Chapter 4 Control Flow}

\textbf{Exercise 4.3} (the sign of a number):

\begin{salamen}
func sign(n : int) : str:
    if n < 0: ret "negative"
    else n == 0: ret "zero"
    else: ret "positive"
    end
end

func main:
    println sign(-5), sign(0), sign(42)
end
\end{salamen}

\begin{progout}
negative zero positive
\end{progout}

\section*{Chapter 5 Loops}

\textbf{Exercise 5.2} (sum 1 to 100):

\begin{salamen}
func main:
    mut total : int = 0
    repeat 1 to 100 with i:
        total += i
    end
    println total
end
\end{salamen}

\begin{progout}
5050
\end{progout}

\section*{Chapter 6 Functions}

\textbf{Exercise 6.4} (recursive sum):

\begin{salamen}
func sum_to(n : int) : int:
    if n <= 0: ret 0 end       // base case
    ret n + sum_to(n - 1)      // recursive step
end

func main:
    println sum_to(10)
end
\end{salamen}

\begin{progout}
55
\end{progout}

\section*{Chapter 7 Arrays and Vectors}

\textbf{Exercise 7.3} (a vector of squares):

\begin{salamen}
func main:
    mut squares : Vector<int> = Vector {}
    repeat 1 to 10 with i:
        squares.push(i * i)
    end
    println "count:", squares.len()
    mut j : int = 0
    until j < squares.len():
        print squares.get(j)[0]
        print " "
        j += 1
    end
    println ""
    squares.free()
end
\end{salamen}

\begin{progout}
count: 10
1 4 9 16 25 36 49 64 81 100
\end{progout}

\section*{Chapter 8 Strings}

\textbf{Exercise 8.4} (palindrome test):

\begin{salamen}
import "str"

func is_palindrome(s : str) : bool:
    ret str.Equals(s, str.Reverse(s))
end

func main:
    println is_palindrome("level")
    println is_palindrome("salam")
end
\end{salamen}

\begin{progout}
true
false
\end{progout}

\section*{Chapter 9 Structs and Methods}

\textbf{Exercise 9.1} (a circle's area):

\begin{salamen}
struct Circle:
    radius : f64 = 0.0
    func area() : f64:
        ret 3.14159 * this.radius * this.radius
    end
end

func main:
    c : Circle = Circle { radius = 2.0 }
    println "area:", c.area()
end
\end{salamen}

\begin{progout}
area: 12.5664
\end{progout}

\section*{Chapter 10 Enums}

\textbf{Exercise 10.3} (turn right):

\begin{salamen}
enum Direction: North, East, South, West end

func turn_right(d : Direction) : Direction:
    if d == Direction.North: ret Direction.East
    else d == Direction.East: ret Direction.South
    else d == Direction.South: ret Direction.West
    else: ret Direction.North
    end
end

func main:
    mut d : Direction = Direction.North
    d = turn_right(d)
    println "now facing index:", d as int      // East = 1
end
\end{salamen}

\begin{progout}
now facing index: 1
\end{progout}

\section*{Chapter 11 Interfaces}

\textbf{Exercise 11.1} (a greeter interface):

\begin{salamen}
interface Greeter:
    func greet() : str
end

struct Formal:
    func greet() : str: ret "Good evening." end
end

struct Casual:
    func greet() : str: ret "Hey!" end
end

func announce<T: Greeter>(g : T):
    println g.greet()
end

func main:
    announce(Formal {})
    announce(Casual {})
end
\end{salamen}

\begin{progout}
Good evening.
Hey!
\end{progout}

\section*{Chapter 12 Generics}

\textbf{Exercise 12.1} (a generic \code{Min}):

\begin{salamen}
func Min<T>(a : T, b : T) : T:
    if a < b: ret a end
    ret b
end

func main:
    println Min(10, 20)
    println Min(2.5, 1.5)
end
\end{salamen}

\begin{progout}
10
1.5
\end{progout}

\section*{Chapter 13 Closures}

\textbf{Exercise 13.3} (a multiplier factory):

\begin{salamen}
func make_multiplier(factor : int) : func(int) int:
    ret (x : int) => x * factor
end

func main:
    times3 : auto = make_multiplier(3)
    times10 : auto = make_multiplier(10)
    println times3(5)
    println times10(5)
end
\end{salamen}

\begin{progout}
15
50
\end{progout}

\section*{Chapter 14 Errors and Defer}

\textbf{Exercise 14.2} (find the first negative):

\begin{salamen}
func find_first_negative(a : int[5]) : int:
    repeat 5 with i:
        if a[i] < 0:
            ret i
        end
    end
    ret -1
end

func main:
    data : int[5] = [3, 8, -2, 5, -7]
    println "first negative at:", find_first_negative(data)
end
\end{salamen}

\begin{progout}
first negative at: 2
\end{progout}

\section*{Chapter 16 Standard Library, Part 1}

\textbf{Exercise 16.3} (the hypotenuse):

\begin{salamen}
import "math"

func main:
    a : f64 = 3.0
    b : f64 = 4.0
    println "hypotenuse:", math.Sqrt(a * a + b * b)
end
\end{salamen}

\begin{progout}
hypotenuse: 5
\end{progout}

\section*{Chapter 17 Standard Library, Part 2}

\textbf{Exercise 17.1} (a map of ages):

\begin{salamen}
func main:
    mut ages : HashMap<str, int> = HashMap {}
    ages.insert("Sara", 30)
    ages.insert("Ali", 25)
    ages.insert("Mina", 41)
    println "Ali is", ages.get("Ali")
    println "have Reza?", ages.has("Reza")
    ages.free()
end
\end{salamen}

\begin{progout}
Ali is 25
have Reza? false
\end{progout}

\section*{Chapter 21 Testing}

\textbf{Exercise 21.1} (testing factorial):

\begin{salamen}
import "testing"

func factorial(n : int) : int:
    if n <= 1: ret 1 end
    ret n * factorial(n - 1)
end

func test_factorial():
    testing.AssertEqInt(factorial(0), 1)
    testing.AssertEqInt(factorial(5), 120)
    testing.AssertEqInt(factorial(3), 6)
end

func main:
    test_factorial()
    testing.Summary()
end
\end{salamen}

\begin{progout}
tests: 3 passed, 0 failed
\end{progout}

\section*{Beyond these}

The remaining exercises follow the same patterns shown here and in the chapters.
If you get stuck, re-read the relevant section, run a smaller version of the
problem, and add \code{println} probes to see what your code is actually doing ---
the very habits from Chapter~21. Every program in this book was written by someone
who started exactly where you are now.


\end{document}
